% Documento de Ingeniería del Proyecto – Fase 1 · Equipo B
% Compilar desde docs/documento/:
%   pdflatex -interaction=nonstopmode IS1-Fase1-EquipoB.tex   (dos veces, por el índice)
% Cada sección vive en secciones/ para que cada rol la edite en su rama sin conflictos.
\newcommand{\versiondoc}{1.1}
\newcommand{\fechadoc}{29/10/2026}
% Preámbulo del Documento de Ingeniería del Proyecto – Fase 1 (Equipo B).
% Formato del machote 0.3: carta, márgenes de 2.5 cm, 11 pt, páginas numeradas.
\PassOptionsToPackage{unicode}{hyperref}
\PassOptionsToPackage{hyphens}{url}
\documentclass[11pt,letterpaper,spanish]{article}
\usepackage[margin=2.5cm]{geometry}
\usepackage[T1]{fontenc}
\usepackage[utf8]{inputenc}
\usepackage{textcomp}
\usepackage{lmodern}
\usepackage{xcolor}
\usepackage{amsmath,amssymb}
\usepackage{microtype}
\usepackage{parskip}
\usepackage{longtable,booktabs,array}
\usepackage{calc}
\newcounter{none} % tablas sin número (formato heredado de la conversión)
\usepackage{etoolbox}
\makeatletter
\patchcmd\longtable{\par}{\if@noskipsec\mbox{}\fi\par}{}{}
\makeatother
\usepackage{footnotehyper}
\makesavenoteenv{longtable}
\usepackage{graphicx}
\usepackage{float}
\usepackage[labelformat=empty,font=small]{caption}
\usepackage{soul}
\usepackage[bidi=default,shorthands=off]{babel}
\usepackage{newunicodechar}
\newunicodechar{☐}{$\square$}
\newunicodechar{✓}{\checkmark}
\newunicodechar{≤}{$\leq$}
\newunicodechar{≥}{$\geq$}
\newunicodechar{→}{$\rightarrow$}
\usepackage{fancyhdr}
\pagestyle{fancy}
\fancyhf{}
\fancyhead[L]{\small IS1 · Equipo B · Sistema de Órdenes de Servicio y Citas}
\fancyhead[R]{\small v\versiondoc}
\fancyfoot[C]{\thepage}
\setlength{\emergencystretch}{3em}
\providecommand{\tightlist}{\setlength{\itemsep}{0pt}\setlength{\parskip}{0pt}}
\setcounter{secnumdepth}{-\maxdimen} % la numeración va en el título, como en el machote
\setcounter{tocdepth}{2}
\usepackage{bookmark}
\usepackage{xurl}
\urlstyle{same}
\hypersetup{pdflang={es}, hidelinks, pdftitle={IS1-Fase1-EquipoB-v\versiondoc}}

% ---------------------------------------------------------------------------
% Figuras. El machote pide número, título y ruta del archivo fuente.
%   \figura{archivo}{opciones de tamaño}{número}{título}{ruta fuente}
\newcommand{\figura}[5]{%
  \begin{figure}[H]\centering
  \includegraphics[#2,keepaspectratio]{#1}
  \caption{\textbf{Figura #3.} #4. Fuente: \texttt{#5}}
  \end{figure}}

% Figura pendiente: se deja el lugar, la descripción y la fuente, en rojo,
% sin la imagen. Se reemplaza por \figura cuando exista la imagen.
%   \figurapendiente{número}{título}{descripción de lo que debe mostrar}{ruta fuente}
\newcommand{\figurapendiente}[4]{%
  \begin{figure}[H]\centering
  \fcolorbox{red}{red!4}{\parbox{0.92\linewidth}{\color{red}\small
    \textbf{PENDIENTE: imagen de la Figura #1}\par\smallskip
    \textbf{Debe mostrar:} #3\par\smallskip
    \textbf{Archivo fuente:} \texttt{#4}}}
  \caption{\textcolor{red}{\textbf{Figura #1.} #2. Fuente: \texttt{#4}}}
  \end{figure}}

% Texto pendiente (dato que falta o que debe confirmarse con el cliente).
\newcommand{\pendiente}[1]{\textcolor{red}{\textbf{[Pendiente:} #1\textbf{]}}}


\begin{document}

\begin{titlepage}
\centering
{\large Universidad del Istmo · Campus Ixtepec\par}
Ingeniería en Desarrollo de Software y Sistemas Inteligentes\par
Ingeniería de Software I (IDSSI-27) · Quinto semestre\par
\vspace{3cm}
{\LARGE\bfseries Documento de Ingeniería del Proyecto -- Fase 1\par}
\vspace{1cm}
{\Large Sistema de Órdenes de Servicio y Citas\par}
\vspace{0.5cm}
Servicio Automotriz Hermanos Ruiz, S.A. de C.V.\par
\vfill
\begin{tabular}{rl}
\textbf{Equipo:} & B (5 integrantes)\\
\textbf{Cliente:} & Servicio Automotriz Hermanos Ruiz\\
\textbf{Producto:} & Sistema de Órdenes de Servicio y Citas\\
\textbf{Versión del documento:} & v\versiondoc\\
\textbf{Fecha:} & \fechadoc\\
\end{tabular}
\end{titlepage}

\tableofcontents
\clearpage

\section{0. Cómo usar este machote}\label{cuxf3mo-usar-este-machote}

Este documento es la plantilla del \textbf{Documento} \textbf{de}
\textbf{Ingeniería} \textbf{del} \textbf{Proyecto} \textbf{--}
\textbf{Fase} \textbf{1} que cada equipo entrega al profesor para su
evaluación. Reúne en un solo archivo lo que el equipo produce durante
los sprints 0 a 6 y que se encuentra en la carpeta docs/ del
repositorio.

\subsection{0.1 Reglas de uso}\label{reglas-de-uso}

\begin{itemize}
\item
  Conserven la numeración y el orden de las secciones. Si una sección no
  aplica, déjenla con la leyenda "No aplica" y una justificación de una
  línea.
\item
  Los recuadros grises marcados como \textbf{"Instrucción"} explican qué
  debe contener cada sección. \textbf{Bórrenlos antes de entregar.}
\item
  El texto gris entre corchetes, como {[}Escriba aquí\ldots{]}, es un
  marcador: reemplácenlo por contenido propio.
\item
  Las filas de ejemplo en las tablas solo muestran el formato;
  sustitúyanlas por los datos de su proyecto.
\item
  Los diagramas se insertan como imágenes legibles y su archivo fuente
  (draw.io, PlantUML o Mermaid) queda en el repositorio. Cada diagrama
  lleva número, título y ruta de su fuente: \emph{Figura 4.2. Máquina de
  estados de Pedido. Fuente: docs/04-analisis/estados/pedido.puml}.
\item
  Cuando el contenido sea muy extenso (por ejemplo, las 35 historias de
  usuario o los 25 casos de prueba), el documento incluye la tabla
  resumen y el detalle completo puede quedar en un anexo o en el
  repositorio, con su enlace.
\end{itemize}

\subsection{0.2 Entrega acumulativa}\label{entrega-acumulativa}

El documento \textbf{crece sprint a sprint}. En cada Sprint Review se
entrega la versión vigente con las secciones correspondientes completas;
el profesor evalúa solo las secciones del corte, pero puede pedir
corregir las anteriores.

{\def\LTcaptype{none} % do not increment counter
\begin{longtable}[]{@{}
  >{\raggedright\arraybackslash}p{(\linewidth - 6\tabcolsep) * \real{0.1923}}
  >{\raggedright\arraybackslash}p{(\linewidth - 6\tabcolsep) * \real{0.2031}}
  >{\raggedright\arraybackslash}p{(\linewidth - 6\tabcolsep) * \real{0.3909}}
  >{\raggedright\arraybackslash}p{(\linewidth - 6\tabcolsep) * \real{0.2137}}@{}}
\toprule\noalign{}
\begin{minipage}[b]{\linewidth}\raggedright
\textbf{Corte}
\end{minipage} & \begin{minipage}[b]{\linewidth}\raggedright
\textbf{Versión del documento}
\end{minipage} & \begin{minipage}[b]{\linewidth}\raggedright
\textbf{Secciones que deben estar completas}
\end{minipage} & \begin{minipage}[b]{\linewidth}\raggedright
\textbf{Evaluación en la que cuenta}
\end{minipage} \\
\midrule\noalign{}
\endhead
\bottomrule\noalign{}
\endlastfoot
Review 0 (sem. 5) & v0.1 & Portada, 1, 2.1-- 2.4, 3 & Parcial 1 \\
Review 1 (sem. 7) & v0.2 & + 4.1-- 4.5 (borrador), 4.6, 4.7 & Parcial
1 \\
Review 2 (sem. 9) & v1.0 (SRS firmado) & + 4 completa, 5.1--5.3, 2.5 &
Parcial 2 \\
Review 3 (sem. 11) & v1.1 & + 6 & Parcial 3 \\
Review 4 (sem. 13) & v1.2 & + 7.1--7.5, 9.1--9.3 & Parcial 3 \\
Review 5 (sem. 15) & v1.3 & + 7.6--7.8, 8, 9.4 & Ordinaria \\
Demo final (sem. 16) & v2.0 (final de la Fase 1) & Todo el documento,
incluidas 10, 11, 12, 13 y anexos & Ordinaria \\
\end{longtable}
}

\subsection{0.3 Formato de entrega}\label{formato-de-entrega}

\begin{itemize}
\item
  \textbf{Archivos:} este documento en Word (.docx) y en PDF, en
  docs/entregas/ del repositorio, y la liga del PDF en la minuta de la
  Review.
\item
  \textbf{Nombre del archivo:}
  IS1-Fase1-Equipo\textless A\textbar B\textgreater-v\textless versión\textgreater.pdf,
  por ejemplo IS1-Fase1-EquipoA-v1.0.pdf.
\item
  \textbf{Formato:} carta, márgenes de 2.5 cm, fuente de 11 puntos,
  encabezados numerados, numeración de páginas, figuras y tablas
  numeradas.
\item
  \textbf{Citas y referencias:} formato APA 7. Toda norma, libro o dato
  de terceros se cita; las referencias deben ser verificables.
\item
  \textbf{Uso de IA:} se declara en la sección 2.6. Cualquier integrante
  debe poder explicar cualquier parte del documento.
\end{itemize}

\begin{quote}
Si la versión del documento no coincide con lo que hay en el repositorio
y en el tablero, el profesor toma como válido el repositorio.
\end{quote}

\section{Portada}\label{portada}

{\def\LTcaptype{none} % do not increment counter
\begin{longtable}[]{@{}
  >{\raggedright\arraybackslash}p{(\linewidth - 2\tabcolsep) * \real{0.2990}}
  >{\raggedright\arraybackslash}p{(\linewidth - 2\tabcolsep) * \real{0.7010}}@{}}
\toprule\noalign{}
\begin{minipage}[b]{\linewidth}\raggedright
\textbf{Campo}
\end{minipage} & \begin{minipage}[b]{\linewidth}\raggedright
\textbf{Dato}
\end{minipage} \\
\midrule\noalign{}
\endhead
\bottomrule\noalign{}
\endlastfoot
Asignatura & Ingeniería de Software I (IDSSI-27) \\
Documento & Documento de Ingeniería del Proyecto -- Fase 1 \\
Equipo & B (5 integrantes) \\
Cliente & Servicio Automotriz Hermanos Ruiz \\
Nombre del producto & Sistema de Órdenes de Servicio y Citas \\
Versión del documento & v\versiondoc \\
Fecha & \fechadoc \\
Repositorio & \url{https://github.com/AbdiasChavez/Florentino} (repositorio original del equipo:
\url{https://github.com/AyagDan/SistemaDeServiciosYCitas}) \\
\end{longtable}
}

\subsection{Control de versiones}\label{control-de-versiones}

{\def\LTcaptype{none} % do not increment counter
\begin{longtable}[]{@{}
  >{\raggedright\arraybackslash}p{(\linewidth - 8\tabcolsep) * \real{0.1067}}
  >{\raggedright\arraybackslash}p{(\linewidth - 8\tabcolsep) * \real{0.1496}}
  >{\raggedright\arraybackslash}p{(\linewidth - 8\tabcolsep) * \real{0.2138}}
  >{\raggedright\arraybackslash}p{(\linewidth - 8\tabcolsep) * \real{0.3590}}
  >{\raggedright\arraybackslash}p{(\linewidth - 8\tabcolsep) * \real{0.1709}}@{}}
\toprule\noalign{}
\begin{minipage}[b]{\linewidth}\raggedright
\textbf{Versión}
\end{minipage} & \begin{minipage}[b]{\linewidth}\raggedright
\textbf{Fecha}
\end{minipage} & \begin{minipage}[b]{\linewidth}\raggedright
\textbf{Autor(es)}
\end{minipage} & \begin{minipage}[b]{\linewidth}\raggedright
\textbf{Descripción del cambio}
\end{minipage} & \begin{minipage}[b]{\linewidth}\raggedright
\textbf{Revisó}
\end{minipage} \\
\midrule\noalign{}
\endhead
\bottomrule\noalign{}
\endlastfoot
v0.1 & 22/10/2026 & Isis Yamilet Meléndez Hernández & Integración del
avance al machote oficial del profesor. & Ángel Yael Avendaño Guzmán \\
v0.1 & 29/10/2026 & Abdías Daniel Chávez Cruz & Correcciones de la
revisión interna: secciones 1 a 3 completas; identificadores RN/RF/HU/CU
unificados; estados del backlog corregidos; usuarios de GitHub en lugar
de correos personales; RN-11 a RN-13 y RF-11 agregados; base de datos
inicial (db/schema.sql). & Pendiente (otro integrante) \\
v1.1 (borrador) & 29/10/2026 & Abdías Daniel Chávez Cruz, con Claude Code
(ver 2.6), por rol: PO, SM, QA Lead, Developers & Migración a LaTeX por
secciones; sección 5 completa (plan, 35 CPA, matriz); sección 6 (modelos
de análisis, diagrama de clases v4); 7.1--7.3 y 9.1--9.3; 10.1 parcial;
figuras faltantes marcadas en rojo. La firma del SRS (Review 2) sigue
pendiente. & Pendiente (otro integrante) \\
\end{longtable}
}

\subsection{Aprobación del cliente}\label{aprobaciuxf3n-del-cliente}

No aplica

Esta sección no aplica para la versión v0.1 (Parcial 1) debido a que los
hitos de firma formal del SRS y del Release v0.1.0 se evalúan y
registran en la Review 2 y en la Demo Final, respectivamente.

{\def\LTcaptype{none} % do not increment counter
\begin{longtable}[]{@{}
  >{\raggedright\arraybackslash}p{(\linewidth - 8\tabcolsep) * \real{0.2137}}
  >{\raggedright\arraybackslash}p{(\linewidth - 8\tabcolsep) * \real{0.2349}}
  >{\raggedright\arraybackslash}p{(\linewidth - 8\tabcolsep) * \real{0.2522}}
  >{\raggedright\arraybackslash}p{(\linewidth - 8\tabcolsep) * \real{0.1601}}
  >{\raggedright\arraybackslash}p{(\linewidth - 8\tabcolsep) * \real{0.1390}}@{}}
\toprule\noalign{}
\begin{minipage}[b]{\linewidth}\raggedright
\textbf{Hito}
\end{minipage} & \begin{minipage}[b]{\linewidth}\raggedright
\textbf{Resultado (aceptado / aceptado con observaciones / rechazado)}
\end{minipage} & \begin{minipage}[b]{\linewidth}\raggedright
\textbf{Observaciones}
\end{minipage} & \begin{minipage}[b]{\linewidth}\raggedright
\textbf{Firma del cliente}
\end{minipage} & \begin{minipage}[b]{\linewidth}\raggedright
\textbf{Fecha}
\end{minipage} \\
\midrule\noalign{}
\endhead
\bottomrule\noalign{}
\endlastfoot
SRS v1.0 (Review 2) & & & & \\
Release v0.1.0 (Demo final) & & & & \\
\end{longtable}
}

\section{1. Información general del
proyecto}\label{informaciuxf3n-general-del-proyecto}

\subsection{1.1 Integrantes y roles}\label{integrantes-y-roles}

{\def\LTcaptype{none} % do not increment counter
\begin{longtable}[]{@{}
  >{\raggedright\arraybackslash}p{(\linewidth - 8\tabcolsep) * \real{0.2456}}
  >{\raggedright\arraybackslash}p{(\linewidth - 8\tabcolsep) * \real{0.1282}}
  >{\raggedright\arraybackslash}p{(\linewidth - 8\tabcolsep) * \real{0.2012}}
  >{\raggedright\arraybackslash}p{(\linewidth - 8\tabcolsep) * \real{0.1686}}
  >{\raggedright\arraybackslash}p{(\linewidth - 8\tabcolsep) * \real{0.2564}}@{}}
\toprule\noalign{}
\begin{minipage}[b]{\linewidth}\raggedright
\textbf{Nombre completo}
\end{minipage} & \begin{minipage}[b]{\linewidth}\raggedright
\textbf{Matrícula}
\end{minipage} & \begin{minipage}[b]{\linewidth}\raggedright
\textbf{Usuario GitHub}
\end{minipage} & \begin{minipage}[b]{\linewidth}\raggedright
\textbf{Rol(es) y sprints}
\end{minipage} & \begin{minipage}[b]{\linewidth}\raggedright
\textbf{Áreas principales de contribución}
\end{minipage} \\
\midrule\noalign{}
\endhead
\bottomrule\noalign{}
\endlastfoot
Abdías Daniel Chávez Cruz & 0224170024 & AbdiasChavez & Product Owner
(S0--S6)

Developer & Gestión del Product Backlog, relación con el cliente,
diagrama del proceso actual (3.5), base de datos. \\
Ángel Yael Avendaño Guzmán & 0224170023 & AyagDan & Scrum Master
(S0--S6)

Developer & Administración del repositorio, revisión e integración de
Pull Requests, documento inicial. \\
Dulce María Estudillo Rosado & 0224170028 & DulceRosado & Developer
(S0--S6) & Requisitos funcionales y no funcionales (4.4, 4.6),
stakeholders (3.4). \\
Isis Alhelí González Maldonado & {[}Pendiente{]} &
IsisAlheliGonzalezMaldonado & Developer (S0--S6) & Historias de usuario
y criterios de aceptación (4.5). \\
Isis Yamilet Meléndez Hernández & 0224170037 & IsisMelendezHernandez &
Developer (S0--S6) & Integración del machote, descripción general (4.3),
reglas de negocio (4.7), protección de datos (4.8). \\
\end{longtable}
}

\subsection{1.2 Enlaces del proyecto}\label{enlaces-del-proyecto}

{\def\LTcaptype{none} % do not increment counter
\begin{longtable}[]{@{}
  >{\raggedright\arraybackslash}p{(\linewidth - 2\tabcolsep) * \real{0.3631}}
  >{\raggedright\arraybackslash}p{(\linewidth - 2\tabcolsep) * \real{0.6369}}@{}}
\toprule\noalign{}
\begin{minipage}[b]{\linewidth}\raggedright
\textbf{Recurso}
\end{minipage} & \begin{minipage}[b]{\linewidth}\raggedright
\textbf{Enlace}
\end{minipage} \\
\midrule\noalign{}
\endhead
\bottomrule\noalign{}
\endlastfoot
Repositorio & \url{https://github.com/AbdiasChavez/Florentino} (original:
\url{https://github.com/AyagDan/SistemaDeServiciosYCitas}) \\
Integración continua & \url{https://github.com/AbdiasChavez/Florentino/actions} \\
Tablero (GitHub Projects / Jira) & \pendiente{crear el tablero en GitHub Projects con las HU-01 a
HU-30 como issues}. \\
Prototipo de UI (Figma / Penpot) & \pendiente{se agrega al crear el prototipo (Sprint 4)}. \\
Ambiente de pruebas (staging) & No aplica todavía (desde el Sprint
4). \\
Release v0.1.0 & No aplica todavía (Demo final). \\
Canal de comunicación del equipo & Pendiente: registrar el canal
acordado (ver CONTRIBUTING.md). \\
\end{longtable}
}

\subsection{1.3 Resumen ejecutivo}\label{resumen-ejecutivo}

Servicio Automotriz Hermanos Ruiz opera hoy con cuadernos, órdenes en
papel y llamadas. Esto provoca disputas por cobros que el cliente dice
no haber autorizado, reclamos por daños sin evidencia del estado en que
llegó el vehículo, muchas llamadas para preguntar por el avance y poca
visibilidad de qué vehículo ocupa cada una de las 4 bahías.

El Sistema de Órdenes de Servicio y Citas digitaliza el ciclo de
atención: la cita, la recepción del vehículo con checklist y
fotografías, el diagnóstico y la asignación de bahía, la cotización que
el cliente autoriza desde su celular con un código de verificación, y el
cobro y la entrega. El cliente consulta el estado de su vehículo con un
enlace seguro, sin instalar aplicaciones.

Estado actual (v\versiondoc): se levantaron los requisitos con el
cliente (entrevista del 19/10/2026) y se definió el alcance. Ya funciona,
en una versión de desarrollo sin pantallas todavía, el corazón del
sistema: abrir una orden con checklist y fotografías, asignar bahía y
mecánico, registrar el diagnóstico, enviar la cotización al cliente,
autorizarla con el código de verificación, cobrar y entregar sin saldo
pendiente. Cada una de estas reglas se comprueba automáticamente con 35
pruebas. Falta: las pantallas para el personal y para el cliente, el envío
real de correos, la operación sin internet, el almacén y el ambiente de
pruebas en línea. La Fase 2 agregará reportes gerenciales, consulta de
existencias y de garantías, y el tablero simplificado de bahías.

\section{2. Gestión del proyecto con
Scrum}\label{gestiuxf3n-del-proyecto-con-scrum}

\subsection{2.1 Acuerdos de trabajo}\label{acuerdos-de-trabajo}

Los acuerdos completos están en CONTRIBUTING.md del repositorio. En
resumen:

\textbf{• Comunicación:} canal del equipo (por confirmar) para lo
urgente y GitHub (issues y Pull Requests) para todo lo técnico. Tiempo
máximo de respuesta: 24 horas en días hábiles.

\textbf{• Reuniones:} Daily breve de lunes a viernes (presencial o por
mensaje) y los eventos de Scrum al inicio y al cierre de cada sprint de
dos semanas.

\textbf{• Decisiones:} por consenso; si no se logra, decide el Product
Owner en temas de alcance y prioridad, y el Scrum Master en temas de
proceso.

\textbf{• Ramas:} main está protegida; cada integrante trabaja en
rama/\textless NOMBRE\textgreater{} y se integra solo mediante Pull
Request con al menos una revisión de otro integrante.

\textbf{• Commits:} en español, en presente y con referencia a la
historia (por ejemplo, «Agrega checklist de recepción (HU-01)»).

\textbf{• Incumplimiento:} se registra en la retrospectiva; a la segunda
falta en el mismo sprint, el Scrum Master lo informa al profesor.

\subsection{2.2 Definición de Listo (DoR) y Definición de Terminado
(DoD)}\label{definiciuxf3n-de-listo-dor-y-definiciuxf3n-de-terminado-dod}

{\def\LTcaptype{none} % do not increment counter
\begin{longtable}[]{@{}
  >{\raggedright\arraybackslash}p{(\linewidth - 2\tabcolsep) * \real{0.5001}}
  >{\raggedright\arraybackslash}p{(\linewidth - 2\tabcolsep) * \real{0.4999}}@{}}
\toprule\noalign{}
\begin{minipage}[b]{\linewidth}\raggedright
\textbf{Definición de Listo}
\end{minipage} & \begin{minipage}[b]{\linewidth}\raggedright
\textbf{Definición de Terminado}
\end{minipage} \\
\midrule\noalign{}
\endhead
\bottomrule\noalign{}
\endlastfoot
• La historia sigue el formato Como\ldots{} quiero\ldots{} para\ldots{}
y tiene ID e issue en el tablero.

• Tiene al menos 2 criterios de aceptación en formato
Dado/Cuando/Entonces.

• El equipo la estimó en puntos y cabe en un sprint (8 puntos o menos).

• Sus reglas de negocio (RN) y requisitos (RF/RNF) relacionados están
identificados.

• Sus dependencias (otras historias, datos, wireframe) están resueltas o
planificadas.

• El Product Owner la priorizó (MoSCoW). & • El código está en main
mediante un Pull Request revisado por otro integrante.

• Cumple todos sus criterios de aceptación y el Product Owner lo
verificó.

• Tiene pruebas automatizadas y la integración continua pasa sin
errores.

• Se probó solo con datos ficticios; no hay datos personales reales ni
contraseñas en el repositorio.

• La documentación afectada (SRS, modelos, README) está actualizada.

• Está desplegada en el ambiente de pruebas (desde el Sprint 4). \\
\end{longtable}
}

\subsection{2.3 Resumen de sprints}\label{resumen-de-sprints}

{\def\LTcaptype{none} % do not increment counter
\begin{longtable}[]{@{}
  >{\raggedright\arraybackslash}p{(\linewidth - 10\tabcolsep) * \real{0.0856}}
  >{\raggedright\arraybackslash}p{(\linewidth - 10\tabcolsep) * \real{0.1495}}
  >{\raggedright\arraybackslash}p{(\linewidth - 10\tabcolsep) * \real{0.3270}}
  >{\raggedright\arraybackslash}p{(\linewidth - 10\tabcolsep) * \real{0.1601}}
  >{\raggedright\arraybackslash}p{(\linewidth - 10\tabcolsep) * \real{0.1389}}
  >{\raggedright\arraybackslash}p{(\linewidth - 10\tabcolsep) * \real{0.1389}}@{}}
\toprule\noalign{}
\begin{minipage}[b]{\linewidth}\raggedright
\textbf{Sprint}
\end{minipage} & \begin{minipage}[b]{\linewidth}\raggedright
\textbf{Fechas}
\end{minipage} & \begin{minipage}[b]{\linewidth}\raggedright
\textbf{Objetivo del sprint}
\end{minipage} & \begin{minipage}[b]{\linewidth}\raggedright
\textbf{¿Se cumplió? (sí / parcial / no)}
\end{minipage} & \begin{minipage}[b]{\linewidth}\raggedright
\textbf{Puntos comprometidos}
\end{minipage} & \begin{minipage}[b]{\linewidth}\raggedright
\textbf{Puntos terminados}
\end{minipage} \\
\midrule\noalign{}
\endhead
\bottomrule\noalign{}
\endlastfoot
S0 & Por registrar & Conocer al cliente, levantar los requisitos
iniciales y preparar el repositorio y el documento base. & En curso &
--- & --- \\
S1 & & & & & \\
S2 & & & & & \\
S3 & & & & & \\
S4 & & & & & \\
S5 & & & & & \\
S6 & & & & & \\
\end{longtable}
}

\subsubsection{Evidencia de los eventos}\label{evidencia-de-los-eventos}

{\def\LTcaptype{none} % do not increment counter
\begin{longtable}[]{@{}
  >{\raggedright\arraybackslash}p{(\linewidth - 8\tabcolsep) * \real{0.0962}}
  >{\raggedright\arraybackslash}p{(\linewidth - 8\tabcolsep) * \real{0.2136}}
  >{\raggedright\arraybackslash}p{(\linewidth - 8\tabcolsep) * \real{0.2137}}
  >{\raggedright\arraybackslash}p{(\linewidth - 8\tabcolsep) * \real{0.3162}}
  >{\raggedright\arraybackslash}p{(\linewidth - 8\tabcolsep) * \real{0.1604}}@{}}
\toprule\noalign{}
\begin{minipage}[b]{\linewidth}\raggedright
\textbf{Sprint}
\end{minipage} & \begin{minipage}[b]{\linewidth}\raggedright
\textbf{Minuta de Review}
\end{minipage} & \begin{minipage}[b]{\linewidth}\raggedright
\textbf{Retrospectiva}
\end{minipage} & \begin{minipage}[b]{\linewidth}\raggedright
\textbf{Acciones de mejora acordadas}
\end{minipage} & \begin{minipage}[b]{\linewidth}\raggedright
\textbf{¿Se aplicaron?}
\end{minipage} \\
\midrule\noalign{}
\endhead
\bottomrule\noalign{}
\endlastfoot
S0 & Pendiente (al cerrar el Sprint 0) & Pendiente (al cerrar el Sprint
0) & --- & --- \\
& & & & \\
& & & & \\
\end{longtable}
}

\subsection{2.4 Registro de riesgos}\label{registro-de-riesgos}

{\def\LTcaptype{none} % do not increment counter
\begin{longtable}[]{@{}
  >{\raggedright\arraybackslash}p{(\linewidth - 10\tabcolsep) * \real{0.0856}}
  >{\raggedright\arraybackslash}p{(\linewidth - 10\tabcolsep) * \real{0.2990}}
  >{\raggedright\arraybackslash}p{(\linewidth - 10\tabcolsep) * \real{0.1282}}
  >{\raggedright\arraybackslash}p{(\linewidth - 10\tabcolsep) * \real{0.1069}}
  >{\raggedright\arraybackslash}p{(\linewidth - 10\tabcolsep) * \real{0.2735}}
  >{\raggedright\arraybackslash}p{(\linewidth - 10\tabcolsep) * \real{0.1067}}@{}}
\toprule\noalign{}
\begin{minipage}[b]{\linewidth}\raggedright
\textbf{ID}
\end{minipage} & \begin{minipage}[b]{\linewidth}\raggedright
\textbf{Riesgo}
\end{minipage} & \begin{minipage}[b]{\linewidth}\raggedright
\textbf{Probabilidad}
\end{minipage} & \begin{minipage}[b]{\linewidth}\raggedright
\textbf{Impacto}
\end{minipage} & \begin{minipage}[b]{\linewidth}\raggedright
\textbf{Respuesta y responsable}
\end{minipage} & \begin{minipage}[b]{\linewidth}\raggedright
\textbf{Estado}
\end{minipage} \\
\midrule\noalign{}
\endhead
\bottomrule\noalign{}
\endlastfoot
R-01 & Un integrante abandona el curso a mitad del semestre & Media &
Alto & Mitigar: conocimiento compartido mediante revisión cruzada de PR;
responsable: SM & Abierto \\
R-02 & La operación sin conexión y la sincronización posterior (RNF-08,
RNF-09) resultan más complejas de lo estimado & Alta & Alto & Mitigar:
prototipo técnico en el Sprint 1 y limitar lo offline a recepción y
diagnóstico; responsable: equipo de desarrollo & Abierto \\
R-03 & El cliente no está disponible para validar requisitos o asistir a
las Reviews & Media & Alto & Mitigar: agendar sesiones con anticipación
y validar por correo cuando no pueda asistir; responsable: PO &
Abierto \\
R-04 & Cambios de alcance frecuentes por parte del cliente & Media &
Medio & Mitigar: registrar cada solicitud de cambio (2.5) con análisis
de impacto; responsable: PO & Abierto \\
R-05 & El servicio gratuito de correo limita el envío de enlaces y
códigos 2FA & Media & Alto & Mitigar: probar los límites del plan
gratuito en el Sprint 1 y dejar SMS como opción (RF-02.3); responsable:
equipo de desarrollo & Abierto \\
R-06 & Las fotografías de recepción saturan el almacenamiento del
hospedaje gratuito & Media & Medio & Mitigar: comprimir las imágenes en
el dispositivo y fijar un tamaño máximo por foto; responsable: equipo de
desarrollo & Abierto \\
R-07 & Poca experiencia del equipo con FastAPI y PWA & Alta & Medio &
Mitigar: capacitación en S0--S1 y trabajo en pares en las primeras
historias; responsable: SM & Abierto \\
R-08 & Conflictos al editar en paralelo el documento .docx en ramas
distintas & Alta & Medio & Mitigar: un solo integrador por entrega y
avances por sección en docs/; responsable: SM & Abierto \\
R-09 & Se publican datos personales reales en el repositorio público &
Media & Alto & Evitar: solo datos ficticios y revisión de datos
personales en cada PR; responsable: todos & Abierto \\
\end{longtable}
}

\subsection{2.5 Solicitudes de cambio}\label{solicitudes-de-cambio}

{\def\LTcaptype{none} % do not increment counter
\begin{longtable}[]{@{}
  >{\raggedright\arraybackslash}p{(\linewidth - 10\tabcolsep) * \real{0.1282}}
  >{\raggedright\arraybackslash}p{(\linewidth - 10\tabcolsep) * \real{0.1068}}
  >{\raggedright\arraybackslash}p{(\linewidth - 10\tabcolsep) * \real{0.2137}}
  >{\raggedright\arraybackslash}p{(\linewidth - 10\tabcolsep) * \real{0.2309}}
  >{\raggedright\arraybackslash}p{(\linewidth - 10\tabcolsep) * \real{0.2137}}
  >{\raggedright\arraybackslash}p{(\linewidth - 10\tabcolsep) * \real{0.1067}}@{}}
\toprule\noalign{}
\begin{minipage}[b]{\linewidth}\raggedright
\textbf{ID / issue}
\end{minipage} & \begin{minipage}[b]{\linewidth}\raggedright
\textbf{Fecha}
\end{minipage} & \begin{minipage}[b]{\linewidth}\raggedright
\textbf{Solicitud del cliente}
\end{minipage} & \begin{minipage}[b]{\linewidth}\raggedright
\textbf{Impacto (requisitos, modelos, pruebas, esfuerzo)}
\end{minipage} & \begin{minipage}[b]{\linewidth}\raggedright
\textbf{Decisión y justificación}
\end{minipage} & \begin{minipage}[b]{\linewidth}\raggedright
\textbf{Release}
\end{minipage} \\
\midrule\noalign{}
\endhead
\bottomrule\noalign{}
\endlastfoot
--- & --- & No se han recibido solicitudes de cambio del cliente a la
fecha. & --- & --- & --- \\
\end{longtable}
}

\subsection{2.6 Declaración de uso de inteligencia
artificial}\label{declaraciuxf3n-de-uso-de-inteligencia-artificial}

{\def\LTcaptype{none} % do not increment counter
\begin{longtable}[]{@{}
  >{\raggedright\arraybackslash}p{(\linewidth - 8\tabcolsep) * \real{0.1708}}
  >{\raggedright\arraybackslash}p{(\linewidth - 8\tabcolsep) * \real{0.2137}}
  >{\raggedright\arraybackslash}p{(\linewidth - 8\tabcolsep) * \real{0.2351}}
  >{\raggedright\arraybackslash}p{(\linewidth - 8\tabcolsep) * \real{0.2307}}
  >{\raggedright\arraybackslash}p{(\linewidth - 8\tabcolsep) * \real{0.1497}}@{}}
\toprule\noalign{}
\begin{minipage}[b]{\linewidth}\raggedright
\textbf{Herramienta}
\end{minipage} & \begin{minipage}[b]{\linewidth}\raggedright
\textbf{Sección o artefacto}
\end{minipage} & \begin{minipage}[b]{\linewidth}\raggedright
\textbf{Para qué se usó}
\end{minipage} & \begin{minipage}[b]{\linewidth}\raggedright
\textbf{Cómo se verificó}
\end{minipage} & \begin{minipage}[b]{\linewidth}\raggedright
\textbf{Responsable}
\end{minipage} \\
\midrule\noalign{}
\endhead
\bottomrule\noalign{}
\endlastfoot
Claude (Anthropic), mediante Claude Code & Revisión de consistencia del
documento v0.1; borradores de 1.3, 2.1--2.4, 3.1--3.3, 3.6, 4.1 y 4.2;
correcciones de trazabilidad en 4.4--4.9; esquema de base de datos
(db/schema.sql). & Detectar inconsistencias y proponer redacción a
partir del anteproyecto, la entrevista del 19/10/2026 y el diagrama de
clases v4. & Revisión del Product Owner contra las fuentes; falta la
revisión de otro integrante antes de la entrega. & Abdías Daniel Chávez
Cruz \\
Por declarar & Secciones redactadas antes de la v0.1 y Propuesta de
Solución Tecnológica & Cada autor declara la herramienta que usó o que
no usó IA. & --- & Cada integrante \\
\end{longtable}
}

\section{3. Visión y viabilidad del
producto}\label{visiuxf3n-y-viabilidad-del-producto}

\subsection{3.1 Problema y contexto del
cliente}\label{problema-y-contexto-del-cliente}

Servicio Automotriz Hermanos Ruiz, S.A. de C.V. es un taller que atiende
vehículos ligeros (automóviles particulares, motocicletas y camionetas)
de clientes particulares y de pequeñas flotillas empresariales. Tiene 4
bahías de trabajo y un estacionamiento para 5 vehículos en espera. Hoy
las citas se apuntan en un cuaderno, las órdenes se llenan en papel y la
comunicación con el cliente es por teléfono o WhatsApp (ver el proceso
actual en 3.5).

Según el anteproyecto y la entrevista con el cliente del 19/10/2026, los
problemas principales son:

\textbf{• Disputas por cobros:} las autorizaciones son verbales o por
mensaje y no queda registro de qué se aprobó ni por cuánto («dijo que
sí, pero cobra otra cosa»).

\textbf{• Reclamos por daños:} no hay checklist ni fotografías del
estado en que llega el vehículo.

\textbf{• Exceso de llamadas:} el cliente llama para saber el avance y
la recepcionista interrumpe su trabajo para averiguarlo.

\textbf{• Poca visibilidad:} la carga de trabajo y las bahías se asignan
de palabra, y hay vehículos detenidos sin justificación.

\textbf{• Historial disperso:} las órdenes se archivan en carpetas y es
difícil consultar reparaciones y garantías anteriores.

Hallazgos de la entrevista que no estaban en el anteproyecto:

• El cliente quiere una doble verificación con código de 6 dígitos al
autorizar cotizaciones, en forma remota o presencial (RN-02).

• El cliente puede aprobar o rechazar por separado los conceptos
opcionales de una cotización (RN-03).

• Si una refacción tarda varios días, el vehículo pasa al
estacionamiento para liberar la bahía (RN-06).

• Las citas pedidas en línea no apartan lugar hasta que la recepcionista
las confirma (RN-08).

• Varios mecánicos pueden colaborar en una orden, pero siempre hay un
responsable (RN-09).

Además, hay cortes de luz e internet, por lo que el personal necesita
seguir capturando sin conexión, y el jefe de mecánicos trabajará con 1 o
2 tabletas en un ambiente con ruido, polvo y manos sucias (4.3).

\subsection{3.2 Objetivos del producto e
indicadores}\label{objetivos-del-producto-e-indicadores}

{\def\LTcaptype{none} % do not increment counter
\begin{longtable}[]{@{}
  >{\raggedright\arraybackslash}p{(\linewidth - 8\tabcolsep) * \real{0.2349}}
  >{\raggedright\arraybackslash}p{(\linewidth - 8\tabcolsep) * \real{0.2030}}
  >{\raggedright\arraybackslash}p{(\linewidth - 8\tabcolsep) * \real{0.1710}}
  >{\raggedright\arraybackslash}p{(\linewidth - 8\tabcolsep) * \real{0.1708}}
  >{\raggedright\arraybackslash}p{(\linewidth - 8\tabcolsep) * \real{0.2202}}@{}}
\toprule\noalign{}
\begin{minipage}[b]{\linewidth}\raggedright
\textbf{Objetivo}
\end{minipage} & \begin{minipage}[b]{\linewidth}\raggedright
\textbf{Indicador}
\end{minipage} & \begin{minipage}[b]{\linewidth}\raggedright
\textbf{Situación actual}
\end{minipage} & \begin{minipage}[b]{\linewidth}\raggedright
\textbf{Meta}
\end{minipage} & \begin{minipage}[b]{\linewidth}\raggedright
\textbf{Cómo se medirá}
\end{minipage} \\
\midrule\noalign{}
\endhead
\bottomrule\noalign{}
\endlastfoot
Eliminar las disputas por cobros no autorizados & \% de órdenes cerradas
cuya cotización tiene autorización digital registrada (fecha, hora, IP y
monto) & Autorizaciones verbales o por mensaje, sin registro & 100 \% de
las órdenes cerradas & Reporte mensual del sistema: órdenes cerradas con
autorización confirmada / total de órdenes cerradas \\
Reducir las llamadas de seguimiento a recepción & Llamadas por día de
clientes que preguntan por el avance de su vehículo & Por medir:
bitácora de llamadas de una semana antes de liberar el portal del
cliente & Al menos 50 \% menos que la medición inicial & Bitácora de
llamadas de recepción, una semana antes y una semana después de la
liberación \\
Contar con evidencia del estado del vehículo al recibirlo & \% de
recepciones con checklist completo y las 4 fotografías obligatorias &
Órdenes en papel sin checklist ni fotografías & 100 \% de las
recepciones & Reporte mensual del sistema de recepciones completas \\
\end{longtable}
}

\subsection{3.3 Alcance}\label{alcance}

{\def\LTcaptype{none} % do not increment counter
\begin{longtable}[]{@{}
  >{\raggedright\arraybackslash}p{(\linewidth - 4\tabcolsep) * \real{0.3333}}
  >{\raggedright\arraybackslash}p{(\linewidth - 4\tabcolsep) * \real{0.3333}}
  >{\raggedright\arraybackslash}p{(\linewidth - 4\tabcolsep) * \real{0.3333}}@{}}
\toprule\noalign{}
\begin{minipage}[b]{\linewidth}\raggedright
\textbf{Dentro del alcance de la Fase 1}
\end{minipage} & \begin{minipage}[b]{\linewidth}\raggedright
\textbf{Dentro del alcance de la Fase 2}
\end{minipage} & \begin{minipage}[b]{\linewidth}\raggedright
\textbf{Fuera de alcance}
\end{minipage} \\
\midrule\noalign{}
\endhead
\bottomrule\noalign{}
\endlastfoot
• Acceso del personal por rol (HU-27)

• Citas en línea y agenda con confirmación manual (HU-03, HU-04, HU-06)

• Recepción con checklist y fotografías, también sin conexión (HU-01,
HU-25, HU-26)

• Asignación de bahía y mecánico, y diagnóstico (HU-07 a HU-09)

• Cotización desglosada y autorización del cliente con código 2FA
(HU-02, HU-12 a HU-16)

• Solicitud de refacciones y proveedores (HU-10, HU-17, HU-18)

• Aviso de fin de trabajos, cobro, entrega y crédito a flotillas (HU-20
a HU-23)

• Portal del cliente responsivo y aviso de privacidad (HU-28, HU-30) & •
Búsqueda rápida de clientes y vehículos (HU-05)

• Tablero simplificado de bahías para tableta (HU-11)

• Consulta de existencias de refacciones (HU-19)

• Consulta de garantías por el cliente (HU-24)

• Panel y reportes gerenciales (HU-29, RF-11)

• Envío del enlace por SMS (RF-02.3) & • Facturación electrónica (CFDI)
y contabilidad

• Pagos en línea con tarjeta

• Aplicación móvil nativa (iOS o Android)

• Vehículos pesados, tráileres y maquinaria (RN-07)

• Nómina y control de asistencia del personal \\
\end{longtable}
}

\subsection{3.4 Stakeholders}\label{stakeholders}

{\def\LTcaptype{none} % do not increment counter
\begin{longtable}[]{@{}
  >{\raggedright\arraybackslash}p{(\linewidth - 8\tabcolsep) * \real{0.1922}}
  >{\raggedright\arraybackslash}p{(\linewidth - 8\tabcolsep) * \real{0.2031}}
  >{\raggedright\arraybackslash}p{(\linewidth - 8\tabcolsep) * \real{0.2349}}
  >{\raggedright\arraybackslash}p{(\linewidth - 8\tabcolsep) * \real{0.1496}}
  >{\raggedright\arraybackslash}p{(\linewidth - 8\tabcolsep) * \real{0.2202}}@{}}
\toprule\noalign{}
\begin{minipage}[b]{\linewidth}\raggedright
\textbf{Stakeholder}
\end{minipage} & \begin{minipage}[b]{\linewidth}\raggedright
\textbf{Rol en el negocio}
\end{minipage} & \begin{minipage}[b]{\linewidth}\raggedright
\textbf{Interés en el sistema}
\end{minipage} & \begin{minipage}[b]{\linewidth}\raggedright
\textbf{Poder / interés}
\end{minipage} & \begin{minipage}[b]{\linewidth}\raggedright
\textbf{Estrategia de participación}
\end{minipage} \\
\midrule\noalign{}
\endhead
\bottomrule\noalign{}
\endlastfoot
Dueños del taller & Propietarios; deciden la inversión y aprueban el
proyecto. & Reducir disputas por cobros, cuidar la rentabilidad y la
reputación, y no pagar licencias. & Alto poder -- alto interés &
Gestionar de cerca mediante entrevistas y revisiones de avance en cada
entrega, validación de requisitos y aprobación formal del alcance. \\
Gerente & Supervisa la operación y autoriza entregas a crédito a
flotillas. & Consultar tableros de control, controlar saldos y autorizar
créditos con respaldo. & Alto poder -- alto interés & Gestionar de cerca
mediante sesiones de definición de reglas de crédito y de reportes, y
pruebas de aceptación de su flujo. \\
Recepcionista & Registra ingresos, confirma citas y atiende a clientes.
& Reducir llamadas de seguimiento y tener el registro de recepción y
fotografías en un solo lugar. & Bajo poder -- alto interés & Mantener
informada e involucrada, participando en las pruebas de usabilidad,
capacitación y retroalimentación continua. \\
Jefe de mecánicos & Centraliza diagnóstico, cotización y asignación de
bahías en la tableta. & Flujos cortos, botones grandes y operación sin
conexión. & Alto poder -- alto interés & Gestionar de cerca los
prototipos en el entrono real (tableta, pruebas con manos sucias) y
ajustes por iteración. \\
Mecánicos de bahía y de apoyo & Realizan reparaciones, inspecciones y
cubrimientos de turnos. & Recibir órdenes claras y enfocarse en el
trabajo físico & Bajo poder -- alto interés & Mantener informado con
capacitaciones breves y un canal para reportar dificultades. \\
Encargado de almacén & Controla requisiciones, proveedores, costos e
inventario de consumibles. & Directorio de proveedores, registro de
costos y fechas, y tickets de compra por orden & Bajo poder -- alto
interés & Mantener informado e involucrado mediante entrevistas sobre su
flujo de compras y pruebas del módulo M6. \\
Clientes particulares & Dueños de autos, motocicletas y vehículos. & Ver
el estado de su vehículo, autorizar presupuestos y recibir un cobro
transparente & Bajo poder -- alto interés & Mantener informado con las
pruebas de usabilidad con algunos clientes y alguna encuesta breve de
satisfacción al finalizar el servicio. \\
Clientes de flotillas empresariales & Empresas con vehículos que pueden
recibir entrega a crédito. & Condiciones de crédito claras y seguimiento
de varios vehículos & Poder medio -- interés medio & Mantener
satisfechos al definir con el gerente las reglas del crédito y validar
el flujo con un representante. \\
Equipo de desarrollo & Diseñan, construyen y prueban la plataforma. &
Entregar un sistema que cumpla los requisitos y la norma ISO/IEC 25010 &
Alto poder -- alto interés & Gestionar de cerca con reuniones
periódicas, control de versiones y revisiones internas de calidad. \\
Docente evaluador & Revisa el documento y el producto en curso. & Que el
proyecto cumpla los criterios y entregas del curso & Alto poder -- alto
interés & Mantener satisfecho entregando a tiempo, pidiendo asesoría en
los avances y pidiendo aclaraciones de dudas. \\
\end{longtable}
}

\subsection{3.5 Proceso actual (as-is)}\label{proceso-actual-as-is}

\figura{figuras/proceso-as-is.png}{height=0.8\textheight}{3.1}{Diagrama de actividades del proceso actual (as-is) de
atención del taller, con los problemas identificados en cada etapa.
Elaboración propia}{docs/03-vision/proceso-as-is.puml}

El diagrama de actividades muestra el proceso actual de atención del
taller antes de implementar la solución tecnológica. Se organiza en
cuatro carriles (cliente, recepcionista, mecánico y registro manual) y
recorre el flujo completo: solicitud de cita, recepción del vehículo,
diagnóstico, cotización, autorización, reparación, entrega y archivo del
historial. Los recuadros rojos señalan los problemas detectados en cada
etapa, como el riesgo de errores en el registro de citas, los datos
incompletos en las órdenes en papel, la información dispersa e ilegible,
las demoras en la cotización manual, el exceso de llamadas y un
historial difícil de consultar.

\subsection{3.6 Estudio de viabilidad}\label{estudio-de-viabilidad}

{\def\LTcaptype{none} % do not increment counter
\begin{longtable}[]{@{}
  >{\raggedright\arraybackslash}p{(\linewidth - 4\tabcolsep) * \real{0.1708}}
  >{\raggedright\arraybackslash}p{(\linewidth - 4\tabcolsep) * \real{0.5513}}
  >{\raggedright\arraybackslash}p{(\linewidth - 4\tabcolsep) * \real{0.2779}}@{}}
\toprule\noalign{}
\begin{minipage}[b]{\linewidth}\raggedright
\textbf{Tipo}
\end{minipage} & \begin{minipage}[b]{\linewidth}\raggedright
\textbf{Análisis y evidencia}
\end{minipage} & \begin{minipage}[b]{\linewidth}\raggedright
\textbf{Conclusión (viable / viable con condiciones / no viable)}
\end{minipage} \\
\midrule\noalign{}
\endhead
\bottomrule\noalign{}
\endlastfoot
Técnica & La pila propuesta (Python con FastAPI, PostgreSQL y una PWA
con Service Workers e IndexedDB) es de código abierto y está bien
documentada. El mayor riesgo técnico es la operación sin conexión con
sincronización posterior (RNF-08, RNF-09, R-02), que se validará con un
prototipo en el Sprint 1. & Viable con condiciones: validar la
sincronización sin conexión en el Sprint 1. \\
Económica & Costo de licencias de \$0: todas las herramientas son de
código abierto o de planes educativos (RNF-07). El taller usa las
computadoras que ya tiene en recepción y almacén, y 1 o 2 tabletas para
el jefe de mecánicos. Los planes gratuitos de correo y almacenamiento
tienen límites (R-05, R-06). & Viable con condiciones: el taller aporta
las tabletas y el sistema respeta los límites de los planes
gratuitos. \\
Operativa & El personal tiene experiencia tecnológica baja a media
(4.3.1). Se atiende con flujos de pocos pasos y botones grandes (RNF-02,
RNF-03), capacitación breve y pruebas de usabilidad con el personal. Los
dueños y el gerente apoyan el proyecto porque reduce disputas y llamadas
(3.4). & Viable con condiciones: capacitar y probar con el personal
antes de la Demo final. \\
Legal & Se tratan datos de identificación y contacto de los clientes
(nombre, teléfono y correo) y datos de sus vehículos, incluidas
fotografías; no se tratan datos sensibles. Las finalidades y medidas de
seguridad (control de acceso por rol, HTTPS, tokens temporales y aviso
de privacidad) están en 4.8, conforme a la Ley Federal de Protección de
Datos Personales en Posesión de los Particulares (DOF 20-03-2025). En
desarrollo y pruebas solo se usan datos ficticios (RNF-05). & Viable \\
De calendario & El semestre tiene 7 sprints (S0 a S6) hasta la Demo
final de la semana 16. El backlog tiene 30 historias; el release v0.1.0
incluye el flujo núcleo y las historias Should o Could pueden pasar a la
Fase 2 si la velocidad del equipo no alcanza. & Viable con condiciones:
priorizar las historias Must. \\
\end{longtable}
}

\textbf{Conclusión general:} el proyecto es viable con condiciones:
validar en el Sprint 1 la operación sin conexión, que el taller aporte
las tabletas, respetar los límites de los planes gratuitos y priorizar
las historias Must del flujo núcleo.

\section{4. Especificación de requisitos de software
(SRS)}\label{especificaciuxf3n-de-requisitos-de-software-srs}

\begin{quote}
\emph{\textbf{Instrucción:} Estructura adaptada de ISO/IEC/IEEE
29148:2018. El SRS es lo que el cliente firma en la Review 2. Cada
requisito tiene un identificador único que se utiliza para la
trazabilidad.}
\end{quote}

\subsection{4.1 Propósito, alcance y audiencia del
SRS}\label{propuxf3sito-alcance-y-audiencia-del-srs}

\textbf{Propósito:} este SRS describe qué debe hacer el Sistema de
Órdenes de Servicio y Citas y con qué calidad, para que el cliente lo
valide y lo firme en la Review 2, y para que el equipo lo use como base
del diseño, la implementación y las pruebas.

\textbf{Alcance:} cubre el ciclo de atención del taller (citas,
recepción, diagnóstico, cotización y autorización, refacciones, cobro y
entrega) y el portal web del cliente, según el alcance de la sección
3.3.

\textbf{Audiencia:} los dueños y el gerente del taller (aprobación), el
personal que usará el sistema (validación de sus flujos), el equipo de
desarrollo y el docente evaluador.

\textbf{Convenciones:} cada requisito tiene un identificador único: RF
(funcional), RNF (no funcional), RN (regla de negocio), HU (historia de
usuario) y CU (caso de uso). Estos identificadores se usan en la matriz
de trazabilidad (5.3).

\subsection{4.2 Glosario del dominio}\label{glosario-del-dominio}

{\def\LTcaptype{none} % do not increment counter
\begin{longtable}[]{@{}
  >{\raggedright\arraybackslash}p{(\linewidth - 4\tabcolsep) * \real{0.2349}}
  >{\raggedright\arraybackslash}p{(\linewidth - 4\tabcolsep) * \real{0.5342}}
  >{\raggedright\arraybackslash}p{(\linewidth - 4\tabcolsep) * \real{0.2309}}@{}}
\toprule\noalign{}
\begin{minipage}[b]{\linewidth}\raggedright
\textbf{Término}
\end{minipage} & \begin{minipage}[b]{\linewidth}\raggedright
\textbf{Definición}
\end{minipage} & \begin{minipage}[b]{\linewidth}\raggedright
\textbf{Sinónimos que usa el cliente}
\end{minipage} \\
\midrule\noalign{}
\endhead
\bottomrule\noalign{}
\endlastfoot
Autorización & Aceptación de una cotización por el cliente; registra
fecha, hora, IP, monto y el código de verificación. & Aprobación, firma
digital \\
Bahía & Espacio de trabajo donde se repara un vehículo. El taller tiene
4 y cada una aloja un vehículo a la vez. & --- \\
Checklist de inspección & Lista de verificación del estado de
carrocería, interiores y accesorios al recibir el vehículo. & Lista de
inspección \\
Cita & Solicitud para atender un vehículo en una fecha y hora; queda
«Pendiente» hasta que la recepcionista la confirma. & --- \\
Código 2FA & Código de 6 dígitos, de un solo uso y vigente 5 minutos,
que se envía al correo del cliente para confirmar una autorización. &
Código de verificación, doble verificación \\
Compra & Adquisición de refacciones a un proveedor, respaldada por un
ticket. & Ticket de compra \\
Concepto & Renglón de una cotización: mano de obra, refacción o insumo.
Puede ser opcional. & Rubro \\
Cotización & Costo desglosado de mano de obra, refacciones e insumos de
una orden; vigente 5 días naturales. & Presupuesto \\
Diagnóstico & Descripción técnica de la falla y de las reparaciones
necesarias que registra el jefe de mecánicos. & --- \\
Enlace seguro & Liga web única que recibe el cliente para consultar su
vehículo y autorizar cotizaciones sin usuario ni contraseña. & Token de
acceso \\
Entrega a crédito & Entrega del vehículo con saldo pendiente; solo para
flotillas registradas y con autorización del gerente. & Crédito \\
Flotilla & Empresa cliente con varios vehículos que puede recibir
entrega a crédito. & Flotilla empresarial \\
Folio & Identificador único de una orden de servicio. & --- \\
Fotografía de recepción & Evidencia digital del estado del vehículo al
llegar. Son obligatorias cuatro: frente, trasera, costado izquierdo y
costado derecho. & --- \\
Garantía & Cobertura de 30 días sobre la mano de obra, contados desde la
entrega del vehículo. & --- \\
Insumo & Material consumible que se usa en la reparación, como aceite o
líquidos. & Consumible \\
Jefe de mecánicos & Líder técnico: diagnostica, cotiza, asigna bahías y
mecánicos y solicita refacciones desde la tableta. & Líder técnico \\
Mecánico responsable & Mecánico titular de una orden; otros mecánicos
pueden colaborar sin cambiar la titularidad. & Mecánico principal \\
Orden de servicio & Registro de todo el trabajo sobre un vehículo, desde
la recepción hasta la entrega. & Orden \\
Proveedor & Refaccionaria local a la que el almacén compra piezas. &
Refaccionaria \\
Recepción vehicular & Registro de la llegada del vehículo: kilometraje,
nivel de combustible, checklist y fotografías. & Check-in, registro de
entrada \\
Re-cotización & Nueva versión de la cotización cuando los trabajos
adicionales aumentan el monto más de 10 \%. & --- \\
Refacción & Pieza que se instala en el vehículo. & Pieza, repuesto \\
Saldo pendiente & Diferencia entre el total autorizado de la orden y lo
que el cliente ha pagado. & --- \\
Solicitud de refacción & Pedido de piezas que el jefe de mecánicos hace
al almacén para una orden. & Requisición \\
\end{longtable}
}

\subsection{4.3 Descripción general}\label{descripciuxf3n-general}

\subsubsection{4.3.1 Usuarios y sus
características}\label{usuarios-y-sus-caracteruxedsticas}

El sistema será utilizado por distintos perfiles dentro de la estructura
operativa del negocio y la parte gerencial del taller mecánico, además
de los clientes finales del negocio. La siguiente tabla detalla las
características, nivel técnico y contexto de uso:

{\def\LTcaptype{none} % do not increment counter
\begin{longtable}[]{@{}
  >{\raggedright\arraybackslash}p{(\linewidth - 6\tabcolsep) * \real{0.2138}}
  >{\raggedright\arraybackslash}p{(\linewidth - 6\tabcolsep) * \real{0.2990}}
  >{\raggedright\arraybackslash}p{(\linewidth - 6\tabcolsep) * \real{0.2138}}
  >{\raggedright\arraybackslash}p{(\linewidth - 6\tabcolsep) * \real{0.2734}}@{}}
\toprule\noalign{}
\begin{minipage}[b]{\linewidth}\raggedright
\textbf{Tipo de usuario}
\end{minipage} & \begin{minipage}[b]{\linewidth}\raggedright
\textbf{Tareas principales}
\end{minipage} & \begin{minipage}[b]{\linewidth}\raggedright
\textbf{Experiencia tecnológica}
\end{minipage} & \begin{minipage}[b]{\linewidth}\raggedright
\textbf{Dispositivo y contexto de uso}
\end{minipage} \\
\midrule\noalign{}
\endhead
\bottomrule\noalign{}
\endlastfoot
Recepcionista & Registro de clientes/vehículos, check-in con fotos,
apertura de órdenes, confirmación de citas web, cobranza y entrega. &
Media: Manejo fluido de navegación web, correo electrónico y paquetes de
oficina. & Computadora de escritorio en mostrador de recepción. Entorno
de oficina con teclado/mouse. \\
Jefe de Mecánicos (Líder Técnico) & Diagnóstico técnico, elaboración de
cotizaciones, asignación de bahías/mecánicos, requisición de refacciones
y sellado final. & Baja--Media: Familiarizado con smartphones
personales; poca experiencia con sistemas de gestión. & Tabletas
táctiles exclusivas de taller (1 o 2 dispositivos). Entorno con ruido,
polvo y manipulación con manos sucias/grasosas. \\
Mecánico en Bahía & Ejecución de reparaciones físicas, notificación de
hallazgos adicionales o solicitud de refacciones al Jefe de Mecánicos. &
Baja: Uso de teléfonos inteligentes para uso personal cotidiano. & Sin
dispositivo asignado directo; interactúa a través del Jefe de Mecánicos
para evitar deterioro de hardware \\
Encargado de Almacén & Atender solicitudes de piezas, consultar
directorio de proveedores locales, registrar compras/tickets e historial
de insumos. & Baja--Media\textbf{:} Uso básico de computadora para
consulta de inventarios. & Computadora de escritorio o terminal dedicada
en el área de almacén/refacciones. \\
Gerente (Dueño) & Monitoreo del Dashboard de KPIs (ocupación, ingresos),
aprobación de entregas a crédito a flotillas y excepciones. & Media:
Manejo de smartphones y computadoras para administración del negocio. &
Laptop personal o dispositivo móvil (smartphone/tableta) con acceso
remoto/local \\
Cliente & Solicitud de citas en línea, seguimiento en vivo del servicio
y aprobación/rechazo de cotizaciones mediante código 2FA. & Variable
(Baja a Alta): Desde usuarios sin conocimientos técnicos hasta usuarios
avanzados. & Dispositivo móvil personal (iOS/Android) vía navegador web
progresivo (PWA) sin instalar aplicaciones. \\
\end{longtable}
}

\subsubsection{4.3.2 Restricciones, supuestos y
dependencias}\label{restricciones-supuestos-y-dependencias}

\begin{enumerate}
\def\labelenumi{\arabic{enumi}.}
\item
  \textbf{Restricciones:}
\end{enumerate}

\begin{itemize}
\item
  Presupuesto y Licenciamiento: El desarrollo debe realizarse
  estrictamente con herramientas, lenguajes y frameworks de código
  abierto (\emph{Open Source} como Python y FastAPI), garantizando cero
  costos de licencias o rentas de software.
\item
  Infraestructura Física: El taller cuenta únicamente con 4 bahías de
  trabajo físicas (capacidad estricta de 1 vehículo por bahía) y un
  estacionamiento con capacidad máxima de 5 vehículos en espera.
\item
  Dispositivos Hardware: La captura digital en taller se limita a máximo
  2 tabletas táctiles centralizadas por el Jefe de Mecánicos para
  proteger la inversión y evitar daños.
\end{itemize}

\begin{itemize}
\item
  Alcance del Vehículo\textbf{:} La plataforma atenderá exclusivamente
  vehículos ligeros (automóviles particulares, motocicletas y
  camionetas), excluyendo maquinaria pesada o industrial.
\end{itemize}

\begin{enumerate}
\def\labelenumi{\arabic{enumi}.}
\setcounter{enumi}{1}
\item
  \textbf{Supuestos:}
\end{enumerate}

\begin{itemize}
\item
  Se asume que el taller mantiene activa una conexión de red local
  (Wi-Fi) para la sincronización entre computadoras de mostrador y
  tabletas de taller en condiciones normales.
\item
  Se asume que los clientes cuentan con un correo electrónico para
  recibir el enlace de seguimiento y los códigos de verificación de 6
  dígitos (2FA); el número celular se usa para el envío opcional del
  enlace por SMS (RF-02.3).
\item
  Se asume que los datos ingresados durante la fase de desarrollo y
  pruebas son 100\% ficticios para dar cumplimiento a la Ley Federal de
  Protección de Datos Personales.
\end{itemize}

\begin{enumerate}
\def\labelenumi{\arabic{enumi}.}
\setcounter{enumi}{2}
\item
  \textbf{Dependencias:}
\end{enumerate}

\begin{itemize}
\item
  Infraestructura Eléctrica y Red: Operación dependiente de baterías en
  laptops/tabletas y arquitectura \emph{Local-First / PWA} (Service
  Workers e IndexedDB) para permitir la captura de datos \emph{offline}
  durante cortes de luz o fallas de internet.
\item
  Proveedores Externos de Comunicación: Dependencia del servicio de
  correo electrónico para entregar enlaces y códigos 2FA al cliente, y
  de un servicio de SMS para el envío opcional del enlace.
\item
  Disponibilidad de Refaccionarias Locales: La adquisición de piezas
  fuera de stock depende del directorio de proveedores locales e
  historial de compras del almacén.
\end{itemize}

\subsection{4.4 Requisitos funcionales}\label{requisitos-funcionales}

{\def\LTcaptype{none} % do not increment counter
\begin{longtable}[]{@{}
  >{\raggedright\arraybackslash}p{(\linewidth - 8\tabcolsep) * \real{0.0962}}
  >{\raggedright\arraybackslash}p{(\linewidth - 8\tabcolsep) * \real{0.4444}}
  >{\raggedright\arraybackslash}p{(\linewidth - 8\tabcolsep) * \real{0.1174}}
  >{\raggedright\arraybackslash}p{(\linewidth - 8\tabcolsep) * \real{0.1497}}
  >{\raggedright\arraybackslash}p{(\linewidth - 8\tabcolsep) * \real{0.1923}}@{}}
\toprule\noalign{}
\begin{minipage}[b]{\linewidth}\raggedright
\textbf{ID}
\end{minipage} & \begin{minipage}[b]{\linewidth}\raggedright
\textbf{Requisito funcional}
\end{minipage} & \begin{minipage}[b]{\linewidth}\raggedright
\textbf{Módulo}
\end{minipage} & \begin{minipage}[b]{\linewidth}\raggedright
\textbf{Prioridad (MoSCoW)}
\end{minipage} & \begin{minipage}[b]{\linewidth}\raggedright
\textbf{Historias relacionadas}
\end{minipage} \\
\midrule\noalign{}
\endhead
\bottomrule\noalign{}
\endlastfoot
RF-01 & El sistema deberá permitir autenticar mediante credenciales
seguras al personal del taller según su rol (Gerente, Recepcionista,
Jefe de Mecánicos y Encargado de Almacén). & M1 -- Autenticación y
Seguridad & Must & HU-27 \\
RF-02.1 & El sistema deberá permitir al cliente consultar el estado de
su vehículo y autorizar cotizaciones mediante un enlace web seguro
enviado por correo electrónico o por SMS (RF-02.3), sin instalar
aplicaciones. & M1 -- Autenticación y Seguridad & Must & HU-02, HU-28 \\
RF-02.2 & El sistema deberá de expirar el enlace del cliente de manera
automática una vez que el vehículo sea entregado y saldado. & M1 --
Autenticación y Seguridad & Should & HU-02, HU-23 \\
RF-02.3 & El sistema deberá de permitir enviar el enlace del cliente
también por SMS. & M1 -- Autenticación y Seguridad & Could & HU-02 \\
RF-03.1 & El sistema deberá registrar, buscar y actualizar la
información de clientes y vehículos (placas, marca, modelo, año y VIN
(opcional)). & M3 -- Recepción E Historial Inicial & Must & HU-05,
HU-06 \\
RF-03.2 & El sistema deberá mantener el historial de servicios de cada
vehículo, con fotografías de recepción, diagnósticos técnicos y
garantías aplicables (30 días en mano de obra). & M7 -- Portal del
Cliente, Cobro y Gerencia & Should & HU-23, HU-24 \\
RF-04 & El sistema deberá registrar en la recepción el kilometraje, el
nivel de combustible y un checklist de carrocería, y exigir fotografías
digitales del estado inicial de la unidad. & M3 -- Recepción E Historial
Inicial & Must & HU-01, HU-25, HU-26 \\
RF-05 & El sistema deberá permitir abrir órdenes de servicio asociadas a
una de las 4 bahías y a un mecánico responsable, controlando los estados
operativos del taller. & M5 -- Taller, Bahías y Expediente Técnico &
Must & HU-07, HU-08, HU-09, HU-11 \\
RF-06 & El sistema deberá generar cotizaciones desglosadas (mano de obra
y refacciones/insumos) con vigencia de 5 días naturales y requerir firma
digital validada con un código enviado al correo del cliente, de forma
remota o presencial. & M4 Cotización y Contratos 2FA & Must & HU-02,
HU-12, HU-13, HU-14, HU-15 \\
RF-07 & El sistema deberá bloquear el avance del trabajo y exigir una
nueva autorización con verificación cuando los trabajos adicionales
incrementen el monto inicial en más de 10 \%. & M4 Cotización y
Contratos 2FA & Must & HU-16 \\
RF-08.1 & El sistema deberá permitir al jefe de mecánicos solicitar al
almacén las piezas que requieran los mecánicos, vinculadas a una orden
de servicio. & M6 -- Almacén y Refacciones & Should & HU-10, HU-19 \\
RF-08.2 & El sistema deberá permitir al encargado de almacén mantener un
directorio de proveedores locales (teléfonos o correos) y registrar
costos y fechas de llegada de las piezas. & M6 -- Almacén y Refacciones
& Should & HU-17, HU-18 \\
RF-08.3 & El sistema deberá de emitir un ticket e historial de compras
asociado a cada orden de servicio. & M6 -- Almacén y Refacciones & Could
& HU-18 \\
RF-09 & El sistema deberá ofrecer una agenda digital por día y bahía,
donde las citas solicitadas en línea quedan ``Pendientes'' hasta que la
recepcionista verifique disponibilidad y las confirme manualmente & M2
-- Citas y Solicitudes & Should & HU-03, HU-04, HU-06 \\
RF-10 & El sistema deberá impedir la entrega física del vehículo con
saldo pendiente, a menos que el gerente autorice la entrega a crédito
para vehículos de flotillas empresariales previamente registradas. & M7
-- Portal del Cliente, Cobro y Gerencia & Must & HU-21, HU-22, HU-23 \\
RF-11 & El sistema deberá ofrecer al gerente un panel con las órdenes
abiertas y su estado, la ocupación de las 4 bahías, los tiempos de
estancia y los ingresos acumulados. & M7 -- Portal del Cliente, Cobro y
Gerencia & Should & HU-29 \\
\end{longtable}
}

\subsection{4.5 Historias de usuario (Product
Backlog)}\label{historias-de-usuario-product-backlog}

{\def\LTcaptype{none} % do not increment counter
\begin{longtable}[]{@{}
  >{\raggedright\arraybackslash}p{(\linewidth - 10\tabcolsep) * \real{0.1389}}
  >{\raggedright\arraybackslash}p{(\linewidth - 10\tabcolsep) * \real{0.4124}}
  >{\raggedright\arraybackslash}p{(\linewidth - 10\tabcolsep) * \real{0.1068}}
  >{\raggedright\arraybackslash}p{(\linewidth - 10\tabcolsep) * \real{0.0855}}
  >{\raggedright\arraybackslash}p{(\linewidth - 10\tabcolsep) * \real{0.1282}}
  >{\raggedright\arraybackslash}p{(\linewidth - 10\tabcolsep) * \real{0.1282}}@{}}
\toprule\noalign{}
\begin{minipage}[b]{\linewidth}\raggedright
\textbf{ID / issue}
\end{minipage} & \begin{minipage}[b]{\linewidth}\raggedright
\textbf{Historia (Como\ldots{} quiero\ldots{} para\ldots)}
\end{minipage} & \begin{minipage}[b]{\linewidth}\raggedright
\textbf{Prioridad}
\end{minipage} & \begin{minipage}[b]{\linewidth}\raggedright
\textbf{Puntos}
\end{minipage} & \begin{minipage}[b]{\linewidth}\raggedright
\textbf{Release}
\end{minipage} & \begin{minipage}[b]{\linewidth}\raggedright
\textbf{Estado}
\end{minipage} \\
\midrule\noalign{}
\endhead
\bottomrule\noalign{}
\endlastfoot
HU-01 / \#{[}1{]} & Como recepcionista del taller, quiero registrar la
recepción del vehículo adjuntando kilometraje, nivel de combustible,
checklist de carrocería y fotografías del estado inicial, para contar
con un respaldo digital inalterable que evite reclamos infundados por
daños o faltantes al entregar la unidad. & Must & 3 & v0.1.0 &
Pendiente \\
HU-02 / \#{[}2{]} & Como cliente del taller, quiero recibir un enlace
web seguro con la cotización detallada (mano de obra y refacciones) y
botones claros de aprobación o rechazo, para revisar y autorizar los
trabajos desde mi teléfono sin necesidad de realizar llamadas ni acudir
físicamente al taller & Must & 3 & v0.1.0 & Pendiente \\
HU-03 / \#{[}3{]} & Como cliente del taller, quiero solicitar una cita
en línea seleccionando el día y la hora disponible, para apartar un
espacio sin necesidad de realizar llamadas telefónicas. & Must & 3 &
v0.1.0 & Pendiente \\
HU-04 / \#{[}4{]} & Como recepcionista del taller, quiero consultar la
agenda digital por día y por bahía, para verificar las solicitudes
pendientes de los clientes y realizar su confirmación manual. & Must & 3
& v0.1.0 & Pendiente \\
HU-05 / \#{[}5{]} & Como recepcionista del taller, quiero buscar a un
cliente o vehículo registrado mediante su número telefónico o placas,
para agilizar el proceso de recepción en visitas recurrentes. & Should &
2 & Fase 2 & Pendiente \\
HU-06 / \#{[}6{]} & Como recepcionista del taller, quiero registrar un
vehículo con cita previa seleccionando su folio de origen, para que el
sistema precargue automáticamente los datos del cliente y del auto. &
Should & 2 & v0.1.0 & Pendiente \\
HU-07 / \#{[}7{]} & Como recepcionista del taller, quiero ver la alerta
de "Sin Bahías Disponibles" cuando las 4 bahías físicas estén ocupadas,
para ubicar temporalmente el vehículo en el estacionamiento de
recepción. & Must & 3 & v0.1.0 & Pendiente \\
HU-08 / \#{[}8{]} & Como jefe de mecánicos, quiero asignar la orden de
servicio a una de las 4 bahías físicas disponibles y a un mecánico
responsable, para iniciar formalmente los trabajos en el taller. & Must
& 5 & v0.1.0 & Pendiente \\
HU-09 / \#{[}9{]} & Como jefe de mecánicos, quiero registrar el
diagnóstico técnico y las observaciones iniciales en la tableta, para
cambiar el estado de la orden a "En Diagnóstico". & Must & 3 & v0.1.0 &
Pendiente \\
HU-10 / \#{[}10{]} & Como jefe de mecánicos, quiero registrar las
refacciones e insumos que los mecánicos detecten durante la revisión,
para que el área de almacén gestione su cotización o pedido. & Should &
3 & v0.1.0 & Pendiente \\
HU-11 / \#{[}11{]} & Como jefe de mecánicos, quiero visualizar el
estatus actual de las 4 bahías de trabajo en una interfaz simplificada
para tabletas, para conocer rápidamente qué unidades están listas o en
proceso. & Should & 3 & Fase 2 & Pendiente \\
HU-12 / \#{[}12{]} & Como recepcionista o jefe de mecánicos, quiero
generar un presupuesto desglosado de mano de obra y refacciones con
vigencia de 5 días naturales, para enviarlo de forma transparente al
cliente. & Must & 5 & v0.1.0 & Pendiente \\
HU-13 / \#{[}13{]} & Como cliente del taller, quiero seleccionar o
deseleccionar servicios opcionales dentro de mi cotización antes de
aprobarla, para adaptar el presupuesto a mis necesidades financieras. &
Should & 3 & v0.1.0 & Pendiente \\
HU-14 / \#{[}14{]} & Como cliente del taller, quiero realizar la firma
digital de aprobación mediante un código de verificación de 2 pasos
enviado a mi correo, para autorizar los trabajos de forma segura ya sea
de forma presencial o remota. & Must & 5 & v0.1.0 & Pendiente \\
HU-15 / \#{[}15{]} & Como cliente del taller, quiero presionar un botón
de rechazo en la cotización digital, para cancelar la orden y coordinar
la recolección del vehículo sin autorizar reparaciones. & Should & 2 &
v0.1.0 & Pendiente \\
HU-16 / \#{[}16{]} & Como cliente del taller, quiero que se me pida una
nueva autorización con código 2FA cuando los costos imprevistos superen
el 10 \% del presupuesto autorizado, para no pagar trabajos que no
aprobé. & Must & 5 & v0.1.0 & Pendiente \\
HU-17 / \#{[}17{]} & Como encargado de almacén, quiero consultar y
administrar un directorio de proveedores locales con teléfonos y correos
electrónicos, para solicitar refacciones especiales rápidamente. &
Should & 3 & v0.1.0 & Pendiente \\
HU-18 / \#{[}18{]} & Como encargado de almacén, quiero registrar los
costos, fechas de llegada y tickets de compra de las refacciones
asociadas a una orden, para mantener un control financiero estricto. &
Should & 3 & v0.1.0 & Pendiente \\
HU-19 / \#{[}19{]} & Como jefe de mecánicos, quiero verificar el stock
de refacciones comunes en el inventario interno, para saber si se
dispone de la pieza de forma inmediata. & Could & 2 & Fase 2 &
Pendiente \\
HU-20 / \#{[}20{]} & Como recepcionista, quiero recibir una notificación
automática cuando los trabajos del vehículo finalicen, para avisar al
cliente que su unidad está lista para entrega. & Should & 2 & v0.1.0 &
Pendiente \\
HU-21 / \#{[}21{]} & Como gerente del taller, quiero que el sistema
impida registrar la entrega de un vehículo cuya orden tenga saldo
pendiente, para evitar pérdidas financieras al taller. & Must & 3 &
v0.1.0 & Pendiente \\
HU-22 / \#{[}22{]} & Como gerente del taller, quiero autorizar de forma
excepcional la entrega a crédito únicamente para vehículos de flotillas
empresariales registradas, para cumplir con los convenios comerciales. &
Should & 3 & v0.1.0 & Pendiente \\
HU-23 / \#{[}23{]} & Como recepcionista, quiero registrar el pago total
del servicio y cerrar formalmente la orden, para liberar la bahía de
trabajo y archivar el historial del vehículo. & Must & 3 & v0.1.0 &
Pendiente \\
HU-24 / \#{[}24{]} & Como cliente, quiero consultar la vigencia de la
garantía de 30 días en mano de obra asociada a mi servicio anterior,
para reclamar ajustes en caso de incidencias. & Could & 2 & Fase 2 &
Pendiente \\
HU-25 / \#{[}25{]} & Como recepcionista o jefe de mecánicos, quiero
registrar la llegada de autos, tomar fotografías y hacer diagnósticos
sin conexión a internet usando la batería del equipo, para que la
operación no se detenga ante fallas eléctricas. & Must & 8 & v0.1.0 &
Pendiente \\
HU-26 / \#{[}26{]} & Como recepcionista o jefe de mecánicos, quiero que
los registros y fotografías que capturé sin conexión se sincronicen
automáticamente cuando vuelva la red, para no perder ni duplicar
información. & Must & 8 & v0.1.0 & Pendiente \\
HU-27 / \#{[}27{]} & Como gerente del taller, quiero dar de alta al
personal y asignarle su rol (Gerente, Recepcionista, Jefe de Mecánicos o
Encargado de Almacén), para que cada quien acceda solo a la información
que le corresponde. & Must & 3 & v0.1.0 & Pendiente \\
HU-28 / \#{[}28{]} & Como cliente, quiero visualizar la plataforma web
desde mi teléfono celular de forma responsiva y sin instalar
aplicaciones adicionales, para consultar mi estatus con facilidad. &
Must & 5 & v0.1.0 & Pendiente \\
HU-29 / \#{[}29{]} & Como gerente, quiero generar reportes mensuales de
ingresos, servicios realizados y tiempos de ocupación de las 4 bahías,
para evaluar el rendimiento general del taller. & Could & 5 & Fase 2 &
Pendiente \\
HU-30 / \#{[}30{]} & Como cliente del taller, quiero conocer el aviso de
privacidad antes de proporcionar mis datos, para saber para qué se
usarán conforme a la LFPDPPP. & Must & 3 & v0.1.0 & Pendiente \\
\end{longtable}
}

Detalle de las historias de usuario

\textbf{HU-01: Recepción e Inspección con Evidencia Fotográfica}

{\def\LTcaptype{none} % do not increment counter
\begin{longtable}[]{@{}
  >{\raggedright\arraybackslash}p{(\linewidth - 2\tabcolsep) * \real{0.2778}}
  >{\raggedright\arraybackslash}p{(\linewidth - 2\tabcolsep) * \real{0.7222}}@{}}
\toprule\noalign{}
\begin{minipage}[b]{\linewidth}\raggedright
\textbf{Campo}
\end{minipage} & \begin{minipage}[b]{\linewidth}\raggedright
\textbf{Contenido}
\end{minipage} \\
\midrule\noalign{}
\endhead
\bottomrule\noalign{}
\endlastfoot
ID e issue & HU-01 / \#{[}1{]} \\
Historia & Como recepcionista del taller, quiero registrar la recepción
del vehículo adjuntando kilometraje, nivel de combustible, checklist de
carrocería y fotografías del estado inicial, para contar con un respaldo
digital inalterable que evite reclamos infundados por daños o faltantes
al entregar la unidad. \\
Criterio de aceptación 1 & Dado que la recepcionista se encuentra en el
módulo de recepción vehicular, cuando intenta registrar una unidad
nueva, entonces el sistema exige obligatoriamente adjuntar al menos 4
fotografías (frente, parte trasera, costado izquierdo y costado
derecho). \\
Criterio de aceptación 2 & Dado que se capturan las evidencias, cuando
se finaliza el formulario de inspección, entonces el sistema genera una
firma visual o sello digital con fecha y hora exacta de recepción,
bloqueando el avance al estado "En Diagnóstico" si falta completar algún
punto del checklist obligatorio. \\
Reglas de negocio & RN-12 (Evidencia fotográfica de recepción), RN-07
(Clasificación de vehículos autorizados). \\
Wireframe & Pendiente (sección 4.10) \\
Casos de prueba & Pendiente (sección 5.2) \\
\end{longtable}
}

\textbf{HU-02: Autorización Digital de Cotización por Parte del Cliente}

{\def\LTcaptype{none} % do not increment counter
\begin{longtable}[]{@{}
  >{\raggedright\arraybackslash}p{(\linewidth - 2\tabcolsep) * \real{0.2778}}
  >{\raggedright\arraybackslash}p{(\linewidth - 2\tabcolsep) * \real{0.7222}}@{}}
\toprule\noalign{}
\begin{minipage}[b]{\linewidth}\raggedright
\textbf{Campo}
\end{minipage} & \begin{minipage}[b]{\linewidth}\raggedright
\textbf{Contenido}
\end{minipage} \\
\midrule\noalign{}
\endhead
\bottomrule\noalign{}
\endlastfoot
ID e issue & HU-02 / \#{[}2{]} \\
Historia & Como cliente del taller, quiero recibir un enlace web seguro
con la cotización detallada (mano de obra y refacciones) y botones
claros de aprobación o rechazo, para revisar y autorizar los trabajos
desde mi teléfono sin necesidad de realizar llamadas ni acudir
físicamente al taller \\
Criterio de aceptación 1 & Dado que la cotización ha sido generada por
el taller, cuando el sistema envía el aviso, entonces el cliente recibe
un enlace por correo electrónico (o por SMS cuando se habilite RF-02.3)
que no requiere registro previo ni contraseña, operando mediante un
token temporal firmado. \\
Criterio de aceptación 2 & Dado que el cliente revisa el desglose en su
dispositivo móvil, cuando presiona el botón "Autorizar" e ingresa un
código 2FA válido (HU-14), entonces el sistema registra de forma
inalterable la dirección IP, fecha, hora y el monto final aceptado,
cambiando automáticamente el estado de la orden a "Autorizada" y
notificando al mecánico responsable. \\
Reglas de negocio & RN-01 (Expiración y vigencia de enlaces), RN-02
(Doble verificación en autorizaciones), RN-11 (Vigencia de
cotizaciones). \\
Wireframe & Pendiente (sección 4.10) \\
Casos de prueba & Pendiente (sección 5.2) \\
\end{longtable}
}

\textbf{HU-03: Solicitud de Cita en Línea}

{\def\LTcaptype{none} % do not increment counter
\begin{longtable}[]{@{}
  >{\raggedright\arraybackslash}p{(\linewidth - 2\tabcolsep) * \real{0.2778}}
  >{\raggedright\arraybackslash}p{(\linewidth - 2\tabcolsep) * \real{0.7222}}@{}}
\toprule\noalign{}
\begin{minipage}[b]{\linewidth}\raggedright
\textbf{Campo}
\end{minipage} & \begin{minipage}[b]{\linewidth}\raggedright
\textbf{Contenido}
\end{minipage} \\
\midrule\noalign{}
\endhead
\bottomrule\noalign{}
\endlastfoot
ID e issue & HU-03 / \#{[}3{]} \\
Historia & Como cliente del taller, quiero solicitar una cita en línea
seleccionando el día y la hora disponible, para apartar un espacio sin
necesidad de realizar llamadas telefónicas. \\
Criterio de aceptación 1 & Dado que el cliente accede al portal web,
cuando selecciona fecha, hora y tipo de servicio, entonces el sistema
valida que exista disponibilidad en la agenda del taller. \\
Criterio de aceptación 2 & Dado que se envía la solicitud, cuando se
registra con éxito, entonces queda en estado "Pendiente" a la espera de
la confirmación manual de la recepcionista. \\
Reglas de negocio & RN-08 (Confirmación manual de citas web). \\
Wireframe & Pendiente (sección 4.10) \\
Casos de prueba & Pendiente (sección 5.2) \\
\end{longtable}
}

\textbf{HU-04: Consulta y Confirmación Manual de Citas}

{\def\LTcaptype{none} % do not increment counter
\begin{longtable}[]{@{}
  >{\raggedright\arraybackslash}p{(\linewidth - 2\tabcolsep) * \real{0.2778}}
  >{\raggedright\arraybackslash}p{(\linewidth - 2\tabcolsep) * \real{0.7222}}@{}}
\toprule\noalign{}
\begin{minipage}[b]{\linewidth}\raggedright
\textbf{Campo}
\end{minipage} & \begin{minipage}[b]{\linewidth}\raggedright
\textbf{Contenido}
\end{minipage} \\
\midrule\noalign{}
\endhead
\bottomrule\noalign{}
\endlastfoot
ID e issue & HU-04 / \#{[}4{]} \\
Historia & Como recepcionista del taller, quiero consultar la agenda
digital por día y por bahía, para verificar las solicitudes pendientes
de los clientes y realizar su confirmación manual. \\
Criterio de aceptación 1 & Dado que la recepcionista accede al módulo de
agenda, cuando filtra los registros por el día actual, entonces el
sistema despliega todas las citas con estatus "Pendiente". \\
Criterio de aceptación 2 & Dado que la recepcionista revisa una cita
solicitada, cuando presiona el botón "Confirmar", entonces el estatus
cambia a "Confirmada" y se habilita la asignación de la bahía
correspondiente. \\
Reglas de negocio & RN-08 (Confirmación manual de citas web), RN-05
(Capacidad física del taller). \\
Wireframe & Pendiente (sección 4.10) \\
Casos de prueba & Pendiente (sección 5.2) \\
\end{longtable}
}

\textbf{HU-05: Búsqueda Rápida de Clientes y Vehículos}

{\def\LTcaptype{none} % do not increment counter
\begin{longtable}[]{@{}
  >{\raggedright\arraybackslash}p{(\linewidth - 2\tabcolsep) * \real{0.2778}}
  >{\raggedright\arraybackslash}p{(\linewidth - 2\tabcolsep) * \real{0.7222}}@{}}
\toprule\noalign{}
\begin{minipage}[b]{\linewidth}\raggedright
\textbf{Campo}
\end{minipage} & \begin{minipage}[b]{\linewidth}\raggedright
\textbf{Contenido}
\end{minipage} \\
\midrule\noalign{}
\endhead
\bottomrule\noalign{}
\endlastfoot
ID e issue & HU-05 / \#{[}5{]} \\
Historia & Como recepcionista del taller, quiero buscar a un cliente o
vehículo registrado mediante su número telefónico o placas, para
agilizar el proceso de recepción en visitas recurrentes. \\
Criterio de aceptación 1 & Dado que la recepcionista se encuentra en el
módulo de búsqueda, cuando ingresa los dígitos de las placas o el
teléfono, entonces el sistema filtra instantáneamente la base de datos
de clientes recurrentes. \\
Criterio de aceptación 2 & Dado que se localiza el expediente, cuando el
usuario lo selecciona, entonces el sistema precarga los datos generales
del cliente y del auto para evitar duplicidad de registros. \\
Reglas de negocio & Sin regla de negocio específica; ver RF-03.1. \\
Wireframe & Pendiente (sección 4.10) \\
Casos de prueba & Pendiente (sección 5.2) \\
\end{longtable}
}

\textbf{HU-06: Registro de Vehículos con Cita Previa}

{\def\LTcaptype{none} % do not increment counter
\begin{longtable}[]{@{}
  >{\raggedright\arraybackslash}p{(\linewidth - 2\tabcolsep) * \real{0.2778}}
  >{\raggedright\arraybackslash}p{(\linewidth - 2\tabcolsep) * \real{0.7222}}@{}}
\toprule\noalign{}
\begin{minipage}[b]{\linewidth}\raggedright
\textbf{Campo}
\end{minipage} & \begin{minipage}[b]{\linewidth}\raggedright
\textbf{Contenido}
\end{minipage} \\
\midrule\noalign{}
\endhead
\bottomrule\noalign{}
\endlastfoot
ID e issue & HU-06 / \#{[}6{]} \\
Historia & Como recepcionista del taller, quiero registrar un vehículo
con cita previa seleccionando su folio de origen, para que el sistema
precargue automáticamente los datos del cliente y del auto. \\
Criterio de aceptación 1 & Dado que el cliente llega al taller con cita
agendada, cuando la recepcionista introduce el número de folio de la
cita, entonces el sistema recupera la información asociada. \\
Criterio de aceptación 2 & Dado que los datos están precargados, cuando
se procede con el ingreso, entonces el sistema vincula la cita
directamente con la nueva orden de servicio abierta. \\
Reglas de negocio & RN-08 (Confirmación manual de citas web). \\
Wireframe & Pendiente (sección 4.10) \\
Casos de prueba & Pendiente (sección 5.2) \\
\end{longtable}
}

\textbf{HU-07: Alerta de Bahías Ocupadas}

{\def\LTcaptype{none} % do not increment counter
\begin{longtable}[]{@{}
  >{\raggedright\arraybackslash}p{(\linewidth - 2\tabcolsep) * \real{0.2778}}
  >{\raggedright\arraybackslash}p{(\linewidth - 2\tabcolsep) * \real{0.7222}}@{}}
\toprule\noalign{}
\begin{minipage}[b]{\linewidth}\raggedright
\textbf{Campo}
\end{minipage} & \begin{minipage}[b]{\linewidth}\raggedright
\textbf{Contenido}
\end{minipage} \\
\midrule\noalign{}
\endhead
\bottomrule\noalign{}
\endlastfoot
ID e issue & HU-07 / \#{[}7{]} \\
Historia & Como recepcionista del taller, quiero ver la alerta de "Sin
Bahías Disponibles" cuando las 4 bahías físicas estén ocupadas, para
ubicar temporalmente el vehículo en el estacionamiento de recepción. \\
Criterio de aceptación 1 & Dado que las 4 bahías físicas del taller
registran una orden activa, cuando la recepcionista intenta registrar la
entrada de un nuevo vehículo, entonces el sistema emite una advertencia
visual en pantalla indicando "Sin Bahías Disponibles". \\
Criterio de aceptación 2 & Dado que se activa la alerta de saturación,
cuando el usuario acepta la notificación, entonces el sistema permite
guardar el registro temporal asignando la unidad al estacionamiento de
espera. \\
Reglas de negocio & RN-05 (Capacidad física del taller). \\
Wireframe & Pendiente (sección 4.10) \\
Casos de prueba & Pendiente (sección 5.2) \\
\end{longtable}
}

\textbf{HU-08: Asignación de Bahía y Mecánico}

{\def\LTcaptype{none} % do not increment counter
\begin{longtable}[]{@{}
  >{\raggedright\arraybackslash}p{(\linewidth - 2\tabcolsep) * \real{0.2778}}
  >{\raggedright\arraybackslash}p{(\linewidth - 2\tabcolsep) * \real{0.7222}}@{}}
\toprule\noalign{}
\begin{minipage}[b]{\linewidth}\raggedright
\textbf{Campo}
\end{minipage} & \begin{minipage}[b]{\linewidth}\raggedright
\textbf{Contenido}
\end{minipage} \\
\midrule\noalign{}
\endhead
\bottomrule\noalign{}
\endlastfoot
ID e issue & HU-08 / \#{[}8{]} \\
Historia & Como jefe de mecánicos, quiero asignar la orden de servicio a
una de las 4 bahías físicas disponibles y a un mecánico responsable,
para iniciar formalmente los trabajos en el taller. \\
Criterio de aceptación 1 & Dado que la orden se encuentra abierta,
cuando el jefe de mecánicos selecciona una de las 4 bahías físicas
libres, entonces el sistema actualiza el estado de la bahía a
"Ocupada". \\
Criterio de aceptación 2 & Dado que las 4 bahías están ocupadas, cuando
se intenta una asignación, entonces el sistema emite una alerta
impidiendo el cruce de unidades en la misma bahía. \\
Reglas de negocio & RN-05 (Capacidad física del taller), RN-09
(Asignación y colaboración de mecánicos). \\
Wireframe & Pendiente (sección 4.10) \\
Casos de prueba & Pendiente (sección 5.2) \\
\end{longtable}
}

\textbf{HU-09: Registro de Diagnóstico en Tableta}

{\def\LTcaptype{none} % do not increment counter
\begin{longtable}[]{@{}
  >{\raggedright\arraybackslash}p{(\linewidth - 2\tabcolsep) * \real{0.2778}}
  >{\raggedright\arraybackslash}p{(\linewidth - 2\tabcolsep) * \real{0.7222}}@{}}
\toprule\noalign{}
\begin{minipage}[b]{\linewidth}\raggedright
\textbf{Campo}
\end{minipage} & \begin{minipage}[b]{\linewidth}\raggedright
\textbf{Contenido}
\end{minipage} \\
\midrule\noalign{}
\endhead
\bottomrule\noalign{}
\endlastfoot
ID e issue & HU-09 / \#{[}9{]} \\
Historia & Como jefe de mecánicos, quiero registrar el diagnóstico
técnico y las observaciones iniciales en la tableta, para cambiar el
estado de la orden a "En Diagnóstico". \\
Criterio de aceptación 1 & Dado que el jefe de mecánicos se encuentra
autenticado en la tableta, cuando selecciona una orden con vehículo en
bahía, entonces visualiza un formulario táctil optimizado con campos
amplios para el dictamen técnico. \\
Criterio de aceptación 2 & Dado que se ingresan las observaciones de la
revisión, cuando se presiona "Guardar Diagnóstico", entonces el sistema
actualiza automáticamente el estado de la orden a "En Diagnóstico" y
habilita la creación de cotizaciones. \\
Reglas de negocio & RN-12 (Evidencia fotográfica de recepción).
Usabilidad en tableta: RNF-02 y RNF-03. \\
Wireframe & Pendiente (sección 4.10) \\
Casos de prueba & Pendiente (sección 5.2) \\
\end{longtable}
}

\textbf{HU-10: Solicitud de Refacciones e Insumos}

{\def\LTcaptype{none} % do not increment counter
\begin{longtable}[]{@{}
  >{\raggedright\arraybackslash}p{(\linewidth - 2\tabcolsep) * \real{0.2778}}
  >{\raggedright\arraybackslash}p{(\linewidth - 2\tabcolsep) * \real{0.7222}}@{}}
\toprule\noalign{}
\begin{minipage}[b]{\linewidth}\raggedright
\textbf{Campo}
\end{minipage} & \begin{minipage}[b]{\linewidth}\raggedright
\textbf{Contenido}
\end{minipage} \\
\midrule\noalign{}
\endhead
\bottomrule\noalign{}
\endlastfoot
ID e issue & HU-10 / \#{[}10{]} \\
Historia & Como jefe de mecánicos, quiero registrar las refacciones e
insumos que los mecánicos detecten durante la revisión, para que el área
de almacén gestione su cotización o pedido. \\
Criterio de aceptación 1 & Dado que un mecánico reporta piezas
necesarias en la bahía, cuando el jefe de mecánicos las añade desde la
tableta y las vincula a la orden, entonces la solicitud se envía
automáticamente al módulo de almacén. \\
Criterio de aceptación 2 & Dado que el almacén recibe el requerimiento,
cuando se procesa, entonces el sistema refleja el estatus de la
refacción como "Solicitada" o "Pedida". \\
Reglas de negocio & RN-04 (Re-cotización obligatoria por sobrecostos),
RN-06 (Reubicación por espera de refacciones). \\
Wireframe & Pendiente (sección 4.10) \\
Casos de prueba & Pendiente (sección 5.2) \\
\end{longtable}
}

\textbf{HU-11: Visualización del Estatus de Bahías}

{\def\LTcaptype{none} % do not increment counter
\begin{longtable}[]{@{}
  >{\raggedright\arraybackslash}p{(\linewidth - 2\tabcolsep) * \real{0.2778}}
  >{\raggedright\arraybackslash}p{(\linewidth - 2\tabcolsep) * \real{0.7222}}@{}}
\toprule\noalign{}
\begin{minipage}[b]{\linewidth}\raggedright
\textbf{Campo}
\end{minipage} & \begin{minipage}[b]{\linewidth}\raggedright
\textbf{Contenido}
\end{minipage} \\
\midrule\noalign{}
\endhead
\bottomrule\noalign{}
\endlastfoot
ID e issue & HU-11 / \#{[}11{]} \\
Historia & Como jefe de mecánicos, quiero visualizar el estatus actual
de las 4 bahías de trabajo en una interfaz simplificada para tabletas,
para conocer rápidamente qué unidades están listas o en proceso. \\
Criterio de aceptación 1 & Dado que el jefe de mecánicos accede al
tablero de control, cuando visualiza la sección de bahías, entonces
identifica mediante códigos de color las unidades activas, en
diagnóstico o listas. \\
Criterio de aceptación 2 & Dado que se actualiza el estado de una
reparación, cuando cambia de fase, entonces el tablero refresca la
información en tiempo real para optimizar la supervisión del taller. \\
Reglas de negocio & RN-05 (Capacidad física del taller). Usabilidad en
tableta: RNF-02. \\
Wireframe & Pendiente (sección 4.10) \\
Casos de prueba & Pendiente (sección 5.2) \\
\end{longtable}
}

\textbf{HU-12: Generación de Presupuesto Desglosado}

{\def\LTcaptype{none} % do not increment counter
\begin{longtable}[]{@{}
  >{\raggedright\arraybackslash}p{(\linewidth - 2\tabcolsep) * \real{0.2778}}
  >{\raggedright\arraybackslash}p{(\linewidth - 2\tabcolsep) * \real{0.7222}}@{}}
\toprule\noalign{}
\begin{minipage}[b]{\linewidth}\raggedright
\textbf{Campo}
\end{minipage} & \begin{minipage}[b]{\linewidth}\raggedright
\textbf{Contenido}
\end{minipage} \\
\midrule\noalign{}
\endhead
\bottomrule\noalign{}
\endlastfoot
ID e issue & HU-12 / \#{[}12{]} \\
Historia & Como recepcionista o jefe de mecánicos, quiero generar un
presupuesto desglosado de mano de obra y refacciones con vigencia de 5
días naturales, para enviarlo de forma transparente al cliente. \\
Criterio de aceptación 1 & Dado que el diagnóstico ha concluido, cuando
el personal selecciona la opción de cotizar, entonces el sistema
presenta apartados separados para costos de mano de obra y
refacciones. \\
Criterio de aceptación 2 & Dado que se emite la cotización, entonces el
sistema calcula automáticamente el total con IVA e integra una fecha de
expiración exactamente a los 5 días naturales posteriores a su
emisión. \\
Reglas de negocio & RN-11 (Vigencia de cotizaciones), RN-03 (Selección
granular de rubros). \\
Wireframe & Pendiente (sección 4.10) \\
Casos de prueba & Pendiente (sección 5.2) \\
\end{longtable}
}

\textbf{HU-13: Selección de Servicios Opcionales en Cotización}

{\def\LTcaptype{none} % do not increment counter
\begin{longtable}[]{@{}
  >{\raggedright\arraybackslash}p{(\linewidth - 2\tabcolsep) * \real{0.2778}}
  >{\raggedright\arraybackslash}p{(\linewidth - 2\tabcolsep) * \real{0.7222}}@{}}
\toprule\noalign{}
\begin{minipage}[b]{\linewidth}\raggedright
\textbf{Campo}
\end{minipage} & \begin{minipage}[b]{\linewidth}\raggedright
\textbf{Contenido}
\end{minipage} \\
\midrule\noalign{}
\endhead
\bottomrule\noalign{}
\endlastfoot
ID e issue & HU-13 / \#{[}13{]} \\
Historia & Como cliente del taller, quiero seleccionar o deseleccionar
servicios opcionales dentro de mi cotización antes de aprobarla, para
adaptar el presupuesto a mis necesidades financieras. \\
Criterio de aceptación 1 & Dado que el cliente revisa el presupuesto
desglosado en su celular, cuando marca o desmarca las casillas de
servicios opcionales o sugeridos, entonces el sistema recalcula el costo
total en tiempo real. \\
Criterio de aceptación 2 & Dado que se ajustan los opcionales, cuando el
cliente procede a la firma, entonces el sistema guarda únicamente los
conceptos seleccionados como parte oficial de la orden. \\
Reglas de negocio & RN-03 (Selección granular de rubros). \\
Wireframe & Pendiente (sección 4.10) \\
Casos de prueba & Pendiente (sección 5.2) \\
\end{longtable}
}

\textbf{HU-14: Firma Digital y Autorización por 2FA}

{\def\LTcaptype{none} % do not increment counter
\begin{longtable}[]{@{}
  >{\raggedright\arraybackslash}p{(\linewidth - 2\tabcolsep) * \real{0.2778}}
  >{\raggedright\arraybackslash}p{(\linewidth - 2\tabcolsep) * \real{0.7222}}@{}}
\toprule\noalign{}
\begin{minipage}[b]{\linewidth}\raggedright
\textbf{Campo}
\end{minipage} & \begin{minipage}[b]{\linewidth}\raggedright
\textbf{Contenido}
\end{minipage} \\
\midrule\noalign{}
\endhead
\bottomrule\noalign{}
\endlastfoot
ID e issue & HU-14 / \#{[}14{]} \\
Historia & Como cliente del taller, quiero realizar la firma digital de
aprobación mediante un código de verificación de 2 pasos enviado a mi
correo, para autorizar los trabajos de forma segura ya sea de forma
presencial o remota. \\
Criterio de aceptación 1 & Dado que el cliente accede al enlace web de
cotización, when solicita el código de autorización, entonces el sistema
envía un token numérico de 6 dígitos al correo electrónico registrado
del cliente. \\
Criterio de aceptación 2 & Dado que el cliente ingresa el código
recibido en la interfaz, cuando presiona "Autorizar Presupuesto",
entonces el sistema valida el token, registra la dirección IP y la marca
de tiempo, cambiando la orden a "Autorizada". \\
Reglas de negocio & RN-02 (Doble verificación en autorizaciones). No
repudio: RNF-04. \\
Wireframe & Pendiente (sección 4.10) \\
Casos de prueba & Pendiente (sección 5.2) \\
\end{longtable}
}

\textbf{HU-15: Rechazo de Cotización y Cancelación}

{\def\LTcaptype{none} % do not increment counter
\begin{longtable}[]{@{}
  >{\raggedright\arraybackslash}p{(\linewidth - 2\tabcolsep) * \real{0.2778}}
  >{\raggedright\arraybackslash}p{(\linewidth - 2\tabcolsep) * \real{0.7222}}@{}}
\toprule\noalign{}
\begin{minipage}[b]{\linewidth}\raggedright
\textbf{Campo}
\end{minipage} & \begin{minipage}[b]{\linewidth}\raggedright
\textbf{Contenido}
\end{minipage} \\
\midrule\noalign{}
\endhead
\bottomrule\noalign{}
\endlastfoot
ID e issue & HU-15 / \#{[}15{]} \\
Historia & Como cliente del taller, quiero presionar un botón de rechazo
en la cotización digital, para cancelar la orden y coordinar la
recolección del vehículo sin autorizar reparaciones. \\
Criterio de aceptación 1 & Dado que el cliente evalúa la cotización en
su dispositivo, cuando presiona el botón "Rechazar", entonces el sistema
solicita una confirmación de la acción. \\
Criterio de aceptación 2 & Dado que se confirma el rechazo, cuando se
procesa, entonces la orden cambia a estatus "Rechazada" y notifica a
recepción para coordinar la entrega física del auto sin costo de
reparación. \\
Reglas de negocio & RN-02 (Doble verificación en autorizaciones). Flujo
de rechazo: CU-11 (A1). \\
Wireframe & Pendiente (sección 4.10) \\
Casos de prueba & Pendiente (sección 5.2) \\
\end{longtable}
}

\textbf{HU-16: Re-cotización por Incremento \textgreater10\%}

{\def\LTcaptype{none} % do not increment counter
\begin{longtable}[]{@{}
  >{\raggedright\arraybackslash}p{(\linewidth - 2\tabcolsep) * \real{0.2778}}
  >{\raggedright\arraybackslash}p{(\linewidth - 2\tabcolsep) * \real{0.7222}}@{}}
\toprule\noalign{}
\begin{minipage}[b]{\linewidth}\raggedright
\textbf{Campo}
\end{minipage} & \begin{minipage}[b]{\linewidth}\raggedright
\textbf{Contenido}
\end{minipage} \\
\midrule\noalign{}
\endhead
\bottomrule\noalign{}
\endlastfoot
ID e issue & HU-16 / \#{[}16{]} \\
Historia & Como cliente del taller, quiero que se me pida una nueva
autorización con código 2FA cuando los costos imprevistos superen el 10
\% del presupuesto autorizado, para no pagar trabajos que no aprobé. \\
Criterio de aceptación 1 & Dado que surgen refacciones imprevistas,
cuando el incremento total supera el 10\% del presupuesto aprobado,
entonces el sistema bloquea el flujo y solicita una nueva autorización
con código 2FA. \\
Criterio de aceptación 2 & Dado que el cliente recibe la notificación
del sobrecosto, cuando aprueba mediante el código digital, entonces se
actualiza el monto total e historial de la orden. \\
Reglas de negocio & RN-04 (Re-cotización obligatoria por sobrecostos),
RN-02 (Doble verificación en autorizaciones). \\
Wireframe & Pendiente (sección 4.10) \\
Casos de prueba & Pendiente (sección 5.2) \\
\end{longtable}
}

\textbf{HU-17: Directorio y Gestión de Proveedores Locales}

{\def\LTcaptype{none} % do not increment counter
\begin{longtable}[]{@{}
  >{\raggedright\arraybackslash}p{(\linewidth - 2\tabcolsep) * \real{0.2778}}
  >{\raggedright\arraybackslash}p{(\linewidth - 2\tabcolsep) * \real{0.7222}}@{}}
\toprule\noalign{}
\begin{minipage}[b]{\linewidth}\raggedright
\textbf{Campo}
\end{minipage} & \begin{minipage}[b]{\linewidth}\raggedright
\textbf{Contenido}
\end{minipage} \\
\midrule\noalign{}
\endhead
\bottomrule\noalign{}
\endlastfoot
ID e issue & HU-17 / \#{[}17{]} \\
Historia & Como encargado de almacén, quiero consultar y administrar un
directorio de proveedores locales con teléfonos y correos electrónicos,
para solicitar refacciones especiales rápidamente. \\
Criterio de aceptación 1 & Dado que el almacenista ingresa al módulo de
refacciones, cuando abre el directorio de proveedores, entonces
visualiza la lista con datos de contacto actualizados (teléfonos y
correos). \\
Criterio de aceptación 2 & Dado que se requiere dar de alta un nuevo
proveedor, cuando el almacenista completa el formulario de registro,
entonces el sistema lo almacena para futuras solicitudes de piezas. \\
Reglas de negocio & RN-06 (Reubicación por espera de refacciones). \\
Wireframe & Pendiente (sección 4.10) \\
Casos de prueba & Pendiente (sección 5.2) \\
\end{longtable}
}

\textbf{HU-18: Registro de Costos e Historial de Compras}

{\def\LTcaptype{none} % do not increment counter
\begin{longtable}[]{@{}
  >{\raggedright\arraybackslash}p{(\linewidth - 2\tabcolsep) * \real{0.2778}}
  >{\raggedright\arraybackslash}p{(\linewidth - 2\tabcolsep) * \real{0.7222}}@{}}
\toprule\noalign{}
\begin{minipage}[b]{\linewidth}\raggedright
\textbf{Campo}
\end{minipage} & \begin{minipage}[b]{\linewidth}\raggedright
\textbf{Contenido}
\end{minipage} \\
\midrule\noalign{}
\endhead
\bottomrule\noalign{}
\endlastfoot
ID e issue & HU-18 / \#{[}18{]} \\
Historia & Como encargado de almacén, quiero registrar los costos,
fechas de llegada y tickets de compra de las refacciones asociadas a una
orden, para mantener un control financiero estricto. \\
Criterio de aceptación 1 & Dado que se adquiere una pieza externa con un
proveedor, cuando el almacenista registra el ticket, costo unitario y
fecha estimada de entrega, entonces el sistema lo vincula a la orden de
servicio. \\
Criterio de aceptación 2 & Dado que se ingresa la compra, cuando se
consolida el registro, entonces el sistema actualiza el historial de
adquisiciones de la refacción para auditoría interna. \\
Reglas de negocio & Sin regla de negocio específica; ver RF-08.2 y
RF-08.3. \\
Wireframe & Pendiente (sección 4.10) \\
Casos de prueba & Pendiente (sección 5.2) \\
\end{longtable}
}

\textbf{HU-19: Verificación de Stock Interno}

{\def\LTcaptype{none} % do not increment counter
\begin{longtable}[]{@{}
  >{\raggedright\arraybackslash}p{(\linewidth - 2\tabcolsep) * \real{0.2778}}
  >{\raggedright\arraybackslash}p{(\linewidth - 2\tabcolsep) * \real{0.7222}}@{}}
\toprule\noalign{}
\begin{minipage}[b]{\linewidth}\raggedright
\textbf{Campo}
\end{minipage} & \begin{minipage}[b]{\linewidth}\raggedright
\textbf{Contenido}
\end{minipage} \\
\midrule\noalign{}
\endhead
\bottomrule\noalign{}
\endlastfoot
ID e issue & HU-19 / \#{[}19{]} \\
Historia & Como jefe de mecánicos, quiero verificar el stock de
refacciones comunes en el inventario interno, para saber si se dispone
de la pieza de forma inmediata. \\
Criterio de aceptación 1 & Dado que el jefe de mecánicos busca un insumo
estándar, cuando consulta el inventario en el sistema, entonces
visualiza las existencias disponibles en almacén. \\
Criterio de aceptación 2 & Dado que hay existencia en stock, cuando se
selecciona la pieza, entonces el sistema descuenta automáticamente la
cantidad del inventario interno. \\
Reglas de negocio & Sin regla de negocio específica; ver RF-08.1. \\
Wireframe & Pendiente (sección 4.10) \\
Casos de prueba & Pendiente (sección 5.2) \\
\end{longtable}
}

\textbf{HU-20: Notificación Automática de Fin de Trabajos}

{\def\LTcaptype{none} % do not increment counter
\begin{longtable}[]{@{}
  >{\raggedright\arraybackslash}p{(\linewidth - 2\tabcolsep) * \real{0.2778}}
  >{\raggedright\arraybackslash}p{(\linewidth - 2\tabcolsep) * \real{0.7222}}@{}}
\toprule\noalign{}
\begin{minipage}[b]{\linewidth}\raggedright
\textbf{Campo}
\end{minipage} & \begin{minipage}[b]{\linewidth}\raggedright
\textbf{Contenido}
\end{minipage} \\
\midrule\noalign{}
\endhead
\bottomrule\noalign{}
\endlastfoot
ID e issue & HU-20 / \#{[}20{]} \\
Historia & Como recepcionista, quiero recibir una notificación
automática cuando los trabajos del vehículo finalicen, para avisar al
cliente que su unidad está lista para entrega. \\
Criterio de aceptación 1 & Dado que el mecánico marca la orden como
terminada en bahía, cuando se actualiza el estatus, entonces el sistema
envía una alerta interna a la recepcionista. \\
Criterio de aceptación 2 & Dado que la recepción recibe la alerta,
cuando se verifica el estado, entonces se habilita el aviso automatizado
al cliente de que su vehículo está listo. \\
Reglas de negocio & Sin regla de negocio específica; ver CU-07. \\
Wireframe & Pendiente (sección 4.10) \\
Casos de prueba & Pendiente (sección 5.2) \\
\end{longtable}
}

\textbf{HU-21: Bloqueo de Entrega por Saldo Pendiente}

{\def\LTcaptype{none} % do not increment counter
\begin{longtable}[]{@{}
  >{\raggedright\arraybackslash}p{(\linewidth - 2\tabcolsep) * \real{0.2778}}
  >{\raggedright\arraybackslash}p{(\linewidth - 2\tabcolsep) * \real{0.7222}}@{}}
\toprule\noalign{}
\begin{minipage}[b]{\linewidth}\raggedright
\textbf{Campo}
\end{minipage} & \begin{minipage}[b]{\linewidth}\raggedright
\textbf{Contenido}
\end{minipage} \\
\midrule\noalign{}
\endhead
\bottomrule\noalign{}
\endlastfoot
ID e issue & HU-21 / \#{[}21{]} \\
Historia & Como gerente del taller, quiero que el sistema impida
registrar la entrega de un vehículo cuya orden tenga saldo pendiente,
para evitar pérdidas financieras al taller. \\
Criterio de aceptación 1 & Dado que un vehículo se encuentra en estado
"Lista para Entrega", cuando la recepcionista intenta registrar la
salida y el balance de pago es mayor a cero, entonces el sistema bloquea
la acción de liberación. \\
Criterio de aceptación 2 & Dado que se activa el bloqueo por saldo
pendiente, cuando el sistema detecta que cuenta con un convenio de
crédito autorizado, entonces permite omitir la restricción
exclusivamente bajo permisos del rol Gerente. \\
Reglas de negocio & RN-10 (Restricción de entrega y excepción a
flotillas). \\
Wireframe & Pendiente (sección 4.10) \\
Casos de prueba & Pendiente (sección 5.2) \\
\end{longtable}
}

\textbf{HU-22: Autorización Excepcional de Crédito a Flotillas}

{\def\LTcaptype{none} % do not increment counter
\begin{longtable}[]{@{}
  >{\raggedright\arraybackslash}p{(\linewidth - 2\tabcolsep) * \real{0.2778}}
  >{\raggedright\arraybackslash}p{(\linewidth - 2\tabcolsep) * \real{0.7222}}@{}}
\toprule\noalign{}
\begin{minipage}[b]{\linewidth}\raggedright
\textbf{Campo}
\end{minipage} & \begin{minipage}[b]{\linewidth}\raggedright
\textbf{Contenido}
\end{minipage} \\
\midrule\noalign{}
\endhead
\bottomrule\noalign{}
\endlastfoot
ID e issue & HU-22 / \#{[}22{]} \\
Historia & Como gerente del taller, quiero autorizar de forma
excepcional la entrega a crédito únicamente para vehículos de flotillas
empresariales registradas, para cumplir con los convenios
comerciales. \\
Criterio de aceptación 1 & Dado que un vehículo pertenece a una empresa
con cuenta de flotilla registrada, cuando se intenta liberar sin pago en
caja, entonces el sistema solicita la clave de autorización
gerencial. \\
Criterio de aceptación 2 & Dado que el gerente ingresa su contraseña de
aprobación, cuando valida el convenio comercial, entonces el sistema
permite liberar la unidad y factura el adeudo a crédito empresarial. \\
Reglas de negocio & RN-10 (Restricción de entrega y excepción a
flotillas). \\
Wireframe & Pendiente (sección 4.10) \\
Casos de prueba & Pendiente (sección 5.2) \\
\end{longtable}
}

\textbf{HU-23: Registro de Pago Total y Cierre de Orden}

{\def\LTcaptype{none} % do not increment counter
\begin{longtable}[]{@{}
  >{\raggedright\arraybackslash}p{(\linewidth - 2\tabcolsep) * \real{0.2778}}
  >{\raggedright\arraybackslash}p{(\linewidth - 2\tabcolsep) * \real{0.7222}}@{}}
\toprule\noalign{}
\begin{minipage}[b]{\linewidth}\raggedright
\textbf{Campo}
\end{minipage} & \begin{minipage}[b]{\linewidth}\raggedright
\textbf{Contenido}
\end{minipage} \\
\midrule\noalign{}
\endhead
\bottomrule\noalign{}
\endlastfoot
ID e issue & HU-23 / \#{[}23{]} \\
Historia & Como recepcionista, quiero registrar el pago total del
servicio y cerrar formalmente la orden, para liberar la bahía de trabajo
y archivar el historial del vehículo. \\
Criterio de aceptación 1 & Dado que el cliente cubre el monto total en
caja, cuando la recepcionista registra el método de pago y confirma la
transacción, entonces el sistema salda la cuenta a cero. \\
Criterio de aceptación 2 & Dado que el pago ha sido registrado con
éxito, cuando se ejecuta el cierre, entonces la orden cambia a estado
"Cerrada", liberando de inmediato la bahía física ocupada y generando el
comprobante digital. \\
Reglas de negocio & RN-10 (Restricción de entrega y excepción a
flotillas), RN-13 (Garantía de mano de obra), RN-01 (Expiración y
vigencia de enlaces). \\
Wireframe & Pendiente (sección 4.10) \\
Casos de prueba & Pendiente (sección 5.2) \\
\end{longtable}
}

\textbf{HU-24: Consulta de Vigencia de Garantía}

{\def\LTcaptype{none} % do not increment counter
\begin{longtable}[]{@{}
  >{\raggedright\arraybackslash}p{(\linewidth - 2\tabcolsep) * \real{0.2778}}
  >{\raggedright\arraybackslash}p{(\linewidth - 2\tabcolsep) * \real{0.7222}}@{}}
\toprule\noalign{}
\begin{minipage}[b]{\linewidth}\raggedright
\textbf{Campo}
\end{minipage} & \begin{minipage}[b]{\linewidth}\raggedright
\textbf{Contenido}
\end{minipage} \\
\midrule\noalign{}
\endhead
\bottomrule\noalign{}
\endlastfoot
ID e issue & HU-24 / \#{[}24{]} \\
Historia & Como cliente, quiero consultar la vigencia de la garantía de
30 días en mano de obra asociada a mi servicio anterior, para reclamar
ajustes en caso de incidencias. \\
Criterio de aceptación 1 & Dado que el cliente revisa el expediente de
su servicio pasado en el portal web, cuando consulta el apartado de
garantías, entonces visualiza la fecha exacta de expiración (30 días
naturales en mano de obra). \\
Criterio de aceptación 2 & Dado que la garantía se encuentra vigente,
cuando se solicita una revisión por el mismo concepto, entonces el
sistema valida el estatus sin generar cargos adicionales de mano de
obra. \\
Reglas de negocio & RN-13 (Garantía de mano de obra). \\
Wireframe & Pendiente (sección 4.10) \\
Casos de prueba & Pendiente (sección 5.2) \\
\end{longtable}
}

\textbf{HU-25: Operación Offline del Personal Interno}

{\def\LTcaptype{none} % do not increment counter
\begin{longtable}[]{@{}
  >{\raggedright\arraybackslash}p{(\linewidth - 2\tabcolsep) * \real{0.2778}}
  >{\raggedright\arraybackslash}p{(\linewidth - 2\tabcolsep) * \real{0.7222}}@{}}
\toprule\noalign{}
\begin{minipage}[b]{\linewidth}\raggedright
\textbf{Campo}
\end{minipage} & \begin{minipage}[b]{\linewidth}\raggedright
\textbf{Contenido}
\end{minipage} \\
\midrule\noalign{}
\endhead
\bottomrule\noalign{}
\endlastfoot
ID e issue & HU-25 / \#{[}25{]} \\
Historia & Como recepcionista o jefe de mecánicos, quiero registrar la
llegada de autos, tomar fotografías y hacer diagnósticos sin conexión a
internet usando la batería del equipo, para que la operación no se
detenga ante fallas eléctricas. \\
Criterio de aceptación 1 & Dado que ocurre un corte de energía o caída
de internet en el taller, cuando el personal opera mediante las baterías
de laptops o tabletas, entonces el sistema local permite continuar
capturando registros, checklists y fotografías. \\
Criterio de aceptación 2 & Dado que se trabaja de forma local (offline),
cuando se guardan los datos, entonces el almacenamiento temporal
mediante Service Workers e IndexedDB asegura que no exista pérdida de
información. \\
Reglas de negocio & Sin regla de negocio específica; ver RNF-06 y
RNF-08. \\
Wireframe & Pendiente (sección 4.10) \\
Casos de prueba & Pendiente (sección 5.2) \\
\end{longtable}
}

\textbf{HU-26: Sincronización Automática en la Nube}

{\def\LTcaptype{none} % do not increment counter
\begin{longtable}[]{@{}
  >{\raggedright\arraybackslash}p{(\linewidth - 2\tabcolsep) * \real{0.2778}}
  >{\raggedright\arraybackslash}p{(\linewidth - 2\tabcolsep) * \real{0.7222}}@{}}
\toprule\noalign{}
\begin{minipage}[b]{\linewidth}\raggedright
\textbf{Campo}
\end{minipage} & \begin{minipage}[b]{\linewidth}\raggedright
\textbf{Contenido}
\end{minipage} \\
\midrule\noalign{}
\endhead
\bottomrule\noalign{}
\endlastfoot
ID e issue & HU-26 / \#{[}26{]} \\
Historia & Como recepcionista o jefe de mecánicos, quiero que los
registros y fotografías que capturé sin conexión se sincronicen
automáticamente cuando vuelva la red, para no perder ni duplicar
información. \\
Criterio de aceptación 1 & Dado que se restablece la conexión a internet
en el taller, cuando el sistema detecta conectividad, entonces inicia de
forma automática la sincronización de los datos almacenados
localmente. \\
Criterio de aceptación 2 & Dado que finaliza el proceso de
sincronización, cuando se actualiza la base de datos central en la nube,
entonces la información multimedia y de auditoría se consolida sin
duplicidades. \\
Reglas de negocio & Sin regla de negocio específica; ver RNF-09. \\
Wireframe & Pendiente (sección 4.10) \\
Casos de prueba & Pendiente (sección 5.2) \\
\end{longtable}
}

\textbf{HU-27: Gestión de Roles y Permisos de Seguridad}

{\def\LTcaptype{none} % do not increment counter
\begin{longtable}[]{@{}
  >{\raggedright\arraybackslash}p{(\linewidth - 2\tabcolsep) * \real{0.2778}}
  >{\raggedright\arraybackslash}p{(\linewidth - 2\tabcolsep) * \real{0.7222}}@{}}
\toprule\noalign{}
\begin{minipage}[b]{\linewidth}\raggedright
\textbf{Campo}
\end{minipage} & \begin{minipage}[b]{\linewidth}\raggedright
\textbf{Contenido}
\end{minipage} \\
\midrule\noalign{}
\endhead
\bottomrule\noalign{}
\endlastfoot
ID e issue & HU-27 / \#{[}27{]} \\
Historia & Como gerente del taller, quiero dar de alta al personal y
asignarle su rol (Gerente, Recepcionista, Jefe de Mecánicos o Encargado
de Almacén), para que cada quien acceda solo a la información que le
corresponde. \\
Criterio de aceptación 1 & Dado que un usuario intenta acceder a la
plataforma, cuando introduce sus credenciales, entonces el sistema
autentica su identidad y restringe los menús según su rol asignado
(Gerente, Recepcionista, Jefe de Mecánicos o Encargado de Almacén). \\
Criterio de aceptación 2 & Dado que un usuario sin privilegios intenta
ingresar a una ruta protegida (como reportes gerenciales), entonces el
sistema deniega el acceso y muestra una pantalla de error por permisos
insuficientes. \\
Reglas de negocio & Sin regla de negocio específica; ver RNF-05 y 4.8
(control de acceso por rol). \\
Wireframe & Pendiente (sección 4.10) \\
Casos de prueba & Pendiente (sección 5.2) \\
\end{longtable}
}

\textbf{HU-28: Accesibilidad PWA Responsiva para Clientes}

{\def\LTcaptype{none} % do not increment counter
\begin{longtable}[]{@{}
  >{\raggedright\arraybackslash}p{(\linewidth - 2\tabcolsep) * \real{0.2778}}
  >{\raggedright\arraybackslash}p{(\linewidth - 2\tabcolsep) * \real{0.7222}}@{}}
\toprule\noalign{}
\begin{minipage}[b]{\linewidth}\raggedright
\textbf{Campo}
\end{minipage} & \begin{minipage}[b]{\linewidth}\raggedright
\textbf{Contenido}
\end{minipage} \\
\midrule\noalign{}
\endhead
\bottomrule\noalign{}
\endlastfoot
ID e issue & HU-28 / \#{[}28{]} \\
Historia & Como cliente, quiero visualizar la plataforma web desde mi
teléfono celular de forma responsiva y sin instalar aplicaciones
adicionales, para consultar mi estatus con facilidad. \\
Criterio de aceptación 1 & Dado que el cliente abre el enlace recibido
por SMS o correo en un navegador móvil (iOS/Android), cuando carga la
página, entonces el diseño se adapta responsivamente a la pantalla sin
requerir descargas de tiendas de aplicaciones. \\
Criterio de aceptación 2 & Dado que el cliente navega por el portal web,
cuando consulta el progreso de su vehículo o galería de fotos, entonces
los elementos táctiles responden de manera fluida y adaptada a
dispositivos móviles. \\
Reglas de negocio & RN-01 (Expiración y vigencia de enlaces).
Portabilidad: RNF-01. \\
Wireframe & Pendiente (sección 4.10) \\
Casos de prueba & Pendiente (sección 5.2) \\
\end{longtable}
}

\textbf{HU-29: Generación de Reportes Gerenciales}

{\def\LTcaptype{none} % do not increment counter
\begin{longtable}[]{@{}
  >{\raggedright\arraybackslash}p{(\linewidth - 2\tabcolsep) * \real{0.2778}}
  >{\raggedright\arraybackslash}p{(\linewidth - 2\tabcolsep) * \real{0.7222}}@{}}
\toprule\noalign{}
\begin{minipage}[b]{\linewidth}\raggedright
\textbf{Campo}
\end{minipage} & \begin{minipage}[b]{\linewidth}\raggedright
\textbf{Contenido}
\end{minipage} \\
\midrule\noalign{}
\endhead
\bottomrule\noalign{}
\endlastfoot
ID e issue & HU-29 / \#{[}29{]} \\
Historia & Como gerente, quiero generar reportes mensuales de ingresos,
servicios realizados y tiempos de ocupación de las 4 bahías, para
evaluar el rendimiento general del taller. \\
Criterio de aceptación 1 & Dado que el gerente ingresa al módulo de
reportes, cuando selecciona el intervalo de fechas mensual, entonces el
sistema genera métricas consolidadas de ingresos y servicios
completados. \\
Criterio de aceptación 2 & Dado que se consultan los tiempos operativos,
cuando el sistema procesa las órdenes, entonces muestra gráficos de
ocupación de las 4 bahías para medir la eficiencia del taller. \\
Reglas de negocio & Sin regla de negocio específica; ver RF-11. \\
Wireframe & Pendiente (sección 4.10) \\
Casos de prueba & Pendiente (sección 5.2) \\
\end{longtable}
}

\textbf{HU-30: Cumplimiento de Privacidad y Uso de Datos Ficticios}

{\def\LTcaptype{none} % do not increment counter
\begin{longtable}[]{@{}
  >{\raggedright\arraybackslash}p{(\linewidth - 2\tabcolsep) * \real{0.2778}}
  >{\raggedright\arraybackslash}p{(\linewidth - 2\tabcolsep) * \real{0.7222}}@{}}
\toprule\noalign{}
\begin{minipage}[b]{\linewidth}\raggedright
\textbf{Campo}
\end{minipage} & \begin{minipage}[b]{\linewidth}\raggedright
\textbf{Contenido}
\end{minipage} \\
\midrule\noalign{}
\endhead
\bottomrule\noalign{}
\endlastfoot
ID e issue & HU-30 / \#{[}30{]} \\
Historia & Como cliente del taller, quiero conocer el aviso de
privacidad antes de proporcionar mis datos, para saber para qué se
usarán conforme a la LFPDPPP. \\
Criterio de aceptación 1 & Dado que la recepcionista registra a un
cliente nuevo, cuando captura sus datos personales, entonces el sistema
muestra el aviso de privacidad simplificado (4.8) antes de guardar el
registro. \\
Criterio de aceptación 2 & Dado que el cliente abre su enlace seguro por
primera vez, cuando carga la página, entonces el sistema muestra el
enlace al aviso de privacidad integral y cómo ejercer sus derechos
ARCO. \\
Reglas de negocio & Sin regla de negocio específica; ver 4.8 y
RNF-05. \\
Wireframe & Pendiente (sección 4.10) \\
Casos de prueba & Pendiente (sección 5.2) \\
\end{longtable}
}

{[}Figura 4.1. Story map con el corte del release v0.1.0.{]}

\subsection{4.6 Requisitos no
funcionales}\label{requisitos-no-funcionales}

{\def\LTcaptype{none} % do not increment counter
\begin{longtable}[]{@{}
  >{\raggedright\arraybackslash}p{(\linewidth - 8\tabcolsep) * \real{0.0962}}
  >{\raggedright\arraybackslash}p{(\linewidth - 8\tabcolsep) * \real{0.2136}}
  >{\raggedright\arraybackslash}p{(\linewidth - 8\tabcolsep) * \real{0.2735}}
  >{\raggedright\arraybackslash}p{(\linewidth - 8\tabcolsep) * \real{0.2137}}
  >{\raggedright\arraybackslash}p{(\linewidth - 8\tabcolsep) * \real{0.2031}}@{}}
\toprule\noalign{}
\begin{minipage}[b]{\linewidth}\raggedright
\textbf{ID}
\end{minipage} & \begin{minipage}[b]{\linewidth}\raggedright
\textbf{Característica ISO/IEC 25010:2023}
\end{minipage} & \begin{minipage}[b]{\linewidth}\raggedright
\textbf{Requisito}
\end{minipage} & \begin{minipage}[b]{\linewidth}\raggedright
\textbf{Métrica y umbral}
\end{minipage} & \begin{minipage}[b]{\linewidth}\raggedright
\textbf{Cómo se verificará}
\end{minipage} \\
\midrule\noalign{}
\endhead
\bottomrule\noalign{}
\endlastfoot
RNF-01 & Flexibilidad (instalabilidad y adaptabilidad) & El portal del
cliente deberá ser una PWA responsiva que funcione en navegadores
móviles iOS y Android sin instalar software ni conceder permisos
especiales. & Funciona sin errores bloqueantes en las 2 últimas
versiones de Safari (iOS) y Chrome (Android); 0 permisos solicitados al
cliente; sin desplazamiento horizontal en anchos de 360 a 430 px. &
Pruebas manuales en 2 móviles reales (iOS y Android) con una lista de
verificación de las pantallas del cliente y revisión con las
herramientas de desarrollo del navegador. \\
RNF-02 & Capacidad de interacción & La interfaz de la tableta del jefe
de Mecánicos deberá tener botones grandes, alto contraste y controles
táctiles adaptados al uso con manos sucias y en ambiente ruidoso. &
Botones de al menos 48 × 48 px con separación mínima de 8 px; contraste
de texto de al menos 7:1; al menos 95 \% de aciertos en 20 toques con
guantes de trabajo. & Auditoría con Lighthouse o axe DevTools
(gratuitos), medición del contraste con un verificador y prueba en la
tableta real con guantes. \\
RNF-03 & Capacidad de interacción & La interfaz de la tableta deberá
ofrecer flujos de pocos pasos para las operaciones frecuentes del
taller. & Asignar bahía, actualizar el estado de una orden y sellar un
trabajo en máximo 4 toques y 30 s cada una. & Prueba de usabilidad
cronometrada con al menos 3 usuarios (o compañeros que simulen el rol),
repitiendo cada operación 3 veces. \\
RNF-04 & Seguridad (no repudio e integridad) & Toda autorización de
presupuesto o de sobrecosto mayor al 10 \% deberá registrar marca de
tiempo y metadatos inalterables, incluida la dirección IP, para evitar
disputas con el cliente. & 100 \% de las autorizaciones guardan fecha y
hora UTC, IP, identificador de orden y huella SHA-256 del documento; 100
\% de los intentos de alterar un registro se detectan al verificar la
huella. & Prueba con 20 autorizaciones, revisión de los registros
guardados e intento de modificar directamente la base de datos para
confirmar que la verificación lo detecta. \\
RNF-05 & Seguridad (confidencialidad y cumplimiento legal) & El manejo
de datos personales deberá apegarse a la Ley Federal de Protección de
Datos Personales en Posesión de los Particulares y usar solo datos
ficticios en las pruebas. & 100 \% de las conexiones con HTTPS (TLS 1.2
o superior); datos personales sensibles cifrados en reposo; aviso de
privacidad visible en el 100 \% de los formularios que capturan datos; 0
registros reales en el ambiente de pruebas. & Análisis con testssl.sh o
SSL Labs, lista de verificación de la ley revisada por el equipo y
revisión de los datos de prueba antes de cada entrega. \\
RNF-06 & Eficiencia de desempeño & El registro de recepción se deberá
guardar de forma local con rapidez, aun sin conexión. & Guardar un
registro con 5 fotografías en máximo 2 s (percentil 95) en la tableta o
laptop del taller. & 20 repeticiones cronometradas con la herramienta de
rendimiento del navegador, calculando el percentil 95. \\
RNF-07 & Mantenibilidad (modificabilidad) & El desarrollo deberá usar
únicamente herramientas de código abierto, frameworks de uso libre (como
FastAPI y Python) y planes educativos. & 100 \% de las dependencias con
licencia de código abierto o plan educativo vigente; costo de
licenciamiento de \$0 en esta fase. & Inventario de dependencias con
herramientas de revisión de licencias (por ejemplo, pip-licenses) y
registro de los planes educativos usados. \\
RNF-08 & Fiabilidad (tolerancia a fallos) & El personal interno deberá
poder operar sin internet o sin energía eléctrica, con equipos en
batería, bajo una arquitectura local con sincronización diferida. & Las
7 operaciones definidas (registro de entrada, checklist, fotografías,
diagnósticos, asignación de bahías, consulta de proveedores y
autorizaciones gerenciales) funcionan al 100 \% durante al menos 4 h sin
conexión; 0 datos perdidos. & Casos de prueba de las 7 operaciones con
la red desconectada y un equipo con batería, comparando los registros al
terminar. \\
RNF-09 & Fiabilidad (recuperabilidad) & Al restablecerse la conexión, el
sistema deberá sincronizar automáticamente datos, archivos multimedia y
registros de auditoría con la base de datos central, sin duplicados. &
100 \% de los registros sincronizados, 0 duplicados y 0 pérdidas; 200
registros con 50 fotografías en máximo 5 min. & Crear 200 registros sin
conexión, reconectar y comparar conteos y huellas con la base central;
repetir 3 veces e interrumpir la conexión a mitad de una
sincronización. \\
RNF-10 & Fiabilidad (tolerancia a fallos del cliente) & Ante una falla
de conectividad durante la consulta de estatus o la autorización 2FA, la
interfaz del cliente deberá informar el estado y reintentar
automáticamente. & Mensaje de estado en máximo 3 s; reintento automático
cada 30 s o menos hasta completar; 0 autorizaciones duplicadas o
perdidas. & Cortar el servidor durante la consulta y durante el envío
del código, cronometrar el mensaje y revisar la base de datos tras
restablecer el servicio. \\
\end{longtable}
}

\subsection{4.7 Reglas de negocio y requisitos de
dominio}\label{reglas-de-negocio-y-requisitos-de-dominio}

{\def\LTcaptype{none} % do not increment counter
\begin{longtable}[]{@{}
  >{\raggedright\arraybackslash}p{(\linewidth - 6\tabcolsep) * \real{0.0963}}
  >{\raggedright\arraybackslash}p{(\linewidth - 6\tabcolsep) * \real{0.4871}}
  >{\raggedright\arraybackslash}p{(\linewidth - 6\tabcolsep) * \real{0.2351}}
  >{\raggedright\arraybackslash}p{(\linewidth - 6\tabcolsep) * \real{0.1815}}@{}}
\toprule\noalign{}
\begin{minipage}[b]{\linewidth}\raggedright
\textbf{ID}
\end{minipage} & \begin{minipage}[b]{\linewidth}\raggedright
\textbf{Regla de negocio}
\end{minipage} & \begin{minipage}[b]{\linewidth}\raggedright
\textbf{Fuente (anteproyecto / entrevista con fecha)}
\end{minipage} & \begin{minipage}[b]{\linewidth}\raggedright
\textbf{Requisitos que afecta}
\end{minipage} \\
\midrule\noalign{}
\endhead
\bottomrule\noalign{}
\endlastfoot
RN-01 & Expiración y vigencia de enlaces: El enlace web de seguimiento
permanecerá activo durante la estancia del vehículo en el taller.
Expirará automáticamente al entregar y saldar la unidad. Si el cliente
requiere ver su historial posteriormente, deberá solicitar la generación
de un nuevo enlace. La vigencia del enlace es independiente de la
vigencia de la cotización (RN-11). & Entrevista con el cliente
(19/10/2026) & RF-02.1, RF-02.2 \\
RN-02 & Doble verificación en autorizaciones: Toda aprobación remota o
presencial de cotización exigirá una doble verificación (código 2FA de 6
dígitos, de un solo uso y vigente 5 minutos, enviado al correo
electrónico registrado del cliente) para evitar reclamos por trabajos no
autorizados. & Entrevista con el cliente (19/10/2026) & RF-06, RNF-04 \\
RN-03 & Selección granular de rubros: El cliente puede aprobar o
rechazar de manera independiente los conceptos opcionales desglosados en
el presupuesto (mano de obra o refacciones), ejecutando únicamente lo
autorizado explícitamente. & Entrevista con el cliente (19/10/2026) &
RF-06 \\
RN-04 & Re-cotización obligatoria por sobrecostos (\textgreater10\%):
Cualquier trabajo o repuesto imprevisto durante la reparación que
incremente el costo inicial más de un 10\% detendrá el avance hasta
contar con una nueva autorización validada por el cliente. &
Anteproyecto y propuesta técnica & RF-07, RNF-04 \\
RN-05 & Capacidad física del taller\textbf{:} La atención simultánea
está limitada estrictamente a 4 bahías físicas (1 vehículo por bahía) y
un área de estacionamiento de recepción/espera con capacidad máxima para
5 vehículos. & Anteproyecto y entrevista con el cliente & RF-05,
RF-09 \\
RN-06 & Reubicación por espera de refacciones: Si un vehículo requiere
refacciones externas con demora de entrega de varios días, se liberará
la bahía de trabajo reubicando la unidad en el estacionamiento hasta el
arribo de las piezas. & Entrevista con el cliente (19/10/2026) & RF-05,
RF-08.2 \\
RN-07 & Clasificación de vehículos autorizados: El sistema operará
exclusivamente con vehículos automotrices ligeros (automóviles
particulares, motocicletas y camionetas hasta de 3 toneladas). Se
excluyen bicicletas, tráilers y maquinaria pesada. & Entrevista con el
cliente (19/10/2026) & RF-03.1 \\
RN-08 & Confirmación manual de citas web: Las solicitudes de cita
agendadas desde el portal público ingresarán en estado "Pendiente" y no
reservarán espacio en la agenda hasta que la recepcionista las confirme
manualmente. & Entrevista con el cliente (19/10/2026) & RF-09 \\
RN-09 & Asignación y colaboración de mecánicos\textbf{:} Cada orden
mantendrá un mecánico principal responsable. Otros mecánicos podrán
colaborar en el trabajo técnico sin alterar la titularidad de la orden,
salvo reasignación formal. & Entrevista con el cliente (19/10/2026) &
RF-05 \\
RN-10 & Restricción de entrega y excepción a flotillas: Queda prohibida
la entrega física del vehículo con saldo pendiente. Como única
excepción, el Gerente podrá autorizar la entrega a crédito para
vehículos pertenecientes a flotillas empresariales registradas. &
Anteproyecto y entrevista con el cliente & RF-10 \\
RN-11 & Vigencia de cotizaciones: Toda cotización tiene vigencia de 5
días naturales desde su emisión. Una cotización vencida no puede
autorizarse; el taller debe emitir una nueva versión. & Anteproyecto &
RF-06 \\
RN-12 & Evidencia fotográfica de recepción: Al recibir un vehículo se
exigen el checklist completo y al menos 4 fotografías (frente, trasera,
costado izquierdo y costado derecho); sin ellas la orden no puede pasar
a diagnóstico. & Anteproyecto; el mínimo de 4 fotografías debe validarse
con el cliente & RF-04 \\
RN-13 & Garantía de mano de obra: Los trabajos tienen garantía de 30
días naturales sobre la mano de obra, contados desde la entrega del
vehículo. & RF-03.2; confirmar la fuente con el cliente & RF-03.2 \\
\end{longtable}
}

\subsection{4.8 Requisitos de protección de datos
personales}\label{requisitos-de-protecciuxf3n-de-datos-personales}

En cumplimiento con la Ley Federal de Protección de Datos Personales en
Posesión de los Particulares (LFPDPPP - DOF):

\begin{enumerate}
\def\labelenumi{\arabic{enumi}.}
\item
  Datos Personales Recabados
\end{enumerate}

\begin{quote}
El sistema capturará y procesará los siguientes datos de clientes
particulares y representantes de flotillas:
\end{quote}

\begin{itemize}
\item
  Datos de identificación y contacto: Nombre completo, número telefónico
  celular (WhatsApp/SMS) y correo electrónico.
\item
  Datos del vehículo y de la propiedad: Placas, marca, modelo, año,
  número de serie (VIN - opcional), kilometraje, nivel de combustible,
  fotografías del estado físico de la unidad, historial de servicios e
  importes de facturación.
\end{itemize}

\begin{enumerate}
\def\labelenumi{\arabic{enumi}.}
\setcounter{enumi}{1}
\item
  Finalidad del Tratamiento
\end{enumerate}

\begin{itemize}
\item
  Finalidades primarias: Identificación del cliente, vinculación del
  vehículo a la orden de servicio, generación y envío de enlaces
  temporales de seguimiento, envío de códigos 2FA para validación de
  presupuestos, cobro, emisión de comprobantes de pago y aplicación de
  garantías de mano de obra (30 días).
\item
  Finalidades secundarias: Recordatorios de mantenimiento preventivo y
  consulta de historial a solicitud del cliente.
\end{itemize}

\begin{enumerate}
\def\labelenumi{\arabic{enumi}.}
\setcounter{enumi}{2}
\item
  Medidas de Seguridad y Control de Accesos
\end{enumerate}

\begin{itemize}
\item
  Control de Acceso Basado en Roles (RBAC): La información personal del
  cliente solo estará accesible para los roles de Recepcionista y
  Gerente. Los Mecánicos y Encargados de Almacén únicamente visualizarán
  datos técnicos del vehículo y requerimientos de refacciones.
\item
  Acceso mediante Tokens Temporales: Los clientes accederán a sus
  cotizaciones y estatus vía tokens dinámicos sin necesidad de almacenar
  contraseñas de usuarios externos en la base de datos. Protección de
  Datos en Pruebas: Se garantiza el uso exclusivo de datos sintéticos o
  ficticios (mock data) durante todas las fases de desarrollo, pruebas y
  demostraciones académicas, sin almacenar registros reales en los
  entornos de staging o prueba.
\item
  Seguridad en Tránsito: Las comunicaciones entre navegadores móviles,
  tabletas y servidores se realizarán mediante protocolos cifrados HTTPS
  (TLS 1.2 o superior).
\end{itemize}

\begin{enumerate}
\def\labelenumi{\arabic{enumi}.}
\setcounter{enumi}{3}
\item
  Borrador del Aviso de Privacidad Simplificado
\end{enumerate}

\begin{quote}
AVISO DE PRIVACIDAD SIMPLIFICADO

Servicio Automotriz Hermanos Ruiz, S.A. de C.V., con domicilio en las
instalaciones operativas del taller, es responsable del tratamiento de
los datos personales (contacto, identificación y datos de su vehículo)
proporcionados en la recepción.

Sus datos serán utilizados exclusivamente para la apertura de órdenes de
servicio, seguimiento en vivo, envío de cotizaciones para su
autorización remota o presencial mediante doble verificación (2FA),
cobro, entrega y aplicación de garantías. Sus datos no serán
transferidos ni compartidos con terceros sin su consentimiento
explícito. Para consultar el Aviso de Privacidad Integral o ejercer sus
derechos ARCO (Acceso, Rectificación, Cancelación y Oposición), puede
acudir a la recepción de nuestro taller o solicitarlo a través de
nuestro portal web de atención a clientes.
\end{quote}

\subsection{4.9 Casos de uso}\label{casos-de-uso}

CU-01 --- Confirmación de cita web

{\def\LTcaptype{none} % do not increment counter
\begin{longtable}[]{@{}
  >{\raggedright\arraybackslash}p{(\linewidth - 2\tabcolsep) * \real{0.2351}}
  >{\raggedright\arraybackslash}p{(\linewidth - 2\tabcolsep) * \real{0.7649}}@{}}
\toprule\noalign{}
\begin{minipage}[b]{\linewidth}\raggedright
Campo
\end{minipage} & \begin{minipage}[b]{\linewidth}\raggedright
Contenido
\end{minipage} \\
\midrule\noalign{}
\endhead
\bottomrule\noalign{}
\endlastfoot
ID y nombre & CU-01 Confirmación de cita web \\
Actor principal & Recepcionista \\
Actores secundarios & Cliente; servicio de correo electrónico. \\
Descripción breve & Permite revisar solicitudes de citas realizadas
desde el sitio web, verificar la disponibilidad del taller y confirmar
fecha, hora y bahía preliminar. El sistema actualiza el estado y
notifica al cliente. \\
Precondiciones & 1. El cliente solicitó previamente una cita mediante el
sitio web.

2. La solicitud está registrada y pendiente de revisión.

3. La recepcionista dispone de credenciales válidas. \\
Flujo principal & 1. La recepcionista inicia sesión desde la computadora
de mostrador.

2. El sistema valida las credenciales y muestra el menú principal.

3. La recepcionista abre el módulo de Gestión de Citas.

4. El sistema muestra las solicitudes pendientes.

5. La recepcionista selecciona una solicitud y revisa los datos del
cliente.

6. El sistema muestra la agenda y la disponibilidad de horarios y
bahías.

7. La recepcionista verifica la disponibilidad y selecciona fecha, hora
y bahía preliminar.

8. La recepcionista confirma la programación.

9. El sistema actualiza la cita a «Confirmada» y guarda la programación.

10. El sistema envía una notificación al cliente por correo electrónico.

11. El sistema informa a la recepcionista que la cita fue confirmada. \\
\textbf{Flujos alternos} & A1 (paso 7): Horario o fecha sin
disponibilidad. El sistema muestra que no hay espacio y permite
seleccionar una alternativa; recepción propone la nueva fecha u hora al
cliente y continúa con la confirmación.

A2 (paso 4): No hay solicitudes pendientes. El sistema informa que la
bandeja está vacía y la recepcionista puede salir del módulo. \\
\textbf{Flujos de excepción} & E1 (paso 1): Credenciales incorrectas o
cuenta bloqueada. El sistema rechaza el acceso; después de tres intentos
fallidos consecutivos bloquea temporalmente la cuenta durante 15
minutos.

E2 (paso 9): Error al guardar. El sistema informa que no fue posible
actualizar la cita y conserva el estado anterior.

E3 (paso 10): Fallo de correo. La cita permanece confirmada y el envío
queda para seguimiento o reintento. \\
\textbf{Postcondiciones} & La cita queda registrada con estado
«Confirmada» y con fecha y hora asignadas. El sistema conserva el
resultado del envío de la notificación. \\
\textbf{Reglas de negocio} & Reglas de 4.7: RN-08 (Confirmación manual
de citas web), RN-05 (Capacidad física del taller).

Validaciones del caso de uso:

• Solo se confirma una solicitud registrada.

• La programación debe respetar la disponibilidad de horarios y bahías.

• La cita confirmada debe contar con fecha y hora.

• La notificación se envía al correo registrado. \\
\textbf{Requisitos relacionados} & RF-09 · HU-04 \\
\end{longtable}
}

CU-02 --- Recepción del vehículo y apertura de orden

{\def\LTcaptype{none} % do not increment counter
\begin{longtable}[]{@{}
  >{\raggedright\arraybackslash}p{(\linewidth - 2\tabcolsep) * \real{0.2351}}
  >{\raggedright\arraybackslash}p{(\linewidth - 2\tabcolsep) * \real{0.7649}}@{}}
\toprule\noalign{}
\begin{minipage}[b]{\linewidth}\raggedright
\textbf{Campo}
\end{minipage} & \begin{minipage}[b]{\linewidth}\raggedright
\textbf{Contenido}
\end{minipage} \\
\midrule\noalign{}
\endhead
\bottomrule\noalign{}
\endlastfoot
\textbf{ID y nombre} & CU-02 Recepción del vehículo y apertura de
orden \\
\textbf{Actor principal} & Recepcionista \\
\textbf{Actores secundarios} & Cliente. \\
\textbf{Descripción breve} & Registra la llegada del vehículo, verifica
o captura los datos del cliente y del automóvil, documenta sus
condiciones iniciales y abre una orden de servicio para dar seguimiento
al trabajo. \\
\textbf{Precondiciones} & 1. El cliente se presenta físicamente en el
taller con su vehículo.

2. La recepcionista tiene acceso autorizado al sistema.

3. El sistema permite registrar clientes, vehículos y órdenes. \\
\textbf{Flujo principal} & 1. La recepcionista inicia sesión o continúa
con una sesión activa.

2. Accede al módulo de Registro de Cliente y Vehículo.

3. Busca al cliente por los datos disponibles.

4. El sistema muestra al cliente y sus vehículos registrados, si
existen.

5. La recepcionista verifica o registra los datos del cliente y del
vehículo: placas, marca, modelo y VIN.

6. Accede al módulo de Check-in y Registro de Entrada.

7. El sistema solicita los datos iniciales del vehículo.

8. La recepcionista registra kilometraje y nivel de combustible.

9. Completa la lista de inspección inicial, incluyendo herramientas y
llanta de refacción, entre otros elementos definidos por el taller.

10. Toma o carga fotografías de la carrocería para documentar daños
previos.

11. Revisa los datos y selecciona «Abrir Orden de Servicio».

12. El sistema valida los datos obligatorios y registra la orden con
estado «Abierta».

13. El sistema genera un enlace seguro asociado a la orden para el
cliente.

14. El sistema confirma la apertura de la orden. \\
\textbf{Flujos alternos} & A1 (paso 3): Cliente no registrado. La
recepcionista captura sus datos y continúa con el registro del vehículo.

A2 (paso 5): Vehículo no registrado. La recepcionista captura sus datos
y lo vincula con el cliente.

A3 (paso 10): No se observan daños visibles. La recepcionista registra
el resultado de la inspección y continúa. \\
\textbf{Flujos de excepción} & E1 (paso 5): Datos obligatorios
incompletos. El sistema identifica los campos faltantes y no permite
guardar hasta completarlos.

E2 (paso 5): Vehículo con orden previa abierta. Si las placas
corresponden a una orden activa, el sistema muestra el folio y dirige a
la orden existente para actualizarla o cerrarla antes de abrir otra.

E3 (paso 10): Fotografía inválida o mayor a 10 MB. El sistema muestra un
error de formato o tamaño y solicita seleccionar una imagen compatible o
comprimida.

E4 (paso 12): Error al abrir la orden. El sistema informa que la
operación no se completó y no debe mostrar una orden como abierta si no
fue guardada. \\
\textbf{Postcondiciones} & Quedan registrados el check-in y la orden de
servicio en estado «Abierta»; el vehículo se asocia al cliente y se
genera el enlace seguro correspondiente. \\
\textbf{Reglas de negocio} & Reglas de 4.7: RN-12 (Evidencia fotográfica
de recepción), RN-07 (Clasificación de vehículos autorizados), RN-01
(Expiración y vigencia de enlaces).

Validaciones del caso de uso:

• Cada vehículo debe estar asociado con un cliente.

• Se registra kilometraje y nivel de combustible al ingreso.

• La inspección inicial documenta el estado del vehículo.

• La orden requiere los datos obligatorios completos.

• La orden tiene folio único y estado inicial «Abierta».

• El enlace seguro debe corresponder a la orden correcta. \\
\textbf{Requisitos relacionados} & RF-02.1, RF-03.1, RF-04 · HU-01,
HU-05, HU-06 \\
\end{longtable}
}

\textbf{CU-03 --- Firma de contrato presencial (2FA)}

{\def\LTcaptype{none} % do not increment counter
\begin{longtable}[]{@{}
  >{\raggedright\arraybackslash}p{(\linewidth - 2\tabcolsep) * \real{0.2387}}
  >{\raggedright\arraybackslash}p{(\linewidth - 2\tabcolsep) * \real{0.7158}}@{}}
\toprule\noalign{}
\begin{minipage}[b]{\linewidth}\raggedright
\textbf{Campo}
\end{minipage} & \begin{minipage}[b]{\linewidth}\raggedright
\textbf{Contenido}
\end{minipage} \\
\midrule\noalign{}
\endhead
\bottomrule\noalign{}
\endlastfoot
\textbf{ID y nombre} & CU-03 Firma de contrato presencial (2FA) \\
\textbf{Actor principal} & Recepcionista \\
\textbf{Actores secundarios} & Cliente; servicio de correo
electrónico. \\
\textbf{Descripción breve} & Formaliza en el mostrador la autorización
de una cotización técnica mediante un código de un solo uso enviado al
correo registrado del cliente. \\
\textbf{Precondiciones} & 1. El jefe de mecánicos ya generó el
diagnóstico y la cotización técnica.

2. La cotización está asociada a una orden y disponible para revisión.

3. El cliente decide autorizar de forma presencial.

4. El cliente tiene un correo registrado y accesible.

5. La recepcionista tiene acceso al sistema. \\
\textbf{Flujo principal} & 1. El clnte comunica que desea autorizar
presencialmente la cotización.

2. La recepcionista abre el módulo de Firma de Contrato Presencial.

3. El sistema muestra las órdenes y cotizaciones pendientes de
autorización.

4. La recepcionista selecciona la orden correspondiente.

5. El sistema muestra el desglose de la cotización.

6. El cliente revisa el alcance y el importe cotizado.

7. La recepcionista selecciona «Enviar Código 2FA».

8. El sistema genera un código de verificación y lo envía al correo
registrado.

9. El cliente recibe el código y se lo proporciona a la recepcionista.

10. La recepcionista introduce los seis dígitos.

11. La recepcionista selecciona «Validar y Firmar».

12. El sistema comprueba que el código sea correcto, vigente y no
utilizado.

13. El sistema registra la autorización y cambia la orden a
«Autorizada».

14. El sistema muestra la confirmación de la operación. \\
\textbf{Flujos alternos} & A1 (paso 6): El cliente solicita
aclaraciones. La recepcionista aclara la información disponible o
consulta al jefe de mecánicos; después el cliente puede continuar o
rechazar la cotización.

A2 (paso 9): El cliente no recibe el código. Se verifica el correo y el
estado del envío; si está habilitado, se reenvía el código.

A3 (paso 10): Si el sistema lo permite, el cliente introduce
personalmente el código. \\
\textbf{Flujos de excepción} & E1 (paso 8): Fallo de envío. El sistema
informa que no pudo enviar el código y no permite autorizar hasta
completar la verificación.

E2 (paso 12): Código incorrecto o expirado tras cinco minutos. El
sistema lo rechaza y permite introducir un código vigente o solicitar un
reenvío.

E3 (paso 12): Límite de intentos excedido. El sistema bloquea
temporalmente la validación o exige iniciar nuevamente el proceso.

E4 (paso 13): Error al registrar la autorización. La orden conserva el
estado anterior hasta que se registre correctamente.

E5: El cliente rechaza la cotización. La recepcionista registra el
rechazo; la orden queda «Cancelada/Rechazada» y se coordina el cobro de
revisión y la entrega del vehículo sin reparar, si aplica. \\
\textbf{Postcondiciones} & Si la validación tiene éxito, la autorización
queda registrada y la orden cambia a «Autorizada». Si el cliente rechaza
la cotización, se registra el rechazo y no continúa la reparación
cotizada. \\
\textbf{Reglas de negocio} & Reglas de 4.7: RN-02 (Doble verificación en
autorizaciones), RN-03 (Selección granular de rubros), RN-11 (Vigencia
de cotizaciones).

Validaciones del caso de uso:

• La cotización debe haberse generado y estar asociada a la orden.

• El cliente debe revisar la cotización antes de autorizarla.

• El código es de un solo uso y tiene vigencia definida de cinco
minutos, según la especificación recibida.

• Se rechazan códigos incorrectos, vencidos o ya utilizados.

• La orden solo cambia a «Autorizada» después de validar y registrar la
aceptación. \\
\textbf{Requisitos relacionados} & RF-06 · HU-12, HU-14 \\
\end{longtable}
}

CU-04 --- Cobro y entrega del vehículo

{\def\LTcaptype{none} % do not increment counter
\begin{longtable}[]{@{}
  >{\raggedright\arraybackslash}p{(\linewidth - 2\tabcolsep) * \real{0.2387}}
  >{\raggedright\arraybackslash}p{(\linewidth - 2\tabcolsep) * \real{0.7613}}@{}}
\toprule\noalign{}
\begin{minipage}[b]{\linewidth}\raggedright
\textbf{Campo}
\end{minipage} & \begin{minipage}[b]{\linewidth}\raggedright
\textbf{Contenido}
\end{minipage} \\
\midrule\noalign{}
\endhead
\bottomrule\noalign{}
\endlastfoot
\textbf{ID y nombre} & CU-04 Cobro y entrega del vehículo \\
\textbf{Actor principal} & Recepcionista \\
\textbf{Actores secundarios} & Cliente; gerente (cuando se requiera
autorización de crédito); sistema de registro de pagos. \\
\textbf{Descripción breve} & Registra el pago de la orden, emite el
comprobante, entrega las llaves al cliente y cierra la orden de servicio
cuando se cumplen las condiciones de salida. \\
\textbf{Precondiciones} & 1. El jefe de mecánicos marcó el trabajo como
«Terminado» y el vehículo está listo para entrega.

2. Existe una orden de servicio registrada.

3. La recepcionista tiene permiso para realizar el cobro y registrar la
entrega. \\
\textbf{Flujo principal} & 1. El cliente acude al taller a recoger su
vehículo.

2. La recepcionista abre el módulo de Cobro y Entrega.

3. El sistema muestra el saldo pendiente de la orden.

4. La recepcionista verifica el importe, registra el pago y emite el
comprobante.

5. La recepcionista entrega las llaves al cliente.

6. La recepcionista selecciona «Registrar Entrega».

7. El sistema guarda la orden en el Historial de Servicio del vehículo y
cambia el estado a «Cerrada».

8. El sistema confirma el cierre de la orden. \\
\textbf{Flujos alternos} & A1: Entrega a crédito para flotilla
empresarial. Si el vehículo pertenece a una flotilla y queda saldo
pendiente, la recepcionista solicita autorización al gerente. Si el
gerente autoriza, se continúa con el registro de entrega sin pago
inmediato de contado (véase CU-13).

A2: El cliente paga el saldo pendiente y se reanuda el proceso de cobro
y entrega. \\
\textbf{Flujos de excepción} & E1: Saldo pendiente sin autorización de
crédito. El sistema bloquea la entrega y muestra que se requiere
liquidación completa o autorización expresa de Gerencia.

E2: Sesión expirada por más de 30 minutos de inactividad. El sistema
redirige al inicio de sesión y conserva el borrador del formulario de
cobro; tras autenticarse, la recepcionista retoma el proceso. \\
\textbf{Postcondiciones} & La orden queda «Cerrada», el pago queda
registrado, se emite un comprobante y la orden se conserva en el
historial de servicio del vehículo. \\
\textbf{Reglas de negocio} & Reglas de 4.7: RN-10 (Restricción de
entrega y excepción a flotillas), RN-13 (Garantía de mano de obra),
RN-01 (Expiración y vigencia de enlaces).

Validaciones del caso de uso:

• No se entrega un vehículo con saldo pendiente salvo autorización
expresa de Gerencia para crédito de flotilla.

• El pago debe registrarse en el sistema y producir un comprobante.

• La orden se cierra después de registrar la entrega.

• El cierre debe quedar en el historial de servicio del vehículo. \\
\textbf{Requisitos relacionados} & RF-02.2, RF-03.2, RF-10 · HU-21,
HU-23 \\
\end{longtable}
}

CU-05 --- Diagnóstico, cotización técnica y asignación de trabajo

{\def\LTcaptype{none} % do not increment counter
\begin{longtable}[]{@{}
  >{\raggedright\arraybackslash}p{(\linewidth - 2\tabcolsep) * \real{0.2387}}
  >{\raggedright\arraybackslash}p{(\linewidth - 2\tabcolsep) * \real{0.7613}}@{}}
\toprule\noalign{}
\begin{minipage}[b]{\linewidth}\raggedright
\textbf{Campo}
\end{minipage} & \begin{minipage}[b]{\linewidth}\raggedright
\textbf{Contenido}
\end{minipage} \\
\midrule\noalign{}
\endhead
\bottomrule\noalign{}
\endlastfoot
\textbf{ID y nombre} & CU-05 Diagnóstico, cotización técnica y
asignación de trabajo \\
\textbf{Actor principal} & Jefe de Mecánicos (líder técnico) \\
\textbf{Actores secundarios} & Recepcionista; mecánico asignado; sistema
de agenda de bahías. \\
\textbf{Descripción breve} & Permite inspeccionar técnicamente un
vehículo con orden abierta, elaborar el diagnóstico y la cotización de
mano de obra y refacciones, y asignar una bahía y un mecánico. \\
\textbf{Precondiciones} & 1. Existe una orden de servicio en estado
«Abierta».

2. El jefe de mecánicos tiene acceso autorizado desde la tableta del
taller.

3. Hay información suficiente para identificar la orden y el
vehículo. \\
\textbf{Flujo principal} & 1. El jefe de mecánicos inicia sesión en la
tableta del taller.

2. Revisa las órdenes abiertas y localiza el vehículo en el área de
recepción.

3. Inspecciona físicamente el automóvil.

4. Abre el módulo de Diagnóstico y Cotización Técnica.

5. Registra el diagnóstico y desglosa la mano de obra estimada y las
refacciones necesarias.

6. Consulta el Tablero de Bahías y Mecánicos.

7. Asigna la orden a una bahía disponible y a un mecánico de su equipo.

8. Selecciona «Guardar Diagnóstico».

9. El sistema guarda el diagnóstico y cambia la orden a «Diagnosticada».

10. El sistema deja disponible la cotización para que se gestione la
autorización del cliente. \\
\textbf{Flujos alternos} & A1 (paso 7): Las cuatro bahías están
ocupadas. El sistema muestra «Capacidad máxima alcanzada (4/4 bahías
ocupadas)» y coloca la orden en la lista de espera; al liberarse una
bahía, el jefe retoma la asignación.

A2: La orden se diagnostica pero queda en espera de asignación hasta que
exista capacidad disponible. \\
\textbf{Flujos de excepción} & E1: Credenciales incorrectas. El sistema
rechaza el acceso y permite volver a intentar.

E2: Pérdida de Wi‑Fi o sesión pausada al guardar. El sistema informa que
no hay conexión o que la sesión se pausó, conserva una copia local si
está disponible y solicita volver a iniciar sesión al restablecer la
red.

E3: Durante la reparación se detecta una falla adicional cuyo costo
incrementa más del 10 \% respecto de lo autorizado. El sistema detiene
el avance, cambia la orden a «Pendiente de Re-cotización» y deshabilita
nuevas requisiciones hasta obtener una nueva autorización mediante
2FA. \\
\textbf{Postcondiciones} & El diagnóstico y la cotización quedan
guardados; la orden cambia a «Diagnosticada». La asignación queda
establecida o la orden queda en lista de espera si todas las bahías
están ocupadas. \\
\textbf{Reglas de negocio} & Reglas de 4.7: RN-05 (Capacidad física del
taller), RN-09 (Asignación y colaboración de mecánicos), RN-04
(Re-cotización obligatoria por sobrecostos), RN-11 (Vigencia de
cotizaciones).

Validaciones del caso de uso:

• El diagnóstico se registra sobre una orden existente en estado
«Abierta».

• La asignación física requiere una bahía disponible y un mecánico
asignado.

• El taller dispone de cuatro bahías; si todas están ocupadas, se
utiliza lista de espera.

• Un incremento superior al 10 \% del costo autorizado requiere
re-cotización y nueva autorización del cliente antes de solicitar
refacciones adicionales. \\
\textbf{Requisitos relacionados} & RF-05, RF-06, RF-07 · HU-07, HU-08,
HU-09, HU-12, HU-16 \\
\end{longtable}
}

CU-06 --- Requisición de refacciones

{\def\LTcaptype{none} % do not increment counter
\begin{longtable}[]{@{}
  >{\raggedright\arraybackslash}p{(\linewidth - 2\tabcolsep) * \real{0.2387}}
  >{\raggedright\arraybackslash}p{(\linewidth - 2\tabcolsep) * \real{0.7613}}@{}}
\toprule\noalign{}
\begin{minipage}[b]{\linewidth}\raggedright
\textbf{Campo}
\end{minipage} & \begin{minipage}[b]{\linewidth}\raggedright
\textbf{Contenido}
\end{minipage} \\
\midrule\noalign{}
\endhead
\bottomrule\noalign{}
\endlastfoot
\textbf{ID y nombre} & CU-06 Requisición de refacciones \\
\textbf{Actor principal} & Jefe de Mecánicos (líder técnico) \\
\textbf{Actores secundarios} & Encargado de almacén; cliente (cuando se
requiere nueva autorización). \\
\textbf{Descripción breve} & Registra y envía al almacén la solicitud de
refacciones necesarias para ejecutar una reparación autorizada. \\
\textbf{Precondiciones} & 1. La orden está autorizada por el cliente.

2. La reparación se encuentra en proceso.

3. El jefe de mecánicos tiene acceso a la tableta y a la orden
activa. \\
\textbf{Flujo principal} & 1. El jefe de mecánicos inicia sesión y abre
el módulo de Requisición de Refacciones en la tableta.

2. Selecciona la orden de servicio activa.

3. Ingresa las refacciones requeridas para realizar el trabajo
autorizado.

4. Revisa la lista de piezas solicitadas.

5. Selecciona «Enviar Solicitud a Almacén».

6. El sistema registra la requisición y notifica al encargado de
almacén.

7. El encargado de almacén puede consultar la solicitud pendiente para
atenderla. \\
\textbf{Flujos alternos} & A1: Si las piezas solicitadas corresponden a
una re-cotización previamente autorizada por el cliente, el jefe
registra la requisición adicional y continúa el flujo normal. \\
\textbf{Flujos de excepción} & E1: Se detecta durante la reparación una
falla oculta que incrementa el costo en más del 10 \% del importe
autorizado. El sistema detiene la orden, la cambia a «Pendiente de
Re-cotización» y deshabilita la solicitud de nuevas refacciones. El jefe
prepara la cotización complementaria; tras una nueva autorización 2FA
por Recepción o por la Web, puede volver a requisar.

E2: Pérdida de conexión o sesión pausada al guardar. El sistema muestra
el aviso correspondiente y, si está disponible, conserva una copia local
para reintentar cuando se restablezca la conexión. \\
\textbf{Postcondiciones} & La requisición queda registrada y disponible
para el encargado de almacén, que recibe una notificación. Si se
requiere re-cotización, las nuevas piezas quedan bloqueadas hasta
obtener autorización. \\
\textbf{Reglas de negocio} & Reglas de 4.7: RN-04 (Re-cotización
obligatoria por sobrecostos), RN-06 (Reubicación por espera de
refacciones).

Validaciones del caso de uso:

• La requisición ordinaria corresponde a una orden autorizada y en
proceso de reparación.

• Las piezas adicionales asociadas a un aumento superior al 10 \%
requieren re-cotización y autorización 2FA.

• La solicitud enviada debe quedar asociada a la orden de servicio. \\
\textbf{Requisitos relacionados} & RF-08.1 · HU-10 \\
\end{longtable}
}

CU-07 --- Finalización de trabajo y sellado de historial

{\def\LTcaptype{none} % do not increment counter
\begin{longtable}[]{@{}
  >{\raggedright\arraybackslash}p{(\linewidth - 2\tabcolsep) * \real{0.2387}}
  >{\raggedright\arraybackslash}p{(\linewidth - 2\tabcolsep) * \real{0.7613}}@{}}
\toprule\noalign{}
\begin{minipage}[b]{\linewidth}\raggedright
\textbf{Campo}
\end{minipage} & \begin{minipage}[b]{\linewidth}\raggedright
\textbf{Contenido}
\end{minipage} \\
\midrule\noalign{}
\endhead
\bottomrule\noalign{}
\endlastfoot
\textbf{ID y nombre} & CU-07 Finalización de trabajo y sellado de
historial \\
\textbf{Actor principal} & Jefe de Mecánicos (líder técnico) \\
\textbf{Actores secundarios} & Mecánico asignado; cliente; servicio de
notificaciones. \\
\textbf{Descripción breve} & Registra las notas técnicas finales,
verifica que las tareas autorizadas estén completadas y finaliza el
trabajo para habilitar la entrega del vehículo. \\
\textbf{Precondiciones} & 1. El mecánico asignado terminó físicamente
las reparaciones.

2. La orden está en reparación y cuenta con diagnóstico y autorización
del cliente.

3. El jefe de mecánicos tiene acceso al sistema desde la tableta. \\
\textbf{Flujo principal} & 1. El jefe de mecánicos inicia sesión en la
tableta y supervisa el vehículo.

2. Abre el módulo de Cierre Técnico y Sellado de Historial.

3. Redacta las notas técnicas finales del trabajo realizado y verifica
su correspondencia con lo autorizado por el cliente.

4. Marca las tareas realizadas como completadas.

5. Selecciona «Finalizar Trabajo».

6. El sistema valida el cierre técnico.

7. El sistema cambia el estado de la orden a «Lista para Entrega».

8. El sistema envía una notificación automática al cliente. \\
\textbf{Flujos alternos} & A1: Si faltan tareas por marcar, el jefe
revisa el estado de las reparaciones, completa las tareas
correspondientes y vuelve a intentar finalizar el trabajo. \\
\textbf{Flujos de excepción} & E1: El jefe intenta cerrar la orden con
tareas autorizadas pendientes. El sistema muestra «Existen tareas
autorizadas pendientes por marcar como concluidas» y no finaliza la
orden.

E2: La sesión se pausa o se pierde la conexión Wi‑Fi antes de guardar.
El sistema informa el problema y permite reanudar el proceso tras
recuperar conexión y autenticarse. \\
\textbf{Postcondiciones} & Las notas técnicas y las tareas completadas
quedan registradas; la orden cambia a «Lista para Entrega» y se envía la
notificación al cliente. \\
\textbf{Reglas de negocio} & Reglas de 4.7: RN-03 (Selección granular de
rubros), RN-13 (Garantía de mano de obra).

Validaciones del caso de uso:

• Solo se finaliza técnicamente una reparación concluida.

• Las tareas completadas deben corresponder con lo autorizado por el
cliente.

• No se cambia la orden a «Lista para Entrega» si quedan tareas
autorizadas pendientes.

• La finalización genera una notificación al cliente. \\
\textbf{Requisitos relacionados} & RF-03.2, RF-05 · HU-20 \\
\end{longtable}
}

CU-08 --- Atender solicitud de refacciones y registrar compra

{\def\LTcaptype{none} % do not increment counter
\begin{longtable}[]{@{}
  >{\raggedright\arraybackslash}p{(\linewidth - 2\tabcolsep) * \real{0.2387}}
  >{\raggedright\arraybackslash}p{(\linewidth - 2\tabcolsep) * \real{0.7613}}@{}}
\toprule\noalign{}
\begin{minipage}[b]{\linewidth}\raggedright
\textbf{Campo}
\end{minipage} & \begin{minipage}[b]{\linewidth}\raggedright
\textbf{Contenido}
\end{minipage} \\
\midrule\noalign{}
\endhead
\bottomrule\noalign{}
\endlastfoot
\textbf{ID y nombre} & CU-08 Atender solicitud de refacciones y
registrar compra \\
\textbf{Actor principal} & Encargado de Almacén \\
\textbf{Actores secundarios} & Jefe de Mecánicos; proveedor local;
sistema de inventario. \\
\textbf{Descripción breve} & Atiende una solicitud de refacciones,
contacta a un proveedor, registra la compra y el ticket, y actualiza la
disponibilidad de las piezas para que continúe la reparación. \\
\textbf{Precondiciones} & 1. El jefe de mecánicos envió una solicitud de
refacciones asociada a una orden.

2. El encargado de almacén tiene acceso a la computadora y al
sistema. \\
\textbf{Flujo principal} & 1. El encargado de almacén inicia sesión en
la computadora de almacén.

2. Abre Solicitudes Pendientes y selecciona el pedido recibido.

3. Consulta el Directorio de Proveedores Locales y localiza un
proveedor.

4. Realiza la compra o pedido por teléfono o correo con el proveedor
seleccionado.

5. Abre el módulo de Registro de Compra y Ticket.

6. Captura los costos de las piezas, el proveedor y la fecha de
recepción.

7. Selecciona «Confirmar Compra y Generar Ticket».

8. El sistema genera el historial de compra, descuenta del inventario si
aplica y marca la refacción como «Recibida».

9. El sistema deja la solicitud registrada para que el taller continúe
la reparación. \\
\textbf{Flujos alternos} & A1: Si existe stock local disponible, el
encargado puede registrar el surtido de las piezas correspondientes sin
efectuar una nueva compra, actualizando el inventario y la solicitud
según corresponda. \\
\textbf{Flujos de excepción} & E1: No hay existencia en almacén y ningún
proveedor local tiene la pieza para entrega inmediata. El encargado
selecciona «Sin disponibilidad de proveedor»; el sistema cambia la orden
a «Detenida por Refacción» y avisa al jefe de mecánicos y al portal del
cliente. Cuando se confirma la llegada de la pieza, el encargado
registra el pedido y marca las piezas como recibidas.

E2: Credenciales incorrectas o tres intentos fallidos. El sistema
rechaza el acceso y bloquea temporalmente la cuenta conforme a la
política descrita.

E3: Cuenta abierta en otro dispositivo o sesión expirada. El sistema
cierra la sesión anterior o informa que la sesión expiró. \\
\textbf{Postcondiciones} & La compra queda en el historial, el ticket se
genera, el inventario se actualiza si aplica y las refacciones quedan
marcadas como «Recibida»; si no hay disponibilidad, la orden queda
«Detenida por Refacción». \\
\textbf{Reglas de negocio} & Reglas de 4.7: RN-06 (Reubicación por
espera de refacciones).

Validaciones del caso de uso:

• Toda compra debe asociarse a la solicitud y a su orden de servicio.

• Se registra el proveedor, los costos y la fecha de recepción.

• El inventario solo se descuenta o actualiza según aplique al
movimiento registrado.

• La falta de pieza disponible debe comunicarse al taller y al portal
del cliente. \\
\textbf{Requisitos relacionados} & RF-08.2, RF-08.3 · HU-17, HU-18 \\
\end{longtable}
}

CU-09 --- Solicitud de cita desde la web

{\def\LTcaptype{none} % do not increment counter
\begin{longtable}[]{@{}
  >{\raggedright\arraybackslash}p{(\linewidth - 2\tabcolsep) * \real{0.2351}}
  >{\raggedright\arraybackslash}p{(\linewidth - 2\tabcolsep) * \real{0.7649}}@{}}
\toprule\noalign{}
\begin{minipage}[b]{\linewidth}\raggedright
\textbf{Campo}
\end{minipage} & \begin{minipage}[b]{\linewidth}\raggedright
\textbf{Contenido}
\end{minipage} \\
\midrule\noalign{}
\endhead
\bottomrule\noalign{}
\endlastfoot
\textbf{ID y nombre} & CU-09 Solicitud de cita desde la web \\
\textbf{Actor principal} & Cliente \\
\textbf{Actores secundarios} & Sistema web; recepcionista (revisa
posteriormente la solicitud). \\
\textbf{Descripción breve} & Permite que el cliente solicite una cita en
el sitio web público sin crear una cuenta ni iniciar sesión,
proporcionando sus datos, los de su vehículo y el horario deseado. \\
\textbf{Precondiciones} & 1. El cliente puede acceder a la página web
pública del taller.

2. El servicio web de solicitud de citas está disponible. \\
\textbf{Flujo principal} & 1. El cliente abre la pantalla de Solicitud
de Cita desde el sitio web, sin registrarse ni introducir contraseña.

2. Llena sus datos personales y la información de su automóvil.

3. Selecciona la fecha y hora deseadas.

4. Revisa la información y selecciona «Solicitar Cita».

5. El sistema registra la solicitud como pendiente de revisión.

6. El sistema muestra un mensaje que indica que Recepción revisará la
petición. \\
\textbf{Flujos alternos} & A1: Si la fecha u hora solicitada no tiene
espacio disponible, Recepción revisa la solicitud y puede proponer al
cliente un día u horario alternativo antes de confirmar la cita. \\
\textbf{Flujos de excepción} & E1: No se puede enviar o registrar la
solicitud. El material proporcionado no especifica el mensaje ni el
procedimiento concreto de recuperación; la solicitud no debe
considerarse recibida hasta que el sistema confirme el registro. \\
\textbf{Postcondiciones} & La solicitud queda registrada para revisión
de Recepción y el cliente recibe confirmación en pantalla de que su
petición fue enviada. \\
\textbf{Reglas de negocio} & Reglas de 4.7: RN-08 (Confirmación manual
de citas web).

Validaciones del caso de uso:

• El cliente puede solicitar la cita sin usuario ni contraseña.

• La solicitud requiere datos del cliente, información del vehículo y
fecha/hora deseadas.

• Enviar una solicitud no equivale a confirmar la cita; Recepción debe
revisarla. \\
\textbf{Requisitos relacionados} & RF-09 · HU-03 \\
\end{longtable}
}

CU-10 --- Consulta del estado del vehículo en vivo

{\def\LTcaptype{none} % do not increment counter
\begin{longtable}[]{@{}
  >{\raggedright\arraybackslash}p{(\linewidth - 2\tabcolsep) * \real{0.2351}}
  >{\raggedright\arraybackslash}p{(\linewidth - 2\tabcolsep) * \real{0.7649}}@{}}
\toprule\noalign{}
\begin{minipage}[b]{\linewidth}\raggedright
\textbf{Campo}
\end{minipage} & \begin{minipage}[b]{\linewidth}\raggedright
\textbf{Contenido}
\end{minipage} \\
\midrule\noalign{}
\endhead
\bottomrule\noalign{}
\endlastfoot
\textbf{ID y nombre} & CU-10 Consulta del estado del vehículo en vivo \\
\textbf{Actor principal} & Cliente \\
\textbf{Actores secundarios} & Sistema web; servicio de generación y
validación de enlaces seguros. \\
\textbf{Descripción breve} & Permite consultar el avance del automóvil
desde un teléfono mediante un enlace seguro asociado a la orden, sin
requerir usuario, contraseña ni descarga de una aplicación. \\
\textbf{Precondiciones} & 1. El vehículo fue recibido en el taller y
cuenta con una orden de servicio.

2. El cliente recibió por correo o SMS el enlace seguro
correspondiente. \\
\textbf{Flujo principal} & 1. El cliente abre desde su teléfono el
enlace seguro recibido.

2. El sistema valida automáticamente el token incluido en el enlace.

3. El sistema muestra la pantalla de Seguimiento en Vivo.

4. El cliente visualiza la barra de progreso y el estado actual del
vehículo.

5. El sistema presenta el avance disponible, que puede incluir
Recepción, Diagnóstico, En Reparación y Listo. \\
\textbf{Flujos alternos} & A1: El cliente vuelve a abrir el mismo enlace
vigente para consultar el estado actualizado de la orden. \\
\textbf{Flujos de excepción} & E1: El enlace ha expirado o el token fue
alterado. El sistema bloquea el acceso y muestra: «Su enlace de acceso
ha expirado o es inválido. Por favor, solicite un nuevo enlace de
seguimiento a la recepción del taller». \\
\textbf{Postcondiciones} & El cliente visualiza el estado asociado a la
orden si el token es válido. No se modifican los datos ni el estado de
la orden mediante esta consulta. \\
\textbf{Reglas de negocio} & Reglas de 4.7: RN-01 (Expiración y vigencia
de enlaces).

Validaciones del caso de uso:

• El acceso del cliente se realiza mediante un token seguro, sin usuario
ni contraseña.

• Un enlace expirado o alterado no debe permitir consultar la
información.

• La información mostrada corresponde a la orden asociada al enlace. \\
\textbf{Requisitos relacionados} & RF-02.1 · HU-28 \\
\end{longtable}
}

CU-11 --- Aprobación remota de cotización (2FA)

{\def\LTcaptype{none} % do not increment counter
\begin{longtable}[]{@{}
  >{\raggedright\arraybackslash}p{(\linewidth - 2\tabcolsep) * \real{0.2317}}
  >{\raggedright\arraybackslash}p{(\linewidth - 2\tabcolsep) * \real{0.7683}}@{}}
\toprule\noalign{}
\begin{minipage}[b]{\linewidth}\raggedright
\textbf{Campo}
\end{minipage} & \begin{minipage}[b]{\linewidth}\raggedright
\textbf{Contenido}
\end{minipage} \\
\midrule\noalign{}
\endhead
\bottomrule\noalign{}
\endlastfoot
\textbf{ID y nombre} & CU-11 Aprobación remota de cotización (2FA) \\
\textbf{Actor principal} & Cliente \\
\textbf{Actores secundarios} & Recepcionista; servicio de correo
electrónico; sistema web. \\
\textbf{Descripción breve} & Permite que el cliente revise y autorice a
distancia la cotización técnica mediante un enlace seguro y un código
2FA de seis dígitos enviado por correo electrónico. \\
\textbf{Precondiciones} & 1. El vehículo ya fue diagnosticado.

2. La cotización fue enviada al cliente mediante un enlace seguro.

3. El enlace está vigente y asociado a la orden correcta.

4. El cliente puede acceder al correo registrado. \\
\textbf{Flujo principal} & 1. El cliente abre el enlace dinámico desde
su teléfono.

2. El sistema valida el enlace y muestra el desglose de mano de obra y
refacciones.

3. El cliente revisa los conceptos y el importe cotizado.

4. Presiona «Solicitar Código 2FA para Autorizar».

5. El sistema envía por correo un código de verificación de seis
dígitos.

6. El cliente introduce el código en la página web.

7. Presiona «Aceptar y Autorizar».

8. El sistema valida el código y registra el monto, la fecha y la hora
de la autorización.

9. El sistema cambia la orden a «Autorizada» y muestra el resultado. \\
\textbf{Flujos alternos} & A1 (paso 3): El cliente no está de acuerdo
con el costo y selecciona «Rechazar Cotización». El sistema solicita
confirmar el rechazo; tras la confirmación cambia la orden a
«Cancelada/Rechazada», notifica a Recepción para coordinar la entrega
del vehículo sin reparar y finaliza la sesión web. \\
\textbf{Flujos de excepción} & E1: Enlace expirado o token alterado. El
sistema bloquea el acceso e indica al cliente que solicite un nuevo
enlace a Recepción.

E2 (paso 8): Código 2FA incorrecto o expirado después de cinco minutos.
El sistema rechaza la autorización e indica que debe solicitarse un
código nuevo.

E3: El código no llega al correo. El material indica que el cliente debe
solicitar un nuevo código; el flujo no concluye hasta que la
verificación sea satisfactoria. \\
\textbf{Postcondiciones} & Si la autorización es válida, queda
registrada la cotización aprobada y la orden pasa a «Autorizada». Si el
cliente la rechaza, la orden queda «Cancelada/Rechazada» y Recepción
recibe aviso para coordinar la entrega sin reparar. \\
\textbf{Reglas de negocio} & Reglas de 4.7: RN-02 (Doble verificación en
autorizaciones), RN-03 (Selección granular de rubros), RN-11 (Vigencia
de cotizaciones).

Validaciones del caso de uso:

• La cotización debe estar asociada a una orden diagnosticada.

• El enlace seguro debe ser vigente y válido.

• El código 2FA es de seis dígitos y su vigencia indicada es de cinco
minutos.

• La orden solo se autoriza después de validar el código y registrar el
consentimiento del cliente.

• El rechazo confirmado debe notificarse a Recepción. \\
\textbf{Requisitos relacionados} & RF-02.1, RF-06 · HU-02, HU-13, HU-14,
HU-15 \\
\end{longtable}
}

CU-12 --- Monitoreo global del taller y KPIs

{\def\LTcaptype{none} % do not increment counter
\begin{longtable}[]{@{}
  >{\raggedright\arraybackslash}p{(\linewidth - 2\tabcolsep) * \real{0.2317}}
  >{\raggedright\arraybackslash}p{(\linewidth - 2\tabcolsep) * \real{0.7683}}@{}}
\toprule\noalign{}
\begin{minipage}[b]{\linewidth}\raggedright
\textbf{Campo}
\end{minipage} & \begin{minipage}[b]{\linewidth}\raggedright
\textbf{Contenido}
\end{minipage} \\
\midrule\noalign{}
\endhead
\bottomrule\noalign{}
\endlastfoot
\textbf{ID y nombre} & CU-12 Monitoreo global del taller y KPIs \\
\textbf{Actor principal} & Gerente (dueño) \\
\textbf{Actores secundarios} & Sistema de gestión del taller. \\
\textbf{Descripción breve} & Proporciona al gerente un tablero general
para visualizar indicadores operativos y financieros del taller,
incluyendo ocupación de bahías, órdenes abiertas, tiempos promedio e
ingresos acumulados. \\
\textbf{Precondiciones} & 1. El gerente requiere revisar la operación
general.

2. Cuenta con perfil gerencial y credenciales válidas.

3. Existen datos operativos registrados en el sistema. \\
\textbf{Flujo principal} & 1. El gerente inicia sesión en su computadora
o teléfono con el perfil gerencial.

2. El sistema valida las credenciales y abre el Dashboard General de
Control.

3. El gerente consulta la ocupación de las cuatro bahías.

4. El gerente visualiza las órdenes abiertas.

5. El gerente consulta los tiempos promedio de estancia.

6. El gerente visualiza los ingresos acumulados y los indicadores
disponibles en tiempo real. \\
\textbf{Flujos alternos} & A1: El gerente consulta el tablero desde otro
dispositivo autorizado y revisa los mismos indicadores según los
permisos de su perfil. \\
\textbf{Flujos de excepción} & E1: Credenciales incorrectas u olvido de
contraseña. El sistema muestra un error y permite iniciar el flujo de
recuperación por correo y verificación de seguridad.

E2: Cambio de red o inactividad. Para proteger la información operativa
y financiera, el sistema requiere reautenticación mediante huella o
contraseña, de acuerdo con el flujo descrito en el material. \\
\textbf{Postcondiciones} & El gerente visualiza los indicadores
disponibles del taller. La consulta no modifica las órdenes ni los datos
operativos. \\
\textbf{Reglas de negocio} & Sin regla de 4.7 específica.

Validaciones del caso de uso:

• Solo un perfil gerencial autorizado puede consultar el Dashboard
General de Control.

• Los indicadores mostrados deben basarse en los datos registrados en el
sistema.

• Tras inactividad o cambio de red, el sistema puede exigir
reautenticación para proteger la información. \\
\textbf{Requisitos relacionados} & RF-11 · HU-29 \\
\end{longtable}
}

CU-13 --- Autorización de entrega a crédito (flotillas)

{\def\LTcaptype{none} % do not increment counter
\begin{longtable}[]{@{}
  >{\raggedright\arraybackslash}p{(\linewidth - 2\tabcolsep) * \real{0.2313}}
  >{\raggedright\arraybackslash}p{(\linewidth - 2\tabcolsep) * \real{0.7687}}@{}}
\toprule\noalign{}
\begin{minipage}[b]{\linewidth}\raggedright
\textbf{Campo}
\end{minipage} & \begin{minipage}[b]{\linewidth}\raggedright
\textbf{Contenido}
\end{minipage} \\
\midrule\noalign{}
\endhead
\bottomrule\noalign{}
\endlastfoot
\textbf{ID y nombre} & CU-13 Autorización de entrega a crédito
(flotillas) \\
\textbf{Actor principal} & Gerente (dueño) \\
\textbf{Actores secundarios} & Recepcionista; cliente
empresarial/flotilla; sistema de control de crédito. \\
\textbf{Descripción breve} & Permite decidir si se libera a crédito un
vehículo de una flotilla empresarial que terminó su reparación pero
conserva saldo pendiente, con revisión de la línea de crédito de la
empresa. \\
\textbf{Precondiciones} & 1. El vehículo de una flotilla empresarial
terminó su reparación.

2. La orden conserva saldo pendiente.

3. Recepción solicitó autorización de crédito para liberar el vehículo.

4. El gerente tiene un perfil autorizado para aprobar excepciones. \\
\textbf{Flujo principal} & 1. La recepcionista solicita al gerente
autorización de entrega a crédito.

2. El gerente inicia sesión y abre Aprobación de Créditos y Excepciones.

3. Selecciona la orden correspondiente a la flotilla empresarial.

4. El sistema muestra la información de la orden y la línea de crédito
aprobada.

5. El gerente revisa la información de crédito y decide autorizar la
entrega.

6. El gerente selecciona «Autorizar Entrega a Crédito».

7. El sistema registra la autorización y habilita a Recepción para
entregar el vehículo sin pago inmediato de contado. \\
\textbf{Flujos alternos} & \begin{minipage}[t]{\linewidth}\raggedright
A1: La flotilla excede su límite autorizado o presenta facturas
vencidas. El sistema muestra una alerta con el saldo vencido y pregunta
si el gerente desea forzar la autorización bajo su responsabilidad.

\begin{quote}
• Si acepta, el sistema registra la excepción gerencial y habilita la
entrega.

• Si rechaza, la entrega continúa bloqueada y se exige liquidar el saldo
pendiente.
\end{quote}
\end{minipage} \\
\textbf{Flujos de excepción} & E1: Credenciales gerenciales incorrectas
o cuenta inaccesible. El gerente debe corregir los datos o recuperar la
contraseña por correo antes de continuar.

E2: El sistema detecta que la empresa excedió el límite de crédito. Si
no existe autorización expresa del gerente, la entrega permanece
bloqueada. \\
\textbf{Postcondiciones} & Si se autoriza, queda registrada la excepción
de crédito y Recepción puede continuar con la entrega. Si se deniega, el
vehículo no se libera hasta liquidar el saldo o conseguir autorización
válida. \\
\textbf{Reglas de negocio} & Reglas de 4.7: RN-10 (Restricción de
entrega y excepción a flotillas).

Validaciones del caso de uso:

• Esta excepción aplica a vehículos de flotillas empresariales.

• Se debe revisar la línea de crédito antes de autorizar.

• El gerente puede forzar una autorización con límite excedido bajo su
responsabilidad, según el flujo definido.

• Si el gerente rechaza la excepción, la entrega permanece bloqueada y
se exige pago del saldo. \\
\textbf{Requisitos relacionados} & RF-10 · HU-22 \\
\end{longtable}
}

\subsubsection{4.9.1 Flujos transversales de autenticación y
sesión}\label{flujos-transversales-de-autenticaciuxf3n-y-sesiuxf3n}

Estos flujos se presentan como apoyo transversal de seguridad, no como
casos de uso de negocio adicionales a los 13 anteriores. Los
comportamientos se transcriben y ordenan a partir del material
entregado.

FT-01 --- Autenticación y sesión de Recepcionista

{\def\LTcaptype{none} % do not increment counter
\begin{longtable}[]{@{}
  >{\raggedright\arraybackslash}p{(\linewidth - 2\tabcolsep) * \real{0.2333}}
  >{\raggedright\arraybackslash}p{(\linewidth - 2\tabcolsep) * \real{0.7667}}@{}}
\toprule\noalign{}
\begin{minipage}[b]{\linewidth}\raggedright
\textbf{Campo}
\end{minipage} & \begin{minipage}[b]{\linewidth}\raggedright
\textbf{Contenido}
\end{minipage} \\
\midrule\noalign{}
\endhead
\bottomrule\noalign{}
\endlastfoot
\textbf{Actor / flujo} & Actor: Recepcionista

Precondición: la recepcionista intenta acceder desde la computadora de
mostrador.

Flujo normal: ingresa usuario/correo y contraseña; el sistema valida los
datos y muestra la pantalla principal de Recepción.

Excepciones: (1) Credenciales incorrectas: muestra «Usuario o contraseña
incorrectos. Intentos restantes: 2». (2) Tres intentos fallidos
consecutivos: bloquea la cuenta por 15 minutos y muestra aviso de
bloqueo. (3) Recuperación: solicita correo corporativo; si no existe
muestra «Correo no registrado»; si existe, envía enlace/token válido por
15 minutos; la recepcionista restablece la contraseña y vuelve al login.
(4) Más de 30 minutos de inactividad: expira la sesión, guarda el
borrador del formulario actual, redirige al login y muestra aviso de
sesión expirada. \\
\end{longtable}
}

FT-02 --- Autenticación y sesión del Jefe de Mecánicos

{\def\LTcaptype{none} % do not increment counter
\begin{longtable}[]{@{}
  >{\raggedright\arraybackslash}p{(\linewidth - 2\tabcolsep) * \real{0.2299}}
  >{\raggedright\arraybackslash}p{(\linewidth - 2\tabcolsep) * \real{0.7701}}@{}}
\toprule\noalign{}
\begin{minipage}[b]{\linewidth}\raggedright
\textbf{Campo}
\end{minipage} & \begin{minipage}[b]{\linewidth}\raggedright
\textbf{Contenido}
\end{minipage} \\
\midrule\noalign{}
\endhead
\bottomrule\noalign{}
\endlastfoot
\textbf{Actor / flujo} & Actor: Jefe de Mecánicos

Precondición: el jefe intenta acceder desde la tableta del taller.

Flujo normal: introduce usuario y contraseña y accede al tablero del
taller.

Excepciones: (1) Error de captura o credenciales: se rechaza el acceso y
se permite reintentar. (2) Recuperación: ingresa el correo, recibe un
token de restablecimiento, crea una nueva contraseña y vuelve al login.
(3) Pérdida de Wi‑Fi o cierre por inactividad durante el guardado:
muestra «Sin conexión / Sesión pausada. Se guardó una copia local.
Vuelva a iniciar sesión al restablecer la red». \\
\end{longtable}
}

FT-03 --- Autenticación y sesión del Encargado de Almacén

{\def\LTcaptype{none} % do not increment counter
\begin{longtable}[]{@{}
  >{\raggedright\arraybackslash}p{(\linewidth - 2\tabcolsep) * \real{0.2290}}
  >{\raggedright\arraybackslash}p{(\linewidth - 2\tabcolsep) * \real{0.7710}}@{}}
\toprule\noalign{}
\begin{minipage}[b]{\linewidth}\raggedright
\textbf{Campo}
\end{minipage} & \begin{minipage}[b]{\linewidth}\raggedright
\textbf{Contenido}
\end{minipage} \\
\midrule\noalign{}
\endhead
\bottomrule\noalign{}
\endlastfoot
\textbf{Actor / flujo} & Actor: Encargado de Almacén

Precondición: el almacenista intenta acceder a la terminal de
inventarios.

Flujo normal: introduce credenciales y accede a la pantalla de
solicitudes.

Excepciones: (1) Tres contraseñas fallidas: bloqueo temporal y
notificación al gerente/sistema. (2) Recuperación por correo corporativo
y restablecimiento de contraseña. (3) Inicio en dos computadoras o
expiración de sesión por inactividad: el sistema cierra la sesión
anterior y muestra «Se ha iniciado sesión en otro dispositivo o la
sesión ha expirado». \\
\end{longtable}
}

FT-04 --- Acceso del Cliente mediante enlace seguro y token

{\def\LTcaptype{none} % do not increment counter
\begin{longtable}[]{@{}
  >{\raggedright\arraybackslash}p{(\linewidth - 2\tabcolsep) * \real{0.2281}}
  >{\raggedright\arraybackslash}p{(\linewidth - 2\tabcolsep) * \real{0.7719}}@{}}
\toprule\noalign{}
\begin{minipage}[b]{\linewidth}\raggedright
\textbf{Campo}
\end{minipage} & \begin{minipage}[b]{\linewidth}\raggedright
\textbf{Contenido}
\end{minipage} \\
\midrule\noalign{}
\endhead
\bottomrule\noalign{}
\endlastfoot
\textbf{Actor / flujo} & Actor: Cliente

Precondición: el cliente intenta consultar el estado del vehículo o
autorizar una cotización desde el celular. No requiere usuario ni
contraseña: la validación se realiza con el token del enlace único.

Excepción: enlace expirado (el enlace expira al entregar y saldar el
vehículo, RN-01) o token alterado. El sistema bloquea el acceso e
indica: «Su enlace de acceso ha expirado o es inválido. Por favor,
solicite un nuevo enlace de seguimiento a la recepción del taller».

Cotización vencida: si pasaron los 5 días naturales de vigencia (RN-11),
el enlace sigue activo, pero la cotización se muestra como «Vencida» y
no puede autorizarse; el cliente pide una nueva versión a Recepción.

Excepción 2FA: código incorrecto o expirado después de cinco minutos. El
sistema rechaza la autorización y solicita un nuevo código. El cliente
también puede rechazar expresamente la cotización; tras confirmar, el
sistema cambia la orden a «Cancelada/Rechazada», notifica a Recepción y
finaliza la sesión web. \\
\end{longtable}
}

FT-05 --- Autenticación y sesión del Gerente

{\def\LTcaptype{none} % do not increment counter
\begin{longtable}[]{@{}
  >{\raggedright\arraybackslash}p{(\linewidth - 2\tabcolsep) * \real{0.2248}}
  >{\raggedright\arraybackslash}p{(\linewidth - 2\tabcolsep) * \real{0.7752}}@{}}
\toprule\noalign{}
\begin{minipage}[b]{\linewidth}\raggedright
\textbf{Campo}
\end{minipage} & \begin{minipage}[b]{\linewidth}\raggedright
\textbf{Contenido}
\end{minipage} \\
\midrule\noalign{}
\endhead
\bottomrule\noalign{}
\endlastfoot
\textbf{Actor / flujo} & Actor: Gerente (dueño)

Precondición: el gerente intenta ingresar al dashboard desde su laptop o
celular.

Flujo normal: inicia sesión con sus credenciales y accede al Dashboard
General.

Excepciones: (1) Credenciales olvidadas o incorrectas: puede recuperar
la cuenta mediante correo personal/gerencial, completar la verificación
de seguridad, cambiar la contraseña y volver a iniciar sesión. (2)
Cambio de red Wi‑Fi o inactividad: el sistema exige reautenticación
mediante huella o contraseña para proteger métricas financieras y
operativas. \\
\end{longtable}
}

\subsection{4.10 Prototipos de baja
fidelidad}\label{prototipos-de-baja-fidelidad}

\begin{quote}
\emph{\textbf{Instrucción:} Wireframes de todas las pantallas del flujo
núcleo, cada uno con la historia o el caso de uso que representa.}
\end{quote}

{[}Figuras 4.x. Wireframes.{]}

\subsection{4.11 Validación de requisitos con el
cliente}\label{validaciuxf3n-de-requisitos-con-el-cliente}

\begin{quote}
\emph{\textbf{Instrucción:} Resumen de las sesiones de validación: qué
se mostró, qué cambió y qué acordó el cliente. Incluyan los defectos de
documentación que reportó el otro equipo en la revisión cruzada y cómo
se atendieron.}
\end{quote}

{\def\LTcaptype{none} % do not increment counter
\begin{longtable}[]{@{}
  >{\raggedright\arraybackslash}p{(\linewidth - 6\tabcolsep) * \real{0.1497}}
  >{\raggedright\arraybackslash}p{(\linewidth - 6\tabcolsep) * \real{0.2564}}
  >{\raggedright\arraybackslash}p{(\linewidth - 6\tabcolsep) * \real{0.3161}}
  >{\raggedright\arraybackslash}p{(\linewidth - 6\tabcolsep) * \real{0.2778}}@{}}
\toprule\noalign{}
\begin{minipage}[b]{\linewidth}\raggedright
\textbf{Fecha}
\end{minipage} & \begin{minipage}[b]{\linewidth}\raggedright
\textbf{Técnica (recorrido, prototipo, revisión)}
\end{minipage} & \begin{minipage}[b]{\linewidth}\raggedright
\textbf{Hallazgos}
\end{minipage} & \begin{minipage}[b]{\linewidth}\raggedright
\textbf{Cambios acordados}
\end{minipage} \\
\midrule\noalign{}
\endhead
\bottomrule\noalign{}
\endlastfoot
& & & \\
& & & \\
\end{longtable}
}

\section{5. Plan de pruebas y
trazabilidad}\label{plan-de-pruebas-y-trazabilidad}

\subsection{5.1 Plan de pruebas}\label{plan-de-pruebas}

\begin{quote}
\emph{\textbf{Instrucción:} Estructura adaptada de ISO/IEC/IEEE
29119-3:2021.}
\end{quote}

{\def\LTcaptype{none} % do not increment counter
\begin{longtable}[]{@{}
  >{\raggedright\arraybackslash}p{(\linewidth - 2\tabcolsep) * \real{0.2778}}
  >{\raggedright\arraybackslash}p{(\linewidth - 2\tabcolsep) * \real{0.7222}}@{}}
\toprule\noalign{}
\begin{minipage}[b]{\linewidth}\raggedright
\textbf{Apartado}
\end{minipage} & \begin{minipage}[b]{\linewidth}\raggedright
\textbf{Contenido}
\end{minipage} \\
\midrule\noalign{}
\endhead
\bottomrule\noalign{}
\endlastfoot
Alcance & {[}Qué se prueba y qué no en la Fase 1{]} \\
Niveles de prueba & {[}Unitarias, integración, sistema, aceptación{]} \\
Enfoque & {[}Manual / automatizado; herramientas{]} \\
Entornos & {[}Local, CI, staging; datos ficticios{]} \\
Criterios de entrada & \\
Criterios de salida & {[}Por ejemplo: 0 defectos S1/S2 abiertos{]} \\
Roles & {[}QA Lead, developers, equipo contrario, cliente{]} \\
Riesgos de las pruebas & \\
Calendario & {[}Por sprint{]} \\
\end{longtable}
}

\subsection{5.2 Casos de prueba de
aceptación}\label{casos-de-prueba-de-aceptaciuxf3n}

\begin{quote}
\emph{\textbf{Instrucción:} Mínimo \textbf{25 casos (Equipo A)} o
\textbf{20 (Equipo B)}, cada uno vinculado a una historia o caso de uso.
Incluyan casos de flujos alternos, valores límite y control de acceso
por rol. La columna de resultado se llena al ejecutarlos (sección 10).}
\end{quote}

{\def\LTcaptype{none} % do not increment counter
\begin{longtable}[]{@{}
  >{\raggedright\arraybackslash}p{(\linewidth - 10\tabcolsep) * \real{0.0963}}
  >{\raggedright\arraybackslash}p{(\linewidth - 10\tabcolsep) * \real{0.1282}}
  >{\raggedright\arraybackslash}p{(\linewidth - 10\tabcolsep) * \real{0.1816}}
  >{\raggedright\arraybackslash}p{(\linewidth - 10\tabcolsep) * \real{0.2308}}
  >{\raggedright\arraybackslash}p{(\linewidth - 10\tabcolsep) * \real{0.2136}}
  >{\raggedright\arraybackslash}p{(\linewidth - 10\tabcolsep) * \real{0.1496}}@{}}
\toprule\noalign{}
\begin{minipage}[b]{\linewidth}\raggedright
\textbf{ID}
\end{minipage} & \begin{minipage}[b]{\linewidth}\raggedright
\textbf{Historia / CU}
\end{minipage} & \begin{minipage}[b]{\linewidth}\raggedright
\textbf{Precondiciones}
\end{minipage} & \begin{minipage}[b]{\linewidth}\raggedright
\textbf{Pasos}
\end{minipage} & \begin{minipage}[b]{\linewidth}\raggedright
\textbf{Resultado esperado}
\end{minipage} & \begin{minipage}[b]{\linewidth}\raggedright
\textbf{Resultado (Pasa / Falla / Bloqueado)}
\end{minipage} \\
\midrule\noalign{}
\endhead
\bottomrule\noalign{}
\endlastfoot
CPA-01 & HU-\ldots{} & & 1. \ldots{}

2. \ldots{} & & \\
& & & & & \\
& & & & & \\
\end{longtable}
}

\subsection{5.3 Matriz de trazabilidad}\label{matriz-de-trazabilidad}

\begin{quote}
\emph{\textbf{Instrucción:} Ningún requisito Must puede quedar sin
historia ni caso de prueba. Se actualiza en cada versión y en el Sprint
6 incluye el resultado de las pruebas.}
\end{quote}

{\def\LTcaptype{none} % do not increment counter
\begin{longtable}[]{@{}
  >{\raggedright\arraybackslash}p{(\linewidth - 10\tabcolsep) * \real{0.1390}}
  >{\raggedright\arraybackslash}p{(\linewidth - 10\tabcolsep) * \real{0.1708}}
  >{\raggedright\arraybackslash}p{(\linewidth - 10\tabcolsep) * \real{0.1389}}
  >{\raggedright\arraybackslash}p{(\linewidth - 10\tabcolsep) * \real{0.1923}}
  >{\raggedright\arraybackslash}p{(\linewidth - 10\tabcolsep) * \real{0.2094}}
  >{\raggedright\arraybackslash}p{(\linewidth - 10\tabcolsep) * \real{0.1496}}@{}}
\toprule\noalign{}
\begin{minipage}[b]{\linewidth}\raggedright
\textbf{Requisito}
\end{minipage} & \begin{minipage}[b]{\linewidth}\raggedright
\textbf{Historias}
\end{minipage} & \begin{minipage}[b]{\linewidth}\raggedright
\textbf{Caso de uso}
\end{minipage} & \begin{minipage}[b]{\linewidth}\raggedright
\textbf{Casos de prueba}
\end{minipage} & \begin{minipage}[b]{\linewidth}\raggedright
\textbf{Componente / módulo}
\end{minipage} & \begin{minipage}[b]{\linewidth}\raggedright
\textbf{Resultado}
\end{minipage} \\
\midrule\noalign{}
\endhead
\bottomrule\noalign{}
\endlastfoot
RF-01 & HU-01, HU-02 & CU-01 & CPA-01, CPA-02 & {[}Paquete o clase{]} &
{[}Pasa{]} \\
& & & & & \\
& & & & & \\
\end{longtable}
}

\section{6. Modelos de análisis}\label{modelos-de-anuxe1lisis}

Los modelos usan los nombres del glosario (4.2) y del SRS, y coinciden con el esquema
\texttt{db/schema.sql} y con el código de \texttt{backend/app/}. Notación UML 2.5.1.

\subsection{6.1 Diagramas de flujo de
datos}\label{diagramas-de-flujo-de-datos}

\figurapendiente{6.1}{Diagrama de flujo de datos de contexto}{el sistema como un único proceso
«Sistema de Órdenes de Servicio y Citas» y sus entidades externas: Cliente (envía solicitud de cita,
selección de conceptos, código 2FA, pago; recibe enlace seguro, cotización, código y estado), Recepcionista
(datos de recepción, checklist, fotografías, pagos, entrega), Jefe de Mecánicos (asignación, diagnóstico,
cotización, fin de trabajo), Gerente (alta de personal, crédito a flotillas; recibe panel) y Servicio de
correo/SMS (recibe mensajes salientes).}{docs/06-analisis/dfd/contexto.drawio}

\figurapendiente{6.2}{Diagrama de flujo de datos de nivel 0}{los procesos 1.0 Gestionar citas, 2.0
Recibir vehículo, 3.0 Asignar y diagnosticar, 4.0 Cotizar y autorizar, 5.0 Cobrar y entregar y 6.0
Consultar estado/panel; y los almacenes D1 Clientes y vehículos, D2 Citas, D3 Órdenes e historial, D4
Cotizaciones y autorizaciones, D5 Pagos y garantías, D6 Tokens de acceso. Los flujos con las entidades
externas deben ser los mismos que en la Figura 6.1 (balanceo).}{docs/06-analisis/dfd/nivel-0.drawio}

\figurapendiente{6.3}{Diagrama de flujo de datos de nivel 1 del proceso 4.0 Cotizar y
autorizar}{subprocesos 4.1 Generar cotización (lee D3, escribe D4), 4.2 Enviar enlace (escribe D6, envía
al servicio de correo), 4.3 Seleccionar conceptos opcionales, 4.4 Emitir código 2FA, 4.5 Validar código y
registrar autorización (escribe D4 y el historial en D3), 4.6 Rechazar cotización. Entradas y salidas
balanceadas con el proceso 4.0 de la Figura 6.2.}{docs/06-analisis/dfd/nivel-1-cotizar.drawio}

\subsection{6.2 Modelo de datos}\label{modelo-de-datos}

El modelo relacional está en 3FN: cada tabla describe una sola entidad, los atributos derivados del
diagrama de clases (/total, /saldoPendiente, /fechaVencimiento, /fechaFin) no se almacenan sino que se
calculan en vistas o columnas generadas, y la especialización \{disjoint, complete\} de Cliente usa una
tabla por subclase. El esquema completo, con sus restricciones y disparadores, está en
\texttt{db/schema.sql} y se valida en PostgreSQL 16 en cada ejecución de la integración continua.

\figurapendiente{6.4}{Modelo entidad-relación}{las 26 tablas de \texttt{db/schema.sql} con sus llaves
primarias y foráneas y la cardinalidad de cada relación; resaltar el núcleo cliente --- vehículo ---
orden\_servicio --- recepcion\_vehicular (checklist, fotografías) --- cotizacion (detalles,
autorizaciones, pagos) --- token\_acceso. Puede generarse con DBeaver o pgAdmin a partir del
esquema.}{db/schema.sql (exportar a docs/06-analisis/er/modelo-er.png)}

\textbf{Diccionario de datos} de las tablas del flujo núcleo (las demás se documentan en los
comentarios de \texttt{db/schema.sql}):

{\footnotesize
\begin{longtable}{@{}L{0.15\linewidth}L{0.17\linewidth}L{0.13\linewidth}L{0.06\linewidth}L{0.2\linewidth}L{0.21\linewidth}@{}}
\toprule
\textbf{Entidad} & \textbf{Atributo} & \textbf{Tipo} & \textbf{Nulo} & \textbf{Restricciones} & \textbf{Descripción} \\
\midrule
\endhead
\bottomrule
\endlastfoot
cliente & id\_cliente & INTEGER & No & PK, identidad & Identificador del cliente \\
 & tipo\_cliente & VARCHAR(10) & No & Particular o Flotilla & Subclase del cliente \\
 & telefono & VARCHAR(20) & No & & Contacto (dato personal) \\
 & correo & VARCHAR(160) & No & & Recibe enlace y códigos 2FA \\
vehiculo & id\_vehiculo & INTEGER & No & PK & Identificador \\
 & id\_cliente & INTEGER & No & FK cliente & Propietario \\
 & placas & VARCHAR(15) & No & Único & Placas en mayúsculas \\
 & vin & VARCHAR(17) & Sí & Único & Número de serie (opcional) \\
 & tipo & tipo\_vehiculo & No & Automovil, Motocicleta, Camioneta & Solo vehículos ligeros (RN-07) \\
orden\_servicio & folio & VARCHAR(20) & No & PK, formato OS-AAAA-NNNN & Folio de la orden \\
 & estado & estado\_orden & No & 7 valores (6.5) & Estado actual \\
 & numero\_bahia & SMALLINT & Sí & FK bahia; único entre órdenes activas & Bahía asignada (RN-05) \\
 & id\_usuario\_responsable & INTEGER & Sí & FK usuario & Mecánico responsable (RN-09) \\
 & entrega\_a\_credito & BOOLEAN & No & Requiere gerente y flotilla & Excepción de RN-10 \\
 & fecha\_entrega & TIMESTAMPTZ & Sí & $\geq$ fecha\_apertura & Momento de la entrega \\
recepcion\_vehicular & kilometraje & INTEGER & No & $\geq 0$ & Lectura del odómetro \\
 & nivel\_combustible & nivel\_combustible & No & 5 niveles & Combustible al ingreso \\
 & id\_cita & INTEGER & Sí & FK cita, único; solo cita Confirmada & Cita de origen \\
fotografia & posicion & posicion\_foto & No & Una por posición salvo «Otra» & Frente, Trasera, Izquierda, Derecha (RN-12) \\
 & archivo & VARCHAR(255) & No & & Ruta en el almacenamiento \\
cotizacion & version & INTEGER & No & $\geq 1$; único con folio & Versión de la cotización \\
 & fecha\_emision & DATE & No & & Inicio de vigencia \\
 & vigencia\_dias & INTEGER & No & $> 0$, por omisión 5 & RN-11 \\
 & estado & estado\_cotizacion & No & Pendiente, Autorizada, Rechazada, Vencida & Estado (6.5) \\
detalle\_cotizacion & tipo & tipo\_concepto & No & ManoDeObra, Refaccion, Insumo & Clase de concepto \\
 & precio\_unitario & NUMERIC(10,2) & No & $\geq 0$ & Precio en pesos \\
 & opcional / seleccionado & BOOLEAN & No & Solo un opcional puede deseleccionarse & RN-03 \\
autorizacion & medio & medio\_autorizacion & No & Remoto exige IP y token & Canal de la autorización \\
 & tipo & tipo\_autorizacion & No & Inicial o IncrementoMayor10 & RN-04 \\
 & monto\_aceptado & NUMERIC(10,2) & No & $\geq 0$ & Monto que el cliente aceptó \\
 & codigo\_verificacion\_hash & VARCHAR(255) & No & PBKDF2; nunca en claro & Código 2FA (RN-02) \\
 & huella\_sha256 & CHAR(64) & Sí & Obligatoria si Confirmada & Detecta alteraciones (RNF-04) \\
token\_acceso & token\_hash & CHAR(64) & No & PK; SHA-256 del token & El token en claro solo viaja en el enlace \\
 & fecha\_expiracion & TIMESTAMPTZ & Sí & & Se fija al entregar y saldar (RN-01) \\
pago & monto & NUMERIC(10,2) & No & $> 0$; solo sobre cotización Autorizada & Abono del cliente \\
\end{longtable}
}

\subsection{6.3 Diagrama de clases del
dominio}\label{diagrama-de-clases-del-dominio}

El diagrama de clases v4 tiene 26 clases del dominio (mínimo del Equipo B: 12), con atributos,
asociaciones, multiplicidades, composiciones, una generalización \{disjoint, complete\} (Cliente →
ClienteParticular | Flotilla), 15 enumeraciones y las restricciones de negocio como notas OCL
informales. Las clases marcadas «OfflineSync» (RecepcionVehicular, ChecklistInspeccion, Diagnostico) se
capturan sin conexión (RNF-08). Se divide en dos láminas por legibilidad; OrdenServicio, Usuario y
TokenAcceso aparecen en ambas como vista parcial.

\figura{../06-analisis/clases/diagrama-clases-v4-1.png}{angle=90,width=\linewidth,height=0.8\textheight}{6.5a}{Diagrama
de clases del dominio (1/2): clientes, citas, recepción y orden}{docs/06-analisis/clases/diagrama-clases-v4-1.png}

\figura{../06-analisis/clases/diagrama-clases-v4-2.png}{angle=90,width=\linewidth,height=0.8\textheight}{6.5b}{Diagrama
de clases del dominio (2/2): cotización, compras, pagos y garantía}{docs/06-analisis/clases/diagrama-clases-v4-2.png}

\pendiente{subir al repositorio el archivo fuente editable del diagrama de clases (draw.io, PlantUML o
StarUML) junto a las imágenes, como pide el machote.}

\subsection{6.4 Diagrama de objetos}\label{diagrama-de-objetos}

\figurapendiente{6.6}{Diagrama de objetos: orden autorizada por el cliente}{instantánea con datos
ficticios del caso CPA-13: \texttt{c1:ClienteParticular} (nombre = «Cliente Ficticio»), \texttt{v1:Vehiculo}
(placas = «PRB-001», Nissan Versa 2019), \texttt{os1:OrdenServicio} (folio = «OS-2026-0001», estado =
Autorizada, bahía 1), \texttt{r1:RecepcionVehicular} (kilometraje = 45210, combustible = Mitad) con su
\texttt{ch1:ChecklistInspeccion} (completo = true) y 4 \texttt{Fotografia}; \texttt{d1:Diagnostico};
\texttt{cot1:Cotizacion} (versión 1, estado = Autorizada, total = 1750.00) con 3 \texttt{DetalleCotizacion};
\texttt{a1:Autorizacion} (medio = Remoto, estado = Confirmada, monto = 1750.00) ligada a
\texttt{t1:TokenAcceso}; \texttt{m1:Usuario} (rol = Mecanico) como responsable.}{docs/06-analisis/objetos/orden-autorizada.puml}

\subsection{6.5 Diagramas de estados}\label{diagramas-de-estados}

Las transiciones de la orden están programadas en una sola tabla (\texttt{TRANSICIONES} en
\texttt{backend/app/servicios/ordenes.py}); cualquier transición que no aparezca aquí es rechazada por
el sistema y toda transición válida queda en HistorialEstado.

\figurapendiente{6.7}{Máquina de estados de la Orden de servicio}{los 7 estados (Abierta,
EnEsperaDeBahia, Diagnosticada, Autorizada, ListaParaEntrega, Cerrada, Cancelada) con las transiciones de
la tabla siguiente, cada una con evento, [condición] y / rol. Renderizar el archivo fuente con
PlantUML.}{docs/06-analisis/estados/orden-servicio.puml}

{\footnotesize
\begin{longtable}{@{}L{0.15\linewidth}L{0.19\linewidth}L{0.31\linewidth}L{0.14\linewidth}L{0.15\linewidth}@{}}
\toprule
\textbf{Estado origen} & \textbf{Evento} & \textbf{Condición} & \textbf{Estado destino} & \textbf{Rol autorizado} \\
\midrule
\endhead
\bottomrule
\endlastfoot
(inicio) & Registrar recepción & El vehículo no tiene otra orden abierta; si no hay bahía, hay lugar en el estacionamiento (RN-05) & Abierta & Recepcionista \\
Abierta & Recepción sin bahía libre & Las 4 bahías están ocupadas & EnEsperaDeBahia & Recepcionista \\
EnEsperaDeBahia & Asignar bahía & Bahía libre; responsable es mecánico activo (RN-09) & Abierta & Jefe de Mecánicos \\
Abierta & Registrar diagnóstico & Checklist completo y 4 fotografías (RN-12); bahía y responsable asignados & Diagnosticada & Jefe de Mecánicos \\
Diagnosticada & Confirmar código 2FA & Cotización vigente (RN-11); código correcto, de un solo uso y con menos de 5 minutos (RN-02) & Autorizada & Cliente (remoto) o Recepcionista (presencial) \\
Diagnosticada & Rechazar cotización & Cotización vigente & Cancelada & Cliente \\
Autorizada & Terminar trabajo & No hay re-cotización pendiente mayor al 10\,\% (RN-04) & ListaParaEntrega & Jefe de Mecánicos \\
ListaParaEntrega & Registrar entrega & Saldo pendiente igual a 0, o entrega a crédito autorizada (RN-10) & Cerrada & Recepcionista \\
Abierta, EnEsperaDeBahia, Diagnosticada & Cancelar & --- & Cancelada & Gerente \\
\end{longtable}
}

\figurapendiente{6.8}{Máquina de estados de la Cotización}{los estados Pendiente, Autorizada,
Rechazada y Vencida, con las transiciones de la tabla siguiente.}{docs/06-analisis/estados/cotizacion.puml}

{\footnotesize
\begin{longtable}{@{}L{0.15\linewidth}L{0.19\linewidth}L{0.31\linewidth}L{0.14\linewidth}L{0.15\linewidth}@{}}
\toprule
\textbf{Estado origen} & \textbf{Evento} & \textbf{Condición} & \textbf{Estado destino} & \textbf{Rol autorizado} \\
\midrule
\endhead
\bottomrule
\endlastfoot
(inicio) & Generar cotización & Orden Diagnosticada (inicial) o Autorizada (re-cotización); al menos un concepto; sin otra cotización pendiente vigente & Pendiente & Recepcionista o Jefe de Mecánicos \\
(inicio) & Re-cotizar & Orden Autorizada y nuevo total $\leq$ total inicial autorizado $\times$ 1.10 (RN-04) & Autorizada & Recepcionista o Jefe de Mecánicos \\
Pendiente & Cambiar selección & Solo conceptos opcionales (RN-03); invalida códigos ya emitidos & Pendiente & Cliente o Recepcionista \\
Pendiente & Confirmar código 2FA & Vigente y código válido (RN-02) & Autorizada & Cliente o Recepcionista \\
Pendiente & Rechazar & Vigente & Rechazada & Cliente \\
Pendiente & Pasa la fecha de vencimiento & Hoy $>$ fecha de emisión + 5 días (RN-11) & Vencida & Sistema \\
\end{longtable}
}

\subsection{6.6 Diagramas de actividades
(to-be)}\label{diagramas-de-actividades-to-be}

\figurapendiente{6.9}{Actividad to-be: recepción y diagnóstico}{carriles Cliente, Recepcionista,
Sistema y Jefe de mecánicos: búsqueda o registro del cliente con aviso de privacidad, captura de
kilometraje, combustible, checklist y 4 fotografías, apertura de la orden, decisión «¿hay bahía libre?»
(EnEsperaDeBahia con alerta), asignación de bahía y responsable, registro del diagnóstico y validación de
RN-12.}{docs/06-analisis/actividades/recepcion-diagnostico.puml}

\figurapendiente{6.10}{Actividad to-be: cotización y autorización con código 2FA}{carriles Jefe de
mecánicos, Sistema y Cliente: generación de la cotización con vigencia de 5 días, envío del enlace seguro,
selección de opcionales, decisión aceptar/rechazar, emisión y validación del código de 6 dígitos, registro
de la autorización con huella SHA-256, ejecución del trabajo y re-cotización bloqueante si el adicional
supera el 10\,\%.}{docs/06-analisis/actividades/cotizacion-autorizacion.puml}

\section{7. Arquitectura y diseño}\label{arquitectura-y-diseuxf1o}

\subsection{7.1 Estilo arquitectónico}\label{estilo-arquitectuxf3nico}

\begin{quote}
\emph{\textbf{Instrucción:} Expliquen el estilo elegido (capas, MVC,
cliente-servidor con API\ldots) y cómo atiende al menos \textbf{4
requisitos no funcionales} de la sección 4.6.}
\end{quote}

{\def\LTcaptype{none} % do not increment counter
\begin{longtable}[]{@{}
  >{\raggedright\arraybackslash}p{(\linewidth - 2\tabcolsep) * \real{0.1922}}
  >{\raggedright\arraybackslash}p{(\linewidth - 2\tabcolsep) * \real{0.8078}}@{}}
\toprule\noalign{}
\begin{minipage}[b]{\linewidth}\raggedright
\textbf{RNF}
\end{minipage} & \begin{minipage}[b]{\linewidth}\raggedright
\textbf{Decisión de arquitectura que lo atiende}
\end{minipage} \\
\midrule\noalign{}
\endhead
\bottomrule\noalign{}
\endlastfoot
& \\
& \\
& \\
& \\
\end{longtable}
}

\subsection{7.2 Diagramas C4}\label{diagramas-c4}

\begin{quote}
\emph{\textbf{Instrucción:} Contexto, Contenedores y Componentes de al
menos un contenedor.}
\end{quote}

{[}Figuras 7.1 a 7.3.{]}

\subsection{7.3 Registros de decisiones de arquitectura
(ADR)}\label{registros-de-decisiones-de-arquitectura-adr}

\begin{quote}
\emph{\textbf{Instrucción:} Mínimo \textbf{4 ADR}: estilo
arquitectónico, pila tecnológica, base de datos y hospedaje. Equipo A:
un ADR adicional sobre el soporte sin conexión. Equipo B: un ADR
adicional sobre el almacenamiento de fotos y el acceso del cliente
mediante un enlace seguro. Cada ADR evalúa al menos dos alternativas.}
\end{quote}

\subsubsection{Plantilla de ADR (repetir por cada
decisión)}\label{plantilla-de-adr-repetir-por-cada-decisiuxf3n}

{\def\LTcaptype{none} % do not increment counter
\begin{longtable}[]{@{}
  >{\raggedright\arraybackslash}p{(\linewidth - 2\tabcolsep) * \real{0.2778}}
  >{\raggedright\arraybackslash}p{(\linewidth - 2\tabcolsep) * \real{0.7222}}@{}}
\toprule\noalign{}
\begin{minipage}[b]{\linewidth}\raggedright
\textbf{Campo}
\end{minipage} & \begin{minipage}[b]{\linewidth}\raggedright
\textbf{Contenido}
\end{minipage} \\
\midrule\noalign{}
\endhead
\bottomrule\noalign{}
\endlastfoot
ID y título & ADR-001 {[}Decisión{]} \\
Fecha y estado & {[}Propuesto / aceptado / reemplazado por
ADR-\ldots{]} \\
Contexto & {[}Qué problema o fuerza obliga a decidir{]} \\
Alternativas evaluadas & 1. \ldots{} 2. \ldots{} \\
Decisión & \\
Consecuencias & {[}Positivas y negativas; qué se vuelve más fácil o más
difícil{]} \\
RNF relacionados & RNF-\ldots{} \\
\end{longtable}
}

\subsection{7.4 Diagramas de artefactos y de
despliegue}\label{diagramas-de-artefactos-y-de-despliegue}

\begin{quote}
\emph{\textbf{Instrucción:} Deben reflejar el ambiente de pruebas real:
nodos, artefactos instalados y protocolos de comunicación.}
\end{quote}

{[}Figuras 7.4 y 7.5.{]}

\subsection{7.5 Diagrama de paquetes}\label{diagrama-de-paquetes}

\begin{quote}
\emph{\textbf{Instrucción:} Organización de módulos y sus dependencias,
sin dependencias circulares.}
\end{quote}

{[}Figura 7.6.{]}

\subsection{7.6 Diagramas de clases de la
solución}\label{diagramas-de-clases-de-la-soluciuxf3n}

\begin{quote}
\emph{\textbf{Instrucción:} Clases del flujo del núcleo tal como existen
en el código (controladores, servicios, repositorios, entidades). Deben
coincidir con el repositorio.}
\end{quote}

{[}Figura 7.7.{]}

\subsection{7.7 Diagramas de secuencia}\label{diagramas-de-secuencia}

\begin{quote}
\emph{\textbf{Instrucción:} Mínimo \textbf{4}, incluido el escenario
principal del flujo del núcleo.}
\end{quote}

{[}Figuras 7.8 a 7.11.{]}

\subsection{7.8 Diagrama de
comunicación}\label{diagrama-de-comunicaciuxf3n}

\begin{quote}
\emph{\textbf{Instrucción:} Al menos uno.}
\end{quote}

{[}Figura 7.12.{]}

\section{8. Diseño de la interfaz de
usuario}\label{diseuxf1o-de-la-interfaz-de-usuario}

\textcolor{red}{\textbf{Pendiente (corte Review 5, Sprints 4 y 5).} La API del flujo núcleo ya existe
(sección 9); el prototipo de alta fidelidad, la evaluación heurística y la prueba con usuarios se hacen
sobre los wireframes de 4.10 una vez validados con el cliente.}

\subsection{8.1 Prototipo de alta
fidelidad}\label{prototipo-de-alta-fidelidad}

\begin{quote}
\emph{\textbf{Instrucción:} Enlace al prototipo navegable y capturas del
flujo principal y de dos flujos adicionales. Expliquen las decisiones de
diseño según el contexto de uso de cada tipo de usuario (dispositivo,
conectividad, experiencia).}
\end{quote}

{[}Enlace y figuras.{]}

\subsection{8.2 Accesibilidad}\label{accesibilidad}

\begin{quote}
\emph{\textbf{Instrucción:} Qué criterios de WCAG 2.2 se consideraron
(contraste, tamaño de objetivos táctiles, textos alternativos,
navegación por teclado) y cómo se verificaron.}
\end{quote}

{[}Escriba aquí.{]}

\subsection{8.3 Evaluación heurística
recibida}\label{evaluaciuxf3n-heuruxedstica-recibida}

\begin{quote}
\emph{\textbf{Instrucción:} Hallazgos que reportó el otro equipo con las
10 heurísticas de Nielsen, y cómo se atendieron.}
\end{quote}

{\def\LTcaptype{none} % do not increment counter
\begin{longtable}[]{@{}
  >{\raggedright\arraybackslash}p{(\linewidth - 8\tabcolsep) * \real{0.1174}}
  >{\raggedright\arraybackslash}p{(\linewidth - 8\tabcolsep) * \real{0.2138}}
  >{\raggedright\arraybackslash}p{(\linewidth - 8\tabcolsep) * \real{0.3056}}
  >{\raggedright\arraybackslash}p{(\linewidth - 8\tabcolsep) * \real{0.1282}}
  >{\raggedright\arraybackslash}p{(\linewidth - 8\tabcolsep) * \real{0.2350}}@{}}
\toprule\noalign{}
\begin{minipage}[b]{\linewidth}\raggedright
\textbf{Issue}
\end{minipage} & \begin{minipage}[b]{\linewidth}\raggedright
\textbf{Heurística}
\end{minipage} & \begin{minipage}[b]{\linewidth}\raggedright
\textbf{Hallazgo}
\end{minipage} & \begin{minipage}[b]{\linewidth}\raggedright
\textbf{Severidad}
\end{minipage} & \begin{minipage}[b]{\linewidth}\raggedright
\textbf{Acción tomada}
\end{minipage} \\
\midrule\noalign{}
\endhead
\bottomrule\noalign{}
\endlastfoot
& & & & \\
& & & & \\
& & & & \\
\end{longtable}
}

\subsection{8.4 Prueba con usuarios y
SUS}\label{prueba-con-usuarios-y-sus}

\begin{quote}
\emph{\textbf{Instrucción:}} \emph{Prueba con} \emph{\textbf{3 a 5
personas}} \emph{ajenas al equipo:} \emph{tareas que realizaron,
problemas observados, puntaje} \emph{promedio del SUS con su cálculo y
cambios derivados.}
\end{quote}

{\def\LTcaptype{none} % do not increment counter
\begin{longtable}[]{@{}
  >{\raggedright\arraybackslash}p{(\linewidth - 10\tabcolsep) * \real{0.2138}}
  >{\raggedright\arraybackslash}p{(\linewidth - 10\tabcolsep) * \real{0.1816}}
  >{\raggedright\arraybackslash}p{(\linewidth - 10\tabcolsep) * \real{0.1174}}
  >{\raggedright\arraybackslash}p{(\linewidth - 10\tabcolsep) * \real{0.1069}}
  >{\raggedright\arraybackslash}p{(\linewidth - 10\tabcolsep) * \real{0.2521}}
  >{\raggedright\arraybackslash}p{(\linewidth - 10\tabcolsep) * \real{0.1281}}@{}}
\toprule\noalign{}
\begin{minipage}[b]{\linewidth}\raggedright
\textbf{Participante (perfil, sin nombre)}
\end{minipage} & \begin{minipage}[b]{\linewidth}\raggedright
\textbf{Tarea}
\end{minipage} & \begin{minipage}[b]{\linewidth}\raggedright
\textbf{¿La completó?}
\end{minipage} & \begin{minipage}[b]{\linewidth}\raggedright
\textbf{Tiempo}
\end{minipage} & \begin{minipage}[b]{\linewidth}\raggedright
\textbf{Problemas observados}
\end{minipage} & \begin{minipage}[b]{\linewidth}\raggedright
\textbf{Puntaje SUS}
\end{minipage} \\
\midrule\noalign{}
\endhead
\bottomrule\noalign{}
\endlastfoot
& & & & & \\
& & & & & \\
& & & & & \\
\end{longtable}
}

\textbf{Puntaje SUS promedio:} {[} {]} · \textbf{Cambios realizados a
partir de la prueba:} {[} {]}

\section{9. Implementación del Release
v0.1.0}\label{implementaciuxf3n-del-release-v0.1.0}

\subsection{9.1 Pila tecnológica}\label{pila-tecnoluxf3gica}

{\def\LTcaptype{none} % do not increment counter
\begin{longtable}[]{@{}
  >{\raggedright\arraybackslash}p{(\linewidth - 4\tabcolsep) * \real{0.2990}}
  >{\raggedright\arraybackslash}p{(\linewidth - 4\tabcolsep) * \real{0.4060}}
  >{\raggedright\arraybackslash}p{(\linewidth - 4\tabcolsep) * \real{0.2950}}@{}}
\toprule\noalign{}
\begin{minipage}[b]{\linewidth}\raggedright
\textbf{Capa}
\end{minipage} & \begin{minipage}[b]{\linewidth}\raggedright
\textbf{Tecnología y versión}
\end{minipage} & \begin{minipage}[b]{\linewidth}\raggedright
\textbf{ADR que la justifica}
\end{minipage} \\
\midrule\noalign{}
\endhead
\bottomrule\noalign{}
\endlastfoot
Interfaz & & \\
Lógica de negocio / API & & \\
Base de datos & & \\
Hospedaje del ambiente de pruebas & & \\
Integración continua & & \\
\end{longtable}
}

\subsection{9.2 Instrucciones de instalación y
ejecución}\label{instrucciones-de-instalaciuxf3n-y-ejecuciuxf3n}

\begin{quote}
\emph{\textbf{Instrucción:} Pasos para que una persona ajena al equipo
instale y ejecute el sistema localmente siguiendo el README. Incluyan
cómo cargar los datos de prueba ficticios.}
\end{quote}

{[}Escriba aquí o indique la sección del README.{]}

\subsection{9.3 Integración continua y flujo de
trabajo}\label{integraciuxf3n-continua-y-flujo-de-trabajo}

\begin{quote}
\emph{\textbf{Instrucción:} Qué ejecuta el pipeline (compilación,
linter, pruebas), reglas de protección de main y estadísticas del
repositorio: número de pull requests fusionados y revisados.}
\end{quote}

{[}Escriba aquí y agregue una captura de pantalla de una ejecución
exitosa.{]}

\subsection{9.4 Evidencia del flujo
núcleo}\label{evidencia-del-flujo-nuxfacleo}

\begin{quote}
\emph{\textbf{Instrucción:} Capturas paso a paso del flujo del núcleo
funcionando en el \textbf{ambiente de pruebas} (no en local), con datos
ficticios. Equipo A: registro de pedido, cálculo del total por tipo de
cliente y consulta de la producción requerida. Equipo B: apertura de
orden con checklist, diagnóstico, cotización con autorización registrada
y lista de órdenes abiertas.}
\end{quote}

{[}Figuras 9.x.{]}

\section{10. Reporte final de pruebas e
incidencias}\label{reporte-final-de-pruebas-e-incidencias}

\subsection{10.1 Resultado de los casos de
prueba}\label{resultado-de-los-casos-de-prueba}

\textcolor{red}{Corte parcial (29/10/2026). El reporte final se completa en el Sprint 6 con la
ejecución de aceptación con el cliente y el QA cruzado.}

{\def\LTcaptype{none} % do not increment counter
\begin{longtable}[]{@{}
  >{\raggedright\arraybackslash}p{(\linewidth - 2\tabcolsep) * \real{0.5984}}
  >{\raggedright\arraybackslash}p{(\linewidth - 2\tabcolsep) * \real{0.4016}}@{}}
\toprule\noalign{}
\begin{minipage}[b]{\linewidth}\raggedright
\textbf{Indicador}
\end{minipage} & \begin{minipage}[b]{\linewidth}\raggedright
\textbf{Valor}
\end{minipage} \\
\midrule\noalign{}
\endhead
\bottomrule\noalign{}
\endlastfoot
Casos de prueba diseñados & 35 (CPA-01 a CPA-35, sección 5.2) \\
Casos ejecutados & 35, automatizados \\
Pasan & 35 \\
Fallan & 0. CPA-11 falló en su primera ejecución (el portal respondía 404 ante una cotización
vencida, en lugar del aviso de RN-11); se corrigió antes de integrar a main. \\
Bloqueados & 0 \\
Pruebas unitarias automatizadas (número y resultado en CI) & 35 pruebas de pytest; se ejecutan en
GitHub Actions en cada push y Pull Request \\
\end{longtable}
}

\subsection{10.2 Incidencias}\label{incidencias}

\begin{quote}
\emph{\textbf{Instrucción:} Incluyan las de todos los orígenes:
internas, QA cruzado y cliente. El tiempo de atención se mide desde que
se registró hasta que se cerró.}
\end{quote}

{\def\LTcaptype{none} % do not increment counter
\begin{longtable}[]{@{}
  >{\raggedright\arraybackslash}p{(\linewidth - 10\tabcolsep) * \real{0.1497}}
  >{\raggedright\arraybackslash}p{(\linewidth - 10\tabcolsep) * \real{0.1495}}
  >{\raggedright\arraybackslash}p{(\linewidth - 10\tabcolsep) * \real{0.1496}}
  >{\raggedright\arraybackslash}p{(\linewidth - 10\tabcolsep) * \real{0.2094}}
  >{\raggedright\arraybackslash}p{(\linewidth - 10\tabcolsep) * \real{0.1497}}
  >{\raggedright\arraybackslash}p{(\linewidth - 10\tabcolsep) * \real{0.1922}}@{}}
\toprule\noalign{}
\begin{minipage}[b]{\linewidth}\raggedright
\textbf{Severidad}
\end{minipage} & \begin{minipage}[b]{\linewidth}\raggedright
\textbf{Registradas}
\end{minipage} & \begin{minipage}[b]{\linewidth}\raggedright
\textbf{Cerradas}
\end{minipage} & \begin{minipage}[b]{\linewidth}\raggedright
\textbf{Rechazadas (con justificación)}
\end{minipage} & \begin{minipage}[b]{\linewidth}\raggedright
\textbf{Abiertas}
\end{minipage} & \begin{minipage}[b]{\linewidth}\raggedright
\textbf{Tiempo promedio de atención}
\end{minipage} \\
\midrule\noalign{}
\endhead
\bottomrule\noalign{}
\endlastfoot
S1 & & & & & \\
S2 & & & & & \\
S3 & & & & & \\
S4 & & & & & \\
\end{longtable}
}

\subsection{10.3 Defectos conocidos que pasan a la Fase
2}\label{defectos-conocidos-que-pasan-a-la-fase-2}

{\def\LTcaptype{none} % do not increment counter
\begin{longtable}[]{@{}
  >{\raggedright\arraybackslash}p{(\linewidth - 8\tabcolsep) * \real{0.1174}}
  >{\raggedright\arraybackslash}p{(\linewidth - 8\tabcolsep) * \real{0.3483}}
  >{\raggedright\arraybackslash}p{(\linewidth - 8\tabcolsep) * \real{0.1282}}
  >{\raggedright\arraybackslash}p{(\linewidth - 8\tabcolsep) * \real{0.1282}}
  >{\raggedright\arraybackslash}p{(\linewidth - 8\tabcolsep) * \real{0.2779}}@{}}
\toprule\noalign{}
\begin{minipage}[b]{\linewidth}\raggedright
\textbf{Issue}
\end{minipage} & \begin{minipage}[b]{\linewidth}\raggedright
\textbf{Descripción}
\end{minipage} & \begin{minipage}[b]{\linewidth}\raggedright
\textbf{Severidad}
\end{minipage} & \begin{minipage}[b]{\linewidth}\raggedright
\textbf{Prioridad}
\end{minipage} & \begin{minipage}[b]{\linewidth}\raggedright
\textbf{Motivo por el que no se corrigió}
\end{minipage} \\
\midrule\noalign{}
\endhead
\bottomrule\noalign{}
\endlastfoot
& & & & \\
& & & & \\
\end{longtable}
}

\subsection{10.4 QA cruzado realizado al otro
equipo}\label{qa-cruzado-realizado-al-otro-equipo}

\begin{quote}
\emph{\textbf{Instrucción:}} \emph{Resumen de los casos ejecutados en}
\emph{el sistema del otro equipo, los defectos reportados y las}
\emph{correcciones verificadas en la ronda 2.}
\end{quote}

{[}Escriba aquí.{]}

\section{11. Paquete de transferencia a la Fase
2}\label{paquete-de-transferencia-a-la-fase-2}

\begin{quote}
\emph{\textbf{Instrucción:} Lo necesario para que un equipo, incluso uno
distinto, continúe el proyecto de Ingeniería de Software II.}
\end{quote}

{\def\LTcaptype{none} % do not increment counter
\begin{longtable}[]{@{}
  >{\raggedright\arraybackslash}p{(\linewidth - 2\tabcolsep) * \real{0.3204}}
  >{\raggedright\arraybackslash}p{(\linewidth - 2\tabcolsep) * \real{0.6796}}@{}}
\toprule\noalign{}
\begin{minipage}[b]{\linewidth}\raggedright
\textbf{Apartado}
\end{minipage} & \begin{minipage}[b]{\linewidth}\raggedright
\textbf{Contenido}
\end{minipage} \\
\midrule\noalign{}
\endhead
\bottomrule\noalign{}
\endlastfoot
Backlog restante priorizado & {[}Enlace y resumen por módulo{]} \\
Deuda técnica conocida & \\
Riesgos abiertos & \\
Decisiones pendientes con el cliente & \\
Accesos y configuración & {[}Qué cuentas y servicios existen; nunca
contraseñas en el documento{]} \\
Recomendaciones & \\
\end{longtable}
}

\section{12. Retrospectiva del proyecto y contribución
individual}\label{retrospectiva-del-proyecto-y-contribuciuxf3n-individual}

\subsection{12.1 Retrospectiva del
proyecto}\label{retrospectiva-del-proyecto}

{\def\LTcaptype{none} % do not increment counter
\begin{longtable}[]{@{}
  >{\raggedright\arraybackslash}p{(\linewidth - 4\tabcolsep) * \real{0.3333}}
  >{\raggedright\arraybackslash}p{(\linewidth - 4\tabcolsep) * \real{0.3333}}
  >{\raggedright\arraybackslash}p{(\linewidth - 4\tabcolsep) * \real{0.3333}}@{}}
\toprule\noalign{}
\begin{minipage}[b]{\linewidth}\raggedright
\textbf{¿Qué funcionó bien?}
\end{minipage} & \begin{minipage}[b]{\linewidth}\raggedright
\textbf{¿Qué no funcionó?}
\end{minipage} & \begin{minipage}[b]{\linewidth}\raggedright
\textbf{¿Qué haríamos distinto en la Fase 2?}
\end{minipage} \\
\midrule\noalign{}
\endhead
\bottomrule\noalign{}
\endlastfoot
\textasciitilde{} {[} {]} & \textasciitilde{} {[} {]} &
\textasciitilde{} {[} {]} \\
\end{longtable}
}

\subsection{12.2 Contribución
individual}\label{contribuciuxf3n-individual}

\begin{quote}
\emph{\textbf{Instrucción:} Cada integrante describe, con sus propias
palabras, su aporte y enlaza la evidencia. El profesor lo contrasta con
el repositorio, el tablero y la coevaluación; esta última se entrega por
separado en un formulario privado.}
\end{quote}

{\def\LTcaptype{none} % do not increment counter
\begin{longtable}[]{@{}
  >{\raggedright\arraybackslash}p{(\linewidth - 6\tabcolsep) * \real{0.1923}}
  >{\raggedright\arraybackslash}p{(\linewidth - 6\tabcolsep) * \real{0.3205}}
  >{\raggedright\arraybackslash}p{(\linewidth - 6\tabcolsep) * \real{0.2735}}
  >{\raggedright\arraybackslash}p{(\linewidth - 6\tabcolsep) * \real{0.2137}}@{}}
\toprule\noalign{}
\begin{minipage}[b]{\linewidth}\raggedright
\textbf{Integrante}
\end{minipage} & \begin{minipage}[b]{\linewidth}\raggedright
\textbf{Principales aportes}
\end{minipage} & \begin{minipage}[b]{\linewidth}\raggedright
\textbf{Evidencia (issues, PR, documentos)}
\end{minipage} & \begin{minipage}[b]{\linewidth}\raggedright
\textbf{Revisiones de código o documentos hechas a otros}
\end{minipage} \\
\midrule\noalign{}
\endhead
\bottomrule\noalign{}
\endlastfoot
& & & \\
& & & \\
& & & \\
& & & \\
& & & \\
& & & \\
\end{longtable}
}

\section{13. Referencias}\label{referencias}

\begin{quote}
\emph{\textbf{Instrucción:} Formato APA 7, en orden alfabético. Solo
fuentes que realmente consultaron y que se pueden verificar. Como mínimo
deben aparecer las normas y guías que usaron:}
\end{quote}

Schwaber, K., \& Sutherland, J. (2020). \emph{La Guía de Scrum}.
\href{https://scrumguides.org/}{\ul{https://scrumguides.org/}}

ISO/IEC/IEEE. (2018). \emph{ISO/IEC/IEEE 29148:2018 Systems and software
engineering --- Life cycle processes --- Requirements engineering}. ISO.

ISO/IEC. (2023). \emph{ISO/IEC 25010:2023 Systems and software
engineering --- SQuaRE --- Product quality model}. ISO.

ISO/IEC/IEEE. (2021). \emph{ISO/IEC/IEEE 29119-3:2021 Software and
systems engineering --- Software testing --- Part 3: Test
documentation}. ISO.

Object Management Group. (2017). \emph{OMG Unified Modeling Language,
Version 2.5.1}.
\href{https://www.omg.org/spec/UML/2.5.1}{\ul{https://www.omg.org/spec/UML/2.5.1}}

{[}Agreguen aquí el resto de sus referencias.{]}

\section{Anexos}\label{anexos}

\begin{quote}
\emph{\textbf{Instrucción:} Numeren los anexos con letras. Contenido
mínimo:}
\end{quote}

\begin{itemize}
\item
  \textbf{Anexo A.} Guiones y minutas de las entrevistas con el cliente
  (mínimo 20 preguntas en el guion inicial).
\item
  \textbf{Anexo B.} Minutas de las Sprint Reviews con aceptaciones y
  rechazos del cliente.
\item
  \textbf{Anexo C.} Detalle completo de las historias de usuario y los
  casos de prueba, si no caben en el cuerpo.
\item
  \textbf{Anexo D.} Manual breve de uso del flujo del núcleo (Equipo A:
  para la auxiliar de oficina; Equipo B: para la recepcionista).
\item
  \textbf{Anexo E.} Aviso de privacidad (borrador).
\end{itemize}

\section{Lista de verificación antes de
entregar}\label{lista-de-verificaciuxf3n-antes-de-entregar}

\begin{quote}
\emph{\textbf{Instrucción:} Revísenla en equipo antes de cada entrega.
Bórrenla del documento final o déjenla marcada como evidencia de
revisión.}
\end{quote}

{\def\LTcaptype{none} % do not increment counter
\begin{longtable}[]{@{}
  >{\raggedright\arraybackslash}p{(\linewidth - 2\tabcolsep) * \real{0.0640}}
  >{\raggedright\arraybackslash}p{(\linewidth - 2\tabcolsep) * \real{0.9360}}@{}}
\toprule\noalign{}
\begin{minipage}[b]{\linewidth}\raggedright
\textbf{✓}
\end{minipage} & \begin{minipage}[b]{\linewidth}\raggedright
\textbf{Verificación}
\end{minipage} \\
\midrule\noalign{}
\endhead
\bottomrule\noalign{}
\endlastfoot
☐ & Se borraron todas las instrucciones y los marcadores entre
corchetes. \\
☐ & La portada, el control de versiones y los enlaces están
actualizados. \\
☐ & Las secciones del corte están completas según la tabla 0.2. \\
☐ & Se alcanzan los mínimos de su equipo: historias, RNF, glosario,
casos de uso, casos de prueba, clases del dominio, ADR, secuencias y
riesgos. \\
☐ & Todo requisito no funcional tiene una métrica y un umbral. \\
☐ & La matriz de trazabilidad no tiene requisitos Must con historia ni
caso de prueba. \\
☐ & Los nombres del glosario, del SRS, de los modelos y del código
coinciden. \\
☐ & Cada figura tiene número, título, es legible y su archivo fuente
está en el repositorio. \\
☐ & Las solicitudes de cambio y las incidencias están registradas junto
con su decisión. \\
☐ & Se declaró el uso de IA. \\
☐ & No hay datos personales reales ni contraseñas en el documento ni en
el repositorio. \\
☐ & Las referencias están en formato APA 7 y son verificables. \\
☐ & El PDF y el .docx están en docs/entregas/ con el nombre correcto. \\
☐ & Otro integrante, distinto del autor principal, revisó el documento
completo. \\
\end{longtable}
}


\end{document}
